% GENERATED FILE -- DO NOT EDIT.
%
% Produced by harness/assemble_contract.py from module.yaml and claims/*.tex.
% Edit the claim files and regenerate; `--check` fails if this file is stale.
\documentclass[11pt,a4paper]{article}
\pdftrailerid{79f65317c18e8b5896fb4bab363d91c2}
%% --- 1. Mathematics (order matters: physics after braket wins \ket, \tr) ---
\usepackage{amsmath}
\usepackage{amssymb}
\usepackage{amsfonts}
\usepackage{amsthm}
\usepackage{mathtools}
\usepackage{dsfont}
\usepackage{braket}
\usepackage{physics}
%% --- 2. Page furniture ------------------------------------------------------
\usepackage{xcolor}
\usepackage{array}    % >{...} column prefixes in the wide tables of some claims
\usepackage{enumitem}
\usepackage{graphicx}
\usepackage{booktabs}
\usepackage{longtable}
\usepackage{xltabular}              % front- and back-matter tables across pages
\usepackage{needspace}
\usepackage{xurl}
\usepackage{tcolorbox}
\tcbuselibrary{skins,breakable}
%% --- 3. Links and navigation (hyperref late; bookmark, refcount after it) ---
\usepackage{xr-hyper}   % before hyperref: the contract's labels
\usepackage{hyperref}
\usepackage{bookmark}
\usepackage{refcount}
%% --- 4. Colours: every distinction also has a shape (rule, caps, size, weight)
\definecolor{nnGrey}{gray}{0.40}      % non-normative text
\definecolor{metaGrey}{gray}{0.30}    % the generated meta line
\definecolor{techGrey}{gray}{0.45}    % heading of a technical claim
\definecolor{mainRule}{RGB}{0,110,140}
\definecolor{reqAccent}{RGB}{165,70,0}
\definecolor{linkBlue}{RGB}{0,62,150}
\definecolor{extLink}{RGB}{120,40,140}  % links into another PDF
\hypersetup{colorlinks=true, linkcolor=linkBlue, citecolor=linkBlue,
  urlcolor=linkBlue, filecolor=extLink, bookmarksnumbered=true,
  bookmarksopen=true, bookmarksopenlevel=1, pdfdisplaydoctitle=true}
% A link into another PDF (\href to a .pdf, \hyperref{file}..., xr-hyper) is a
% relative URI, <file>#nameddest=<destination>, in the file-link colour:
% browser viewers (Edge, Chrome) ignore hyperref's default GoToR action.
\makeatletter
\renewcommand*{\hyper@linkfile}[3]{% anchor text, file name, destination
  \begingroup
    \ifx\\#3\\\def\CQ@target{#2}\else\def\CQ@target{#2\hyper@hash nameddest=#3}\fi
    \let\@urlcolor\@filecolor
    \hyper@linkurl{#1}{\CQ@target}%
  \endgroup}
\makeatother
%% --- 5. The mapped Unicode set (harness/contract_navigation.py, UNICODE_TEX)
\DeclareUnicodeCharacter{03B1}{\ensuremath{\alpha}}
\DeclareUnicodeCharacter{03B2}{\ensuremath{\beta}}
\DeclareUnicodeCharacter{03B3}{\ensuremath{\gamma}}
\DeclareUnicodeCharacter{03B4}{\ensuremath{\delta}}
\DeclareUnicodeCharacter{03B5}{\ensuremath{\varepsilon}}
\DeclareUnicodeCharacter{03F5}{\ensuremath{\epsilon}}
\DeclareUnicodeCharacter{03B6}{\ensuremath{\zeta}}
\DeclareUnicodeCharacter{03B7}{\ensuremath{\eta}}
\DeclareUnicodeCharacter{03B8}{\ensuremath{\theta}}
\DeclareUnicodeCharacter{03D1}{\ensuremath{\vartheta}}
\DeclareUnicodeCharacter{03B9}{\ensuremath{\iota}}
\DeclareUnicodeCharacter{03BA}{\ensuremath{\kappa}}
\DeclareUnicodeCharacter{03BB}{\ensuremath{\lambda}}
\DeclareUnicodeCharacter{03BC}{\ensuremath{\mu}}
\DeclareUnicodeCharacter{03BD}{\ensuremath{\nu}}
\DeclareUnicodeCharacter{03BE}{\ensuremath{\xi}}
\DeclareUnicodeCharacter{03C0}{\ensuremath{\pi}}
\DeclareUnicodeCharacter{03C1}{\ensuremath{\rho}}
\DeclareUnicodeCharacter{03F1}{\ensuremath{\varrho}}
\DeclareUnicodeCharacter{03C3}{\ensuremath{\sigma}}
\DeclareUnicodeCharacter{03C2}{\ensuremath{\varsigma}}
\DeclareUnicodeCharacter{03C4}{\ensuremath{\tau}}
\DeclareUnicodeCharacter{03C5}{\ensuremath{\upsilon}}
\DeclareUnicodeCharacter{03C6}{\ensuremath{\varphi}}
\DeclareUnicodeCharacter{03D5}{\ensuremath{\phi}}
\DeclareUnicodeCharacter{03C7}{\ensuremath{\chi}}
\DeclareUnicodeCharacter{03C8}{\ensuremath{\psi}}
\DeclareUnicodeCharacter{03C9}{\ensuremath{\omega}}
\DeclareUnicodeCharacter{0393}{\ensuremath{\Gamma}}
\DeclareUnicodeCharacter{0394}{\ensuremath{\Delta}}
\DeclareUnicodeCharacter{0398}{\ensuremath{\Theta}}
\DeclareUnicodeCharacter{039B}{\ensuremath{\Lambda}}
\DeclareUnicodeCharacter{039E}{\ensuremath{\Xi}}
\DeclareUnicodeCharacter{03A0}{\ensuremath{\Pi}}
\DeclareUnicodeCharacter{03A3}{\ensuremath{\Sigma}}
\DeclareUnicodeCharacter{03A5}{\ensuremath{\Upsilon}}
\DeclareUnicodeCharacter{03A6}{\ensuremath{\Phi}}
\DeclareUnicodeCharacter{03A8}{\ensuremath{\Psi}}
\DeclareUnicodeCharacter{03A9}{\ensuremath{\Omega}}
\DeclareUnicodeCharacter{00E1}{\'a}
\DeclareUnicodeCharacter{00E0}{\`a}
\DeclareUnicodeCharacter{00E2}{\^a}
\DeclareUnicodeCharacter{00E4}{\"a}
\DeclareUnicodeCharacter{00E5}{\aa{}}
\DeclareUnicodeCharacter{00E9}{\'e}
\DeclareUnicodeCharacter{00E8}{\`e}
\DeclareUnicodeCharacter{00EA}{\^e}
\DeclareUnicodeCharacter{00EB}{\"e}
\DeclareUnicodeCharacter{00ED}{\'{\i}}
\DeclareUnicodeCharacter{00EF}{\"{\i}}
\DeclareUnicodeCharacter{00F3}{\'o}
\DeclareUnicodeCharacter{00F2}{\`o}
\DeclareUnicodeCharacter{00F4}{\^o}
\DeclareUnicodeCharacter{00F6}{\"o}
\DeclareUnicodeCharacter{00F8}{\o{}}
\DeclareUnicodeCharacter{00FA}{\'u}
\DeclareUnicodeCharacter{00F9}{\`u}
\DeclareUnicodeCharacter{00FB}{\^u}
\DeclareUnicodeCharacter{00FC}{\"u}
\DeclareUnicodeCharacter{00F1}{\~n}
\DeclareUnicodeCharacter{00E7}{\c{c}}
\DeclareUnicodeCharacter{00DF}{\ss{}}
\DeclareUnicodeCharacter{00C4}{\"A}
\DeclareUnicodeCharacter{00C5}{\AA{}}
\DeclareUnicodeCharacter{00C9}{\'E}
\DeclareUnicodeCharacter{00D6}{\"O}
\DeclareUnicodeCharacter{00DC}{\"U}
\DeclareUnicodeCharacter{2013}{\textendash{}}
\DeclareUnicodeCharacter{2014}{\textemdash{}}
\DeclareUnicodeCharacter{2018}{`}
\DeclareUnicodeCharacter{2019}{'}
\DeclareUnicodeCharacter{201C}{``}
\DeclareUnicodeCharacter{201D}{''}
\DeclareUnicodeCharacter{2026}{\ldots{}}
\DeclareUnicodeCharacter{00A0}{~}
\DeclareUnicodeCharacter{00A7}{\S{}}
\DeclareUnicodeCharacter{2264}{\ensuremath{\le}}
\DeclareUnicodeCharacter{2265}{\ensuremath{\ge}}
\DeclareUnicodeCharacter{2260}{\ensuremath{\neq}}
\DeclareUnicodeCharacter{2248}{\ensuremath{\approx}}
\DeclareUnicodeCharacter{2261}{\ensuremath{\equiv}}
\DeclareUnicodeCharacter{223C}{\ensuremath{\sim}}
\DeclareUnicodeCharacter{2243}{\ensuremath{\simeq}}
\DeclareUnicodeCharacter{2245}{\ensuremath{\cong}}
\DeclareUnicodeCharacter{221D}{\ensuremath{\propto}}
\DeclareUnicodeCharacter{226A}{\ensuremath{\ll}}
\DeclareUnicodeCharacter{226B}{\ensuremath{\gg}}
\DeclareUnicodeCharacter{2192}{\ensuremath{\to}}
\DeclareUnicodeCharacter{2190}{\ensuremath{\leftarrow}}
\DeclareUnicodeCharacter{2194}{\ensuremath{\leftrightarrow}}
\DeclareUnicodeCharacter{21D2}{\ensuremath{\Rightarrow}}
\DeclareUnicodeCharacter{21D0}{\ensuremath{\Leftarrow}}
\DeclareUnicodeCharacter{21D4}{\ensuremath{\Leftrightarrow}}
\DeclareUnicodeCharacter{21A6}{\ensuremath{\mapsto}}
\DeclareUnicodeCharacter{2208}{\ensuremath{\in}}
\DeclareUnicodeCharacter{2209}{\ensuremath{\notin}}
\DeclareUnicodeCharacter{2286}{\ensuremath{\subseteq}}
\DeclareUnicodeCharacter{2282}{\ensuremath{\subset}}
\DeclareUnicodeCharacter{2287}{\ensuremath{\supseteq}}
\DeclareUnicodeCharacter{2283}{\ensuremath{\supset}}
\DeclareUnicodeCharacter{2229}{\ensuremath{\cap}}
\DeclareUnicodeCharacter{222A}{\ensuremath{\cup}}
\DeclareUnicodeCharacter{2216}{\ensuremath{\setminus}}
\DeclareUnicodeCharacter{2205}{\ensuremath{\emptyset}}
\DeclareUnicodeCharacter{2200}{\ensuremath{\forall}}
\DeclareUnicodeCharacter{2203}{\ensuremath{\exists}}
\DeclareUnicodeCharacter{2115}{\ensuremath{\mathbb{N}}}
\DeclareUnicodeCharacter{2124}{\ensuremath{\mathbb{Z}}}
\DeclareUnicodeCharacter{211A}{\ensuremath{\mathbb{Q}}}
\DeclareUnicodeCharacter{211D}{\ensuremath{\mathbb{R}}}
\DeclareUnicodeCharacter{2102}{\ensuremath{\mathbb{C}}}
\DeclareUnicodeCharacter{00B7}{\ensuremath{\cdot}}
\DeclareUnicodeCharacter{00D7}{\ensuremath{\times}}
\DeclareUnicodeCharacter{2212}{\ensuremath{-}}
\DeclareUnicodeCharacter{00B1}{\ensuremath{\pm}}
\DeclareUnicodeCharacter{2213}{\ensuremath{\mp}}
\DeclareUnicodeCharacter{221A}{\ensuremath{\surd}}
\DeclareUnicodeCharacter{2211}{\ensuremath{\sum}}
\DeclareUnicodeCharacter{220F}{\ensuremath{\prod}}
\DeclareUnicodeCharacter{222B}{\ensuremath{\int}}
\DeclareUnicodeCharacter{221E}{\ensuremath{\infty}}
\DeclareUnicodeCharacter{2202}{\ensuremath{\partial}}
\DeclareUnicodeCharacter{2207}{\ensuremath{\nabla}}
\DeclareUnicodeCharacter{2297}{\ensuremath{\otimes}}
\DeclareUnicodeCharacter{2295}{\ensuremath{\oplus}}
\DeclareUnicodeCharacter{2218}{\ensuremath{\circ}}
\DeclareUnicodeCharacter{2308}{\ensuremath{\lceil}}
\DeclareUnicodeCharacter{2309}{\ensuremath{\rceil}}
\DeclareUnicodeCharacter{230A}{\ensuremath{\lfloor}}
\DeclareUnicodeCharacter{230B}{\ensuremath{\rfloor}}
\DeclareUnicodeCharacter{27E8}{\ensuremath{\langle}}
\DeclareUnicodeCharacter{27E9}{\ensuremath{\rangle}}
\DeclareUnicodeCharacter{2016}{\ensuremath{\|}}
\DeclareUnicodeCharacter{2113}{\ensuremath{\ell}}
\DeclareUnicodeCharacter{00B0}{\ensuremath{^{\circ}}}
\DeclareUnicodeCharacter{00B2}{\ensuremath{^{2}}}
\DeclareUnicodeCharacter{00B3}{\ensuremath{^{3}}}
\DeclareUnicodeCharacter{00B9}{\ensuremath{^{1}}}
\DeclareUnicodeCharacter{2080}{\ensuremath{_{0}}}
\DeclareUnicodeCharacter{2081}{\ensuremath{_{1}}}
\DeclareUnicodeCharacter{2082}{\ensuremath{_{2}}}
\DeclareUnicodeCharacter{207F}{\ensuremath{^{n}}}
\DeclareUnicodeCharacter{2096}{\ensuremath{_{k}}}
%% --- 6. Numbering: one counter per section -----------------------
\numberwithin{equation}{section}
\setcounter{tocdepth}{1}
%% --- 7. The project macro set, then the claim environments -----------------
% spec/macros.tex, read relative to this module's directory (latexmk -cd).
% =============================================================================
%  spec/macros.tex -- the one macro set every generated contract inputs.
%
%  Two parts: part A, the macros spec/notation.md catalogues, and part B,
%  the macros local to one module, grouped by module.
%  harness/assemble_contract.py emits % =============================================================================
%  spec/macros.tex -- the one macro set every generated contract inputs.
%
%  Two parts: part A, the macros spec/notation.md catalogues, and part B,
%  the macros local to one module, grouped by module.
%  harness/assemble_contract.py emits % =============================================================================
%  spec/macros.tex -- the one macro set every generated contract inputs.
%
%  Two parts: part A, the macros spec/notation.md catalogues, and part B,
%  the macros local to one module, grouped by module.
%  harness/assemble_contract.py emits \input{../../macros} for it, so the
%  file is read from spec/modules/<id>/.
%
%  RULES
%    * A macro that spec/notation.md catalogues is a verbatim copy of it
%      (part A). Changing one is a notation change: bump spec/notation.md.
%    * Part B holds the module-local macros. A macro defined by two modules
%      with the same expansion appears once, under the first module (in the
%      order below) that defines it. Two macros with the same name and
%      DIFFERENT expansions are both kept under a module prefix, and the
%      claims write the prefixed names:
%        \Chain  -> \redChain (reductions, \mathfrak{C}: the reduction chain)
%                   \pcChain  (protocol-compiler, \mathfrak{R}: the two-segment
%                             chain; \Chsrc and \Chrec are built on it)
%        \Nmsm   -> \redNmsm  (reductions, \mathcal{N}: a POVM measurement
%                             channel)
%                   \pcNmsm   (\mathcal{N}_{\mathsf A}: Alice's measurement
%                             channel, written by
%                             transmitter/thm:certificate-sufficiency,
%                             transmitter/rem:xreg and
%                             reductions/lem:source-transfer)
%    * Deliberately omitted: a \provenance{tag}{details} macro. Provenance
%      is a header field of every claim file (spec/claim-format.md section
%      3); a claim body that calls such a macro fails to compile, which is
%      the intended signal.
%
%  REQUIRED PACKAGES, loaded by the assembler's preamble before this file:
%    amsmath, amssymb, amsfonts, mathtools, dsfont, braket, physics, xcolor
%    (braket before physics, so that physics wins for \ket, \bra, \tr, ...).
%
%  This file contains only definitions: no \usepackage, no document body.
% =============================================================================


% =============================================================================
%  PART A -- notation macros, the ones spec/notation.md catalogues, each a
%  verbatim copy of its catalogue row; grouped under numbered headings.
% =============================================================================

% -----------------------------------------------------------------------------
% 1. State spaces and operator sets (a "system label" -> spaces of operators)
% -----------------------------------------------------------------------------
\newcommand{\Hs}[1]{\mathcal{H}(#1)}        % Hilbert space of a system label
\newcommand{\Lin}[1]{\mathcal{L}(#1)}       % bounded linear operators on H(#1)
\newcommand{\Linn}[2]{\mathcal{L}(#1,#2)}   % linear maps H(#1) -> H(#2)
\newcommand{\Pos}[1]{\mathrm{Pos}(#1)}      % positive semidefinite operators on H(#1)
\newcommand{\Density}{\mathrm{D}}           % density operators: \Density(A) = D(A)
\newcommand{\Seq}[1]{S_{=}(#1)}             % normalized states (trace = 1)
\newcommand{\Sleq}[1]{S_{\leq}(#1)}         % subnormalized states (0 <= trace <= 1)
\newcommand{\id}{\mathds{1}}                % identity operator
\newcommand{\idon}[1]{\mathds{1}_{#1}}      % identity on a named system, e.g. \idon{C}

% DEPRECATED: applies \mathcal to both arguments itself, so it takes bare
% letters, not system labels. Prefer \Linn; do not use in new contracts.
\providecommand{\Lop}[2]{\mathcal{L}(\mathcal{#1}, \mathcal{#2})}


% -----------------------------------------------------------------------------
% 2. Composite systems
% -----------------------------------------------------------------------------
\newcommand{\dsum}{\oplus}                  % direct sum of system labels
% Tensor product: use \otimes directly.
% n-fold direct sum of an operator space: \Lin{A}^{\oplus n}


% -----------------------------------------------------------------------------
% 3. Events, conditioning and classical registers (for STATES)
%    Convention:  \wedge  -> subnormalized "and-event" component, e.g. \rho_{A\wedge x}
%                 |        -> normalized conditional state,        e.g. \rho_{A|x}
% (Write these directly in subscripts; no extra macros needed.)
% -----------------------------------------------------------------------------


% -----------------------------------------------------------------------------
% 4. Photon-number subspaces / direct-sum blocks
%    Convention:  operators carry the restriction "on top", in brackets;
%                 states carry it as a \wedge subscript (see Sec. 3).
% -----------------------------------------------------------------------------
\newcommand{\Pnum}[1]{P_{#1}}               % photon-number projector (system from context)
\newcommand{\Pnumon}[2]{P_{#1}^{#2}}        % with explicit system superscript, e.g. P_n^C
\newcommand{\blk}[2]{#1^{(#2)}}             % operator block, e.g. \blk{M}{\leqN} = M^{(<=N)}
\newcommand{\leqN}{\leq N}                  % use as \blk{M}{\leqN}
\newcommand{\gtN}{> N}                      % use as \blk{M}{\gtN}
\newcommand{\eqn}{= n}                      % use as \blk{M}{\eqn}


% -----------------------------------------------------------------------------
% 5. System labels used in the protocol
% -----------------------------------------------------------------------------
\newcommand{\Chan}{C}                       % quantum-channel (Eve) input to Bob
\newcommand{\Prb}{P}                        % probe input
\newcommand{\Flag}{{\mathrm{Flg}}}          % flag system (EXISTS ONLY AFTER squashing)
\newcommand{\Loss}{{\mathrm{Lss}}}          % loss system
\newcommand{\Vc}{{\mathrm{Vcm}}}            % vacuum system
\newcommand{\Nout}{N_{\mathrm{out}}}
\newcommand{\Nin}{N_{\mathrm{in}}}


% -----------------------------------------------------------------------------
% 6. Alphabets
% -----------------------------------------------------------------------------
\newcommand{\Yset}{\mathcal{Y}}             % Bob's (coarse-grained) detection alphabet
\newcommand{\Yraw}{\mathcal{Y}_{\mathrm{raw}}}% raw click patterns {0,1}^4
\newcommand{\Yscset}{\mathcal{Y}_{\singlec}}% single-click outcomes
\newcommand{\Ymcset}{\mathcal{Y}_{\mc}}     % multi-click outcomes
\newcommand{\Tset}{\mathcal{T}_{\mathrm{set}}}% private setting set (the probe settings in the optical instance)
\newcommand{\Xset}{\mathcal{X}}             % Alice's alphabet
\newcommand{\Dset}{\{H,V,+,-\}}             % polarization / detector labels
\newcommand{\pass}{\top}                    % pass sifting


% -----------------------------------------------------------------------------
% 7. Mode / outcome abbreviations
% -----------------------------------------------------------------------------
\newcommand{\on}{{\mathrm{on}}}
\newcommand{\off}{{\mathrm{off}}}
\newcommand{\dc}{{\mathrm{dc}}}             % dark count
\newcommand{\sq}{\mathrm{sq}}               % squashed
\newcommand{\mc}{\mathrm{mc}}               % multi-click
\newcommand{\nc}{\mathrm{nc}}               % no-click
\newcommand{\singlec}{\mathrm{sc}}          % single-click
\newcommand{\vac}{\mathrm{vac}}             % vacuum
\newcommand{\kout}{{k_\mathrm{out}}}
\newcommand{\kin}{{k_\mathrm{in}}}


% -----------------------------------------------------------------------------
% 8. Maps, channels and miscellaneous operators
% -----------------------------------------------------------------------------
\newcommand{\Prob}{\mathrm{Prob}}           % probability (use \Prob[\cdot] or \Prob(\cdot))
\newcommand{\lmin}{\lambda_{\min}}          % minimum eigenvalue
\newcommand{\Wflag}{W_{\Flag}}              % flag-space weight
\newcommand{\transpose}{\mathsf{T}}
\newcommand{\innerp}[2]{\left\langle #1 \vert #2 \right\rangle}
\renewcommand{\vec}[1]{\boldsymbol{#1}}     % NB: overrides the LaTeX arrow accent


% -----------------------------------------------------------------------------
% 9. Key-rate numerics
% -----------------------------------------------------------------------------
\newcommand{\cG}{\mathcal{G}}               % protocol map
\newcommand{\cZ}{\mathcal{Z}}               % pinching map
\newcommand{\Amap}{\mathcal{A}_{\mathrm{ann}}} % announcement + measurement map
\newcommand{\Amapnb}{\mathcal{A}_{\mathrm{ann}|\mathrm{nb}}}
\newcommand{\simpl}{}                       % empty suffix for simplified block maps
\newcommand{\acSet}{\mathcal{S}_{\mathrm{acc}}} % acceptance set
\newcommand{\testE}{\mathcal{E}_{\mathrm{test}}} % test-statistics map
\newcommand{\ecLeak}{\mathrm{leak}^{\mathrm{EC}}_{\mathrm{obs}}}
\newcommand{\pgen}{p_{\mathrm{gen}}}
\newcommand{\askeyr}{R_{\infty}}             % asymptotic key rate
\newcommand{\ppass}{p_{\mathrm{pass}}}       % sifting pass probability


% -----------------------------------------------------------------------------
% 10. Finite-size analysis / Renyi entropies
% -----------------------------------------------------------------------------

% --- Round types: test vs. key generation (register T, alphabet \Fset) --------
\newcommand{\tst}{\mathrm{test}}            % test round
\newcommand{\gen}{\mathrm{gen}}             % key-generation round
\newcommand{\Fset}{\mathcal{T}}             % round-type alphabet {test, gen} (not \Tset)
\newcommand{\Chat}{\widehat{C}}             % acceptance-test register
\newcommand{\Chatset}{\widehat{\mathcal{C}}}% its alphabet

% --- Basis / key extraction functions ----------------------------------------
\newcommand{\bA}{b_{A}}                     % Alice's basis-extraction function
\newcommand{\bB}{b_{B}}                     % Bob's basis-extraction function
\newcommand{\kB}{k_{B}}                     % Bob's key-extraction function

% --- Maps --------------------------------------------------------------------
\newcommand{\Mprot}{\mathcal{M}}            % single-round protocol map
\newcommand{\Phimap}{\boldsymbol{\Phi}}     % test-statistics map
\newcommand{\Pdeph}{\mathcal{P}}            % dephasing (pinching) map on \Blreg

% --- Isometries --------------------------------------------------------------
\newcommand{\Vmsm}{V_{\mathrm{msm}}}        % measurement isometry
\newcommand{\Vann}{V_{\mathrm{ann}}}        % announcement isometry
\newcommand{\VI}{V_{I}}                     % Eve's copy of the announcements
\newcommand{\VChat}{V_{\Chat}}              % acceptance-data isometry
\newcommand{\VR}{V_{R}}                     % sifting + key map -> raw key R
\newcommand{\VZ}{V_{Z}}                     % raw key -> key register Z
\newcommand{\Wprot}{W}                      % full protocol isometry

% --- Renyi entropies ---------------------------------------------------------
% SECURITY-CRITICAL: see spec/notation.md §4. Never change an arrow to
% make an inequality work out; never convert silently between the sandwiched
% (tilde) and Petz (overline) families.
\newcommand{\Hsandup}{\widetilde{H}^{\uparrow}_{\alpha}}   % sandwiched, up
\newcommand{\Hsanddown}{\widetilde{H}^{\downarrow}_{\alpha}}% sandwiched, down
\newcommand{\Hpetzup}{\overline{H}^{\uparrow}_{\alpha}}    % Petz, up
\newcommand{\Hpetzdown}{\overline{H}^{\downarrow}_{\alpha}}% Petz, down
\newcommand{\hup}{h^{\uparrow}_{\alpha}}                   % single-round rate
\newcommand{\hqkd}{h^{\uparrow,\mathrm{QKD}}_{\alpha}}     % single-round QKD rate
\newcommand{\hlow}{\underline{h}_{\alpha}}                 % lower bound on \hqkd
% --- spec/notation.md 3.5 and 4.5 -------------------------------------------
% \hlb  is the first-order linearization lower bound on \hlow: the headline
%       key rate (the `first-order` tier of foundations/def:tiers). NOT \hlow
%       itself.
% \tnum is the instance's one declared numerical tolerance. Its
%       near-collision with the declared marginal \tau_A is disclosed in
%       spec/notation.md 3.5; bare \tau is reserved and denotes neither.
\newcommand{\hlb}{\underline{h}^{\mathrm{FO}}_{\alpha}}    % first-order lower bound on \hlow
\newcommand{\tnum}{\tau_{\mathrm{num}}}                    % declared numerical tolerance


% =============================================================================
%  PART B -- module-local macros, grouped by module. Deprecated aliases,
%  defined so that a text written with them compiles, used in no new contract: \Lop,
%  \testE, \Pdeph, \pass, \ecLeak (spec/notation.md section 3.9).
% =============================================================================

% --- meat-core -----------------------------------------------------------------
\newcommand{\SecSet}{\mathbb{S}}            % sector index set of Q
\newcommand{\Qsys}{Q}                       % declared sector-carrying measured system
\newcommand{\Vset}{\mathcal{V}}             % declared joint outcome alphabet of the Q-POVM
\newcommand{\CAset}{\mathcal{C}_{A}}        % Alice announcement alphabet
\newcommand{\CBset}{\mathcal{C}_{B}}        % Bob announcement alphabet
\newcommand{\Party}{\mathfrak{P}}           % declared party partition
\newcommand{\Oset}[1]{\mathcal{O}_{#1}}     % readout alphabet of party #1
\newcommand{\rd}[1]{\pi_{#1}}               % readout map of party #1
\newcommand{\Sadm}{\mathcal{S}_{A}}         % admissible Alice subspace
\newcommand{\nomarg}{\mathsf{none}}         % "no trusted marginal declared"
\newcommand{\Zset}{\mathcal{Z}_{\perp}}     % raw-key alphabet incl. \perp
\newcommand{\Zval}{\mathcal{Z}}             % key alphabet without \perp
\newcommand{\Etil}{\widetilde{E}}           % single-round purifying register
\newcommand{\Ehat}{\widehat{E}}             % Eve's n-round register
\newcommand{\OmAT}{\Omega_{\mathrm{AT}}}    % acceptance event
\newcommand{\epa}{\varepsilon_{\mathrm{PA}}}
\newcommand{\eev}{\varepsilon_{\mathrm{EV}}}
\newcommand{\lEC}{\lambda_{\mathrm{EC}}}
\newcommand{\Decl}{\mathfrak{D}}            % protocol declaration
\newcommand{\CList}{\mathfrak{L}}           % declared constraint list
\newcommand{\PinS}{\mathcal{P}_{\SecSet}}   % sector pinching on Q
\newcommand{\Feas}{\mathcal{W}}             % feasible set of block states
\newcommand{\Chkey}{\Chatset_{\mathrm{key}}}% key-live acceptance values
\newcommand{\Kann}{\mathcal{K}}             % announcement-class alphabet
\newcommand{\hold}{\mathsf{h}}              % declared key holder (A or B)
\newcommand{\sfD}{D}                        % classical KL divergence (base 2)
\newcommand{\Dtwo}{D}                       % quantum relative entropy (base 2)
\newcommand{\Dsand}{\widetilde{D}_{\alpha}} % sandwiched Renyi divergence
\newcommand{\idmap}{\mathrm{id}}            % identity channel
\newcommand{\Aev}{\mathcal{A}}              % Eve's per-round attack channel
\newcommand{\AttSet}{\mathfrak{A}}          % licensed attack-channel set
\newcommand{\Chacc}{\Chatset_{\mathrm{acc}}}% acceptance support of \acSet
\newcommand{\epsstar}{\varepsilon^{\star}}      % distance to an exact POVM
\newcommand{\dfH}{\Delta_{\mathrm{herm}}}       % Hermiticity defect
\newcommand{\dfP}{\Delta_{\mathrm{pos}}}        % negative part
\newcommand{\dfC}{\Delta_{\mathrm{cpl}}}        % completeness defect
\newcommand{\lmoq}{g}                       % acceptance-side gradient block
\newcommand{\lmoJ}{G}                       % feasible-side gradient block
\newcommand{\bacc}{\beta_{\mathrm{acc}}}    % certified acceptance-side minimum
\newcommand{\Itset}{\mathfrak{I}}           % finite index set of iterates (C22)

% --- reductions ----------------------------------------------------------------
\newcommand{\Bsys}{B}                        % physical (pre-squashing) system
\newcommand{\Tcl}[1]{\mathcal{T}(#1)}        % trace-class operators on H(#1)
\newcommand{\Ins}{\mathcal{I}}               % generalized-instrument branch map
\newcommand{\redNmsm}{\mathcal{N}}              % POVM measurement channel
\newcommand{\Eprot}{\mathcal{E}}             % n-round protocol channel
\newcommand{\Grnd}{\mathsf{G}}               % one round map of a protocol
\newcommand{\Kfin}{\mathsf{K}}               % end-of-protocol processing
\newcommand{\Crit}{\mathsf{sec}}             % value of a security criterion
\newcommand{\Pflag}{\Pi_{\Flag}}             % flag-sector projector on Q
\newcommand{\Pleq}{\Pi_{\leqN}}              % preserved-block projector
\newcommand{\Pgt}{\Pi_{\gtN}}                % above-cutoff projector
\newcommand{\Sph}{S^{\mathrm{ph}}}           % family of physical inputs
\newcommand{\Srch}{S^{\mathrm{rch}}}         % reachable squashed states
\newcommand{\Ssq}{S^{\sq}}                   % squashed admissible set
\newcommand{\Ffss}{F^{\mathrm{FSS}}}         % flag-state-squasher target
\newcommand{\Fwp}{F^{\mathrm{WP}}}           % multi-bound WPFSS target
\newcommand{\Mwp}{M^{\mathrm{WP}}}           % renormalized above-cutoff POVM
\newcommand{\Mdec}{{M'}^{\mathrm{ph}}}       % declared physical family on B
\newcommand{\Ppvm}{P^{\Flag}}                % flag PVM
\newcommand{\Lmp}{{\Lambda_{\mathrm{mp}}}}   % measure-and-prepare channel
\newcommand{\Lwp}{{\Lambda^{\mathrm{WP}}}}   % multi-bound WPFSS channel
\newcommand{\Lfss}{{\Lambda^{\mathrm{FSS}}}} % flag-state-squasher channel
\newcommand{\lmax}{\lambda_{\max}}           % maximum eigenvalue
\newcommand{\Lnoise}{\widetilde{\Lambda}}    % flag-noise channel
\newcommand{\cg}{\varpi}                     % declared coarse-graining map
\newcommand{\MCK}{\mathsf{MC}}               % multiclick raw-pattern set
\newcommand{\redChain}{\mathfrak{C}}            % reduction chain
\newcommand{\Lnk}{\ell}                      % a single typed link
\newcommand{\dsq}{\preceq}                      % the squashing relation
\newcommand{\opn}[1]{\left\lVert #1\right\rVert}           % operator (spectral) norm
\newcommand{\mxn}[1]{\left\lVert #1\right\rVert_{\max}}    % entrywise max modulus
\newcommand{\Fld}{\mathbb{F}}                   % exact ordered substrate field
\newcommand{\Fldi}{\mathbb{F}_{i}}              % its Gaussian extension F + iF
\newcommand{\lat}[1]{\big\lfloor #1\big\rfloor_{\kappa}} % round DOWN onto the lattice

% --- protocol-compiler ---------------------------------------------------------
\newcommand{\Spec}{\mathfrak{S}}             % protocol specification (compiler input)
\newcommand{\pcChain}{\mathfrak{R}}            % reduction chain (two typed segments)
\newcommand{\Chsrc}{\pcChain_{\mathrm{src}}}   % its Alice-side source-map segment
\newcommand{\Chrec}{\pcChain_{\mathrm{rec}}}   % its receiver-side squashing segment
\newcommand{\rvirt}{\rho^{\mathrm{v}}}       % virtual (PRE-source-map) signal state
\newcommand{\Fph}{F^{\mathrm{ph}}}           % physical (pre-reduction) POVM element
\newcommand{\Vph}{\mathcal{V}^{\mathrm{ph}}} % physical outcome alphabet
\newcommand{\Tpub}{\mathcal{T}_{\mathrm{pub}}}% public-choice alphabet
\newcommand{\Uset}{\mathcal{U}}              % residual outcome alphabet
\newcommand{\idxm}{\mathsf{idx}}             % label-to-index map
\newcommand{\Asig}{A'}                       % transmitted (signal) system -- the REAL one
\newcommand{\Avirt}{A''}                     % virtual transmitted system (source descriptor of Xi_1)
\newcommand{\Abar}{\bar{A}}                  % Alice's measured register
\newcommand{\Ahat}{\hat{A}}                  % Alice's shield register
\newcommand{\pcNmsm}{\mathcal{N}_{\mathsf A}}  % Alice's measurement channel (see the RULES block)
\newcommand{\Phip}{\Phi^{+}}                 % maximally entangled vector label
\newcommand{\Smap}{\mathfrak{s}}             % the emitted sector map (fraktur: a compile-time object)

% --- receiver ------------------------------------------------------------------
\newcommand{\Modes}{\mathbb{M}}              % declared ordered mode set
\newcommand{\Inp}{\mathbb{I}}                % declared input modes
\newcommand{\InpC}{\mathbb{I}_{\Chan}}       % channel input modes
\newcommand{\InpS}{\mathbb{I}_{\Prb}}        % probe (source) input modes
\newcommand{\Env}{\mathbb{E}}                % dilation (environment) modes
\newcommand{\Port}[1]{\mathbb{P}_{#1}}       % mode set of detector port #1
\newcommand{\Emd}[1]{E_{#1}}                 % coordinate projector on modes #1
\newcommand{\Gmat}{G}                        % instrument (transfer) matrix
\newcommand{\Fock}{\mathfrak{F}}             % Fock / second-quantization functor
\newcommand{\Symn}[1]{\mathrm{Sym}^{#1}}     % symmetric tensor power
\newcommand{\normord}[1]{{:}\mspace{2mu}#1\mspace{2mu}{:}}   % normal ordering
\newcommand{\rk}{\operatorname{rank}}        % rank (physics.sty owns \rank)
\newcommand{\SCK}{\mathsf{SC}}               % single-click pattern set
\newcommand{\Recv}{\mathfrak{R}}             % receiver instrument tuple
\newcommand{\Cat}{\mathfrak{C}}              % component catalogue
\newcommand{\Grph}{\mathfrak{E}}             % declared element graph
\newcommand{\defic}{\mathbb{L}}              % loss deficiency operator
\newcommand{\Pvac}[1]{\Pi_{\vac}(#1)}        % joint no-click element on ports #1
\newcommand{\pmc}{p_{\mc}}                   % uniform high-photon mc weight
\newcommand{\Bnd}{\beta}                     % the scalar bound function beta(n)
\newcommand{\Rot}{R}                         % planar rotation matrix
\newcommand{\Wplate}{W}                      % retarder (wave-plate) matrix
\newcommand{\Bsp}{B}                         % beamsplitter core matrix

% --- statistics ----------------------------------------------------------------
\newcommand{\honest}{p^{\mathrm{hon}}}          % honest per-round distribution
\newcommand{\ecomp}{\varepsilon_{\mathrm{comp}}}% completeness budget
\newcommand{\simplex}[1]{\Delta(#1)}            % probability simplex
\newcommand{\Chatcon}{\Chatset_{\mathrm{con}}}  % constrained outcomes
\newcommand{\Ireg}[3]{\mathcal{I}_{#1}(#2,#3)}  % regularized incomplete beta
\newcommand{\Ncount}[1]{\mathsf{N}_{#1}}        % per-outcome count
\newcommand{\qlo}{\mathbf{q}^{-}}               % box lower bounds
\newcommand{\qhi}{\mathbf{q}^{+}}               % box upper bounds
 for it, so the
%  file is read from spec/modules/<id>/.
%
%  RULES
%    * A macro that spec/notation.md catalogues is a verbatim copy of it
%      (part A). Changing one is a notation change: bump spec/notation.md.
%    * Part B holds the module-local macros. A macro defined by two modules
%      with the same expansion appears once, under the first module (in the
%      order below) that defines it. Two macros with the same name and
%      DIFFERENT expansions are both kept under a module prefix, and the
%      claims write the prefixed names:
%        \Chain  -> \redChain (reductions, \mathfrak{C}: the reduction chain)
%                   \pcChain  (protocol-compiler, \mathfrak{R}: the two-segment
%                             chain; \Chsrc and \Chrec are built on it)
%        \Nmsm   -> \redNmsm  (reductions, \mathcal{N}: a POVM measurement
%                             channel)
%                   \pcNmsm   (\mathcal{N}_{\mathsf A}: Alice's measurement
%                             channel, written by
%                             transmitter/thm:certificate-sufficiency,
%                             transmitter/rem:xreg and
%                             reductions/lem:source-transfer)
%    * Deliberately omitted: a \provenance{tag}{details} macro. Provenance
%      is a header field of every claim file (spec/claim-format.md section
%      3); a claim body that calls such a macro fails to compile, which is
%      the intended signal.
%
%  REQUIRED PACKAGES, loaded by the assembler's preamble before this file:
%    amsmath, amssymb, amsfonts, mathtools, dsfont, braket, physics, xcolor
%    (braket before physics, so that physics wins for \ket, \bra, \tr, ...).
%
%  This file contains only definitions: no \usepackage, no document body.
% =============================================================================


% =============================================================================
%  PART A -- notation macros, the ones spec/notation.md catalogues, each a
%  verbatim copy of its catalogue row; grouped under numbered headings.
% =============================================================================

% -----------------------------------------------------------------------------
% 1. State spaces and operator sets (a "system label" -> spaces of operators)
% -----------------------------------------------------------------------------
\newcommand{\Hs}[1]{\mathcal{H}(#1)}        % Hilbert space of a system label
\newcommand{\Lin}[1]{\mathcal{L}(#1)}       % bounded linear operators on H(#1)
\newcommand{\Linn}[2]{\mathcal{L}(#1,#2)}   % linear maps H(#1) -> H(#2)
\newcommand{\Pos}[1]{\mathrm{Pos}(#1)}      % positive semidefinite operators on H(#1)
\newcommand{\Density}{\mathrm{D}}           % density operators: \Density(A) = D(A)
\newcommand{\Seq}[1]{S_{=}(#1)}             % normalized states (trace = 1)
\newcommand{\Sleq}[1]{S_{\leq}(#1)}         % subnormalized states (0 <= trace <= 1)
\newcommand{\id}{\mathds{1}}                % identity operator
\newcommand{\idon}[1]{\mathds{1}_{#1}}      % identity on a named system, e.g. \idon{C}

% DEPRECATED: applies \mathcal to both arguments itself, so it takes bare
% letters, not system labels. Prefer \Linn; do not use in new contracts.
\providecommand{\Lop}[2]{\mathcal{L}(\mathcal{#1}, \mathcal{#2})}


% -----------------------------------------------------------------------------
% 2. Composite systems
% -----------------------------------------------------------------------------
\newcommand{\dsum}{\oplus}                  % direct sum of system labels
% Tensor product: use \otimes directly.
% n-fold direct sum of an operator space: \Lin{A}^{\oplus n}


% -----------------------------------------------------------------------------
% 3. Events, conditioning and classical registers (for STATES)
%    Convention:  \wedge  -> subnormalized "and-event" component, e.g. \rho_{A\wedge x}
%                 |        -> normalized conditional state,        e.g. \rho_{A|x}
% (Write these directly in subscripts; no extra macros needed.)
% -----------------------------------------------------------------------------


% -----------------------------------------------------------------------------
% 4. Photon-number subspaces / direct-sum blocks
%    Convention:  operators carry the restriction "on top", in brackets;
%                 states carry it as a \wedge subscript (see Sec. 3).
% -----------------------------------------------------------------------------
\newcommand{\Pnum}[1]{P_{#1}}               % photon-number projector (system from context)
\newcommand{\Pnumon}[2]{P_{#1}^{#2}}        % with explicit system superscript, e.g. P_n^C
\newcommand{\blk}[2]{#1^{(#2)}}             % operator block, e.g. \blk{M}{\leqN} = M^{(<=N)}
\newcommand{\leqN}{\leq N}                  % use as \blk{M}{\leqN}
\newcommand{\gtN}{> N}                      % use as \blk{M}{\gtN}
\newcommand{\eqn}{= n}                      % use as \blk{M}{\eqn}


% -----------------------------------------------------------------------------
% 5. System labels used in the protocol
% -----------------------------------------------------------------------------
\newcommand{\Chan}{C}                       % quantum-channel (Eve) input to Bob
\newcommand{\Prb}{P}                        % probe input
\newcommand{\Flag}{{\mathrm{Flg}}}          % flag system (EXISTS ONLY AFTER squashing)
\newcommand{\Loss}{{\mathrm{Lss}}}          % loss system
\newcommand{\Vc}{{\mathrm{Vcm}}}            % vacuum system
\newcommand{\Nout}{N_{\mathrm{out}}}
\newcommand{\Nin}{N_{\mathrm{in}}}


% -----------------------------------------------------------------------------
% 6. Alphabets
% -----------------------------------------------------------------------------
\newcommand{\Yset}{\mathcal{Y}}             % Bob's (coarse-grained) detection alphabet
\newcommand{\Yraw}{\mathcal{Y}_{\mathrm{raw}}}% raw click patterns {0,1}^4
\newcommand{\Yscset}{\mathcal{Y}_{\singlec}}% single-click outcomes
\newcommand{\Ymcset}{\mathcal{Y}_{\mc}}     % multi-click outcomes
\newcommand{\Tset}{\mathcal{T}_{\mathrm{set}}}% private setting set (the probe settings in the optical instance)
\newcommand{\Xset}{\mathcal{X}}             % Alice's alphabet
\newcommand{\Dset}{\{H,V,+,-\}}             % polarization / detector labels
\newcommand{\pass}{\top}                    % pass sifting


% -----------------------------------------------------------------------------
% 7. Mode / outcome abbreviations
% -----------------------------------------------------------------------------
\newcommand{\on}{{\mathrm{on}}}
\newcommand{\off}{{\mathrm{off}}}
\newcommand{\dc}{{\mathrm{dc}}}             % dark count
\newcommand{\sq}{\mathrm{sq}}               % squashed
\newcommand{\mc}{\mathrm{mc}}               % multi-click
\newcommand{\nc}{\mathrm{nc}}               % no-click
\newcommand{\singlec}{\mathrm{sc}}          % single-click
\newcommand{\vac}{\mathrm{vac}}             % vacuum
\newcommand{\kout}{{k_\mathrm{out}}}
\newcommand{\kin}{{k_\mathrm{in}}}


% -----------------------------------------------------------------------------
% 8. Maps, channels and miscellaneous operators
% -----------------------------------------------------------------------------
\newcommand{\Prob}{\mathrm{Prob}}           % probability (use \Prob[\cdot] or \Prob(\cdot))
\newcommand{\lmin}{\lambda_{\min}}          % minimum eigenvalue
\newcommand{\Wflag}{W_{\Flag}}              % flag-space weight
\newcommand{\transpose}{\mathsf{T}}
\newcommand{\innerp}[2]{\left\langle #1 \vert #2 \right\rangle}
\renewcommand{\vec}[1]{\boldsymbol{#1}}     % NB: overrides the LaTeX arrow accent


% -----------------------------------------------------------------------------
% 9. Key-rate numerics
% -----------------------------------------------------------------------------
\newcommand{\cG}{\mathcal{G}}               % protocol map
\newcommand{\cZ}{\mathcal{Z}}               % pinching map
\newcommand{\Amap}{\mathcal{A}_{\mathrm{ann}}} % announcement + measurement map
\newcommand{\Amapnb}{\mathcal{A}_{\mathrm{ann}|\mathrm{nb}}}
\newcommand{\simpl}{}                       % empty suffix for simplified block maps
\newcommand{\acSet}{\mathcal{S}_{\mathrm{acc}}} % acceptance set
\newcommand{\testE}{\mathcal{E}_{\mathrm{test}}} % test-statistics map
\newcommand{\ecLeak}{\mathrm{leak}^{\mathrm{EC}}_{\mathrm{obs}}}
\newcommand{\pgen}{p_{\mathrm{gen}}}
\newcommand{\askeyr}{R_{\infty}}             % asymptotic key rate
\newcommand{\ppass}{p_{\mathrm{pass}}}       % sifting pass probability


% -----------------------------------------------------------------------------
% 10. Finite-size analysis / Renyi entropies
% -----------------------------------------------------------------------------

% --- Round types: test vs. key generation (register T, alphabet \Fset) --------
\newcommand{\tst}{\mathrm{test}}            % test round
\newcommand{\gen}{\mathrm{gen}}             % key-generation round
\newcommand{\Fset}{\mathcal{T}}             % round-type alphabet {test, gen} (not \Tset)
\newcommand{\Chat}{\widehat{C}}             % acceptance-test register
\newcommand{\Chatset}{\widehat{\mathcal{C}}}% its alphabet

% --- Basis / key extraction functions ----------------------------------------
\newcommand{\bA}{b_{A}}                     % Alice's basis-extraction function
\newcommand{\bB}{b_{B}}                     % Bob's basis-extraction function
\newcommand{\kB}{k_{B}}                     % Bob's key-extraction function

% --- Maps --------------------------------------------------------------------
\newcommand{\Mprot}{\mathcal{M}}            % single-round protocol map
\newcommand{\Phimap}{\boldsymbol{\Phi}}     % test-statistics map
\newcommand{\Pdeph}{\mathcal{P}}            % dephasing (pinching) map on \Blreg

% --- Isometries --------------------------------------------------------------
\newcommand{\Vmsm}{V_{\mathrm{msm}}}        % measurement isometry
\newcommand{\Vann}{V_{\mathrm{ann}}}        % announcement isometry
\newcommand{\VI}{V_{I}}                     % Eve's copy of the announcements
\newcommand{\VChat}{V_{\Chat}}              % acceptance-data isometry
\newcommand{\VR}{V_{R}}                     % sifting + key map -> raw key R
\newcommand{\VZ}{V_{Z}}                     % raw key -> key register Z
\newcommand{\Wprot}{W}                      % full protocol isometry

% --- Renyi entropies ---------------------------------------------------------
% SECURITY-CRITICAL: see spec/notation.md §4. Never change an arrow to
% make an inequality work out; never convert silently between the sandwiched
% (tilde) and Petz (overline) families.
\newcommand{\Hsandup}{\widetilde{H}^{\uparrow}_{\alpha}}   % sandwiched, up
\newcommand{\Hsanddown}{\widetilde{H}^{\downarrow}_{\alpha}}% sandwiched, down
\newcommand{\Hpetzup}{\overline{H}^{\uparrow}_{\alpha}}    % Petz, up
\newcommand{\Hpetzdown}{\overline{H}^{\downarrow}_{\alpha}}% Petz, down
\newcommand{\hup}{h^{\uparrow}_{\alpha}}                   % single-round rate
\newcommand{\hqkd}{h^{\uparrow,\mathrm{QKD}}_{\alpha}}     % single-round QKD rate
\newcommand{\hlow}{\underline{h}_{\alpha}}                 % lower bound on \hqkd
% --- spec/notation.md 3.5 and 4.5 -------------------------------------------
% \hlb  is the first-order linearization lower bound on \hlow: the headline
%       key rate (the `first-order` tier of foundations/def:tiers). NOT \hlow
%       itself.
% \tnum is the instance's one declared numerical tolerance. Its
%       near-collision with the declared marginal \tau_A is disclosed in
%       spec/notation.md 3.5; bare \tau is reserved and denotes neither.
\newcommand{\hlb}{\underline{h}^{\mathrm{FO}}_{\alpha}}    % first-order lower bound on \hlow
\newcommand{\tnum}{\tau_{\mathrm{num}}}                    % declared numerical tolerance


% =============================================================================
%  PART B -- module-local macros, grouped by module. Deprecated aliases,
%  defined so that a text written with them compiles, used in no new contract: \Lop,
%  \testE, \Pdeph, \pass, \ecLeak (spec/notation.md section 3.9).
% =============================================================================

% --- meat-core -----------------------------------------------------------------
\newcommand{\SecSet}{\mathbb{S}}            % sector index set of Q
\newcommand{\Qsys}{Q}                       % declared sector-carrying measured system
\newcommand{\Vset}{\mathcal{V}}             % declared joint outcome alphabet of the Q-POVM
\newcommand{\CAset}{\mathcal{C}_{A}}        % Alice announcement alphabet
\newcommand{\CBset}{\mathcal{C}_{B}}        % Bob announcement alphabet
\newcommand{\Party}{\mathfrak{P}}           % declared party partition
\newcommand{\Oset}[1]{\mathcal{O}_{#1}}     % readout alphabet of party #1
\newcommand{\rd}[1]{\pi_{#1}}               % readout map of party #1
\newcommand{\Sadm}{\mathcal{S}_{A}}         % admissible Alice subspace
\newcommand{\nomarg}{\mathsf{none}}         % "no trusted marginal declared"
\newcommand{\Zset}{\mathcal{Z}_{\perp}}     % raw-key alphabet incl. \perp
\newcommand{\Zval}{\mathcal{Z}}             % key alphabet without \perp
\newcommand{\Etil}{\widetilde{E}}           % single-round purifying register
\newcommand{\Ehat}{\widehat{E}}             % Eve's n-round register
\newcommand{\OmAT}{\Omega_{\mathrm{AT}}}    % acceptance event
\newcommand{\epa}{\varepsilon_{\mathrm{PA}}}
\newcommand{\eev}{\varepsilon_{\mathrm{EV}}}
\newcommand{\lEC}{\lambda_{\mathrm{EC}}}
\newcommand{\Decl}{\mathfrak{D}}            % protocol declaration
\newcommand{\CList}{\mathfrak{L}}           % declared constraint list
\newcommand{\PinS}{\mathcal{P}_{\SecSet}}   % sector pinching on Q
\newcommand{\Feas}{\mathcal{W}}             % feasible set of block states
\newcommand{\Chkey}{\Chatset_{\mathrm{key}}}% key-live acceptance values
\newcommand{\Kann}{\mathcal{K}}             % announcement-class alphabet
\newcommand{\hold}{\mathsf{h}}              % declared key holder (A or B)
\newcommand{\sfD}{D}                        % classical KL divergence (base 2)
\newcommand{\Dtwo}{D}                       % quantum relative entropy (base 2)
\newcommand{\Dsand}{\widetilde{D}_{\alpha}} % sandwiched Renyi divergence
\newcommand{\idmap}{\mathrm{id}}            % identity channel
\newcommand{\Aev}{\mathcal{A}}              % Eve's per-round attack channel
\newcommand{\AttSet}{\mathfrak{A}}          % licensed attack-channel set
\newcommand{\Chacc}{\Chatset_{\mathrm{acc}}}% acceptance support of \acSet
\newcommand{\epsstar}{\varepsilon^{\star}}      % distance to an exact POVM
\newcommand{\dfH}{\Delta_{\mathrm{herm}}}       % Hermiticity defect
\newcommand{\dfP}{\Delta_{\mathrm{pos}}}        % negative part
\newcommand{\dfC}{\Delta_{\mathrm{cpl}}}        % completeness defect
\newcommand{\lmoq}{g}                       % acceptance-side gradient block
\newcommand{\lmoJ}{G}                       % feasible-side gradient block
\newcommand{\bacc}{\beta_{\mathrm{acc}}}    % certified acceptance-side minimum
\newcommand{\Itset}{\mathfrak{I}}           % finite index set of iterates (C22)

% --- reductions ----------------------------------------------------------------
\newcommand{\Bsys}{B}                        % physical (pre-squashing) system
\newcommand{\Tcl}[1]{\mathcal{T}(#1)}        % trace-class operators on H(#1)
\newcommand{\Ins}{\mathcal{I}}               % generalized-instrument branch map
\newcommand{\redNmsm}{\mathcal{N}}              % POVM measurement channel
\newcommand{\Eprot}{\mathcal{E}}             % n-round protocol channel
\newcommand{\Grnd}{\mathsf{G}}               % one round map of a protocol
\newcommand{\Kfin}{\mathsf{K}}               % end-of-protocol processing
\newcommand{\Crit}{\mathsf{sec}}             % value of a security criterion
\newcommand{\Pflag}{\Pi_{\Flag}}             % flag-sector projector on Q
\newcommand{\Pleq}{\Pi_{\leqN}}              % preserved-block projector
\newcommand{\Pgt}{\Pi_{\gtN}}                % above-cutoff projector
\newcommand{\Sph}{S^{\mathrm{ph}}}           % family of physical inputs
\newcommand{\Srch}{S^{\mathrm{rch}}}         % reachable squashed states
\newcommand{\Ssq}{S^{\sq}}                   % squashed admissible set
\newcommand{\Ffss}{F^{\mathrm{FSS}}}         % flag-state-squasher target
\newcommand{\Fwp}{F^{\mathrm{WP}}}           % multi-bound WPFSS target
\newcommand{\Mwp}{M^{\mathrm{WP}}}           % renormalized above-cutoff POVM
\newcommand{\Mdec}{{M'}^{\mathrm{ph}}}       % declared physical family on B
\newcommand{\Ppvm}{P^{\Flag}}                % flag PVM
\newcommand{\Lmp}{{\Lambda_{\mathrm{mp}}}}   % measure-and-prepare channel
\newcommand{\Lwp}{{\Lambda^{\mathrm{WP}}}}   % multi-bound WPFSS channel
\newcommand{\Lfss}{{\Lambda^{\mathrm{FSS}}}} % flag-state-squasher channel
\newcommand{\lmax}{\lambda_{\max}}           % maximum eigenvalue
\newcommand{\Lnoise}{\widetilde{\Lambda}}    % flag-noise channel
\newcommand{\cg}{\varpi}                     % declared coarse-graining map
\newcommand{\MCK}{\mathsf{MC}}               % multiclick raw-pattern set
\newcommand{\redChain}{\mathfrak{C}}            % reduction chain
\newcommand{\Lnk}{\ell}                      % a single typed link
\newcommand{\dsq}{\preceq}                      % the squashing relation
\newcommand{\opn}[1]{\left\lVert #1\right\rVert}           % operator (spectral) norm
\newcommand{\mxn}[1]{\left\lVert #1\right\rVert_{\max}}    % entrywise max modulus
\newcommand{\Fld}{\mathbb{F}}                   % exact ordered substrate field
\newcommand{\Fldi}{\mathbb{F}_{i}}              % its Gaussian extension F + iF
\newcommand{\lat}[1]{\big\lfloor #1\big\rfloor_{\kappa}} % round DOWN onto the lattice

% --- protocol-compiler ---------------------------------------------------------
\newcommand{\Spec}{\mathfrak{S}}             % protocol specification (compiler input)
\newcommand{\pcChain}{\mathfrak{R}}            % reduction chain (two typed segments)
\newcommand{\Chsrc}{\pcChain_{\mathrm{src}}}   % its Alice-side source-map segment
\newcommand{\Chrec}{\pcChain_{\mathrm{rec}}}   % its receiver-side squashing segment
\newcommand{\rvirt}{\rho^{\mathrm{v}}}       % virtual (PRE-source-map) signal state
\newcommand{\Fph}{F^{\mathrm{ph}}}           % physical (pre-reduction) POVM element
\newcommand{\Vph}{\mathcal{V}^{\mathrm{ph}}} % physical outcome alphabet
\newcommand{\Tpub}{\mathcal{T}_{\mathrm{pub}}}% public-choice alphabet
\newcommand{\Uset}{\mathcal{U}}              % residual outcome alphabet
\newcommand{\idxm}{\mathsf{idx}}             % label-to-index map
\newcommand{\Asig}{A'}                       % transmitted (signal) system -- the REAL one
\newcommand{\Avirt}{A''}                     % virtual transmitted system (source descriptor of Xi_1)
\newcommand{\Abar}{\bar{A}}                  % Alice's measured register
\newcommand{\Ahat}{\hat{A}}                  % Alice's shield register
\newcommand{\pcNmsm}{\mathcal{N}_{\mathsf A}}  % Alice's measurement channel (see the RULES block)
\newcommand{\Phip}{\Phi^{+}}                 % maximally entangled vector label
\newcommand{\Smap}{\mathfrak{s}}             % the emitted sector map (fraktur: a compile-time object)

% --- receiver ------------------------------------------------------------------
\newcommand{\Modes}{\mathbb{M}}              % declared ordered mode set
\newcommand{\Inp}{\mathbb{I}}                % declared input modes
\newcommand{\InpC}{\mathbb{I}_{\Chan}}       % channel input modes
\newcommand{\InpS}{\mathbb{I}_{\Prb}}        % probe (source) input modes
\newcommand{\Env}{\mathbb{E}}                % dilation (environment) modes
\newcommand{\Port}[1]{\mathbb{P}_{#1}}       % mode set of detector port #1
\newcommand{\Emd}[1]{E_{#1}}                 % coordinate projector on modes #1
\newcommand{\Gmat}{G}                        % instrument (transfer) matrix
\newcommand{\Fock}{\mathfrak{F}}             % Fock / second-quantization functor
\newcommand{\Symn}[1]{\mathrm{Sym}^{#1}}     % symmetric tensor power
\newcommand{\normord}[1]{{:}\mspace{2mu}#1\mspace{2mu}{:}}   % normal ordering
\newcommand{\rk}{\operatorname{rank}}        % rank (physics.sty owns \rank)
\newcommand{\SCK}{\mathsf{SC}}               % single-click pattern set
\newcommand{\Recv}{\mathfrak{R}}             % receiver instrument tuple
\newcommand{\Cat}{\mathfrak{C}}              % component catalogue
\newcommand{\Grph}{\mathfrak{E}}             % declared element graph
\newcommand{\defic}{\mathbb{L}}              % loss deficiency operator
\newcommand{\Pvac}[1]{\Pi_{\vac}(#1)}        % joint no-click element on ports #1
\newcommand{\pmc}{p_{\mc}}                   % uniform high-photon mc weight
\newcommand{\Bnd}{\beta}                     % the scalar bound function beta(n)
\newcommand{\Rot}{R}                         % planar rotation matrix
\newcommand{\Wplate}{W}                      % retarder (wave-plate) matrix
\newcommand{\Bsp}{B}                         % beamsplitter core matrix

% --- statistics ----------------------------------------------------------------
\newcommand{\honest}{p^{\mathrm{hon}}}          % honest per-round distribution
\newcommand{\ecomp}{\varepsilon_{\mathrm{comp}}}% completeness budget
\newcommand{\simplex}[1]{\Delta(#1)}            % probability simplex
\newcommand{\Chatcon}{\Chatset_{\mathrm{con}}}  % constrained outcomes
\newcommand{\Ireg}[3]{\mathcal{I}_{#1}(#2,#3)}  % regularized incomplete beta
\newcommand{\Ncount}[1]{\mathsf{N}_{#1}}        % per-outcome count
\newcommand{\qlo}{\mathbf{q}^{-}}               % box lower bounds
\newcommand{\qhi}{\mathbf{q}^{+}}               % box upper bounds
 for it, so the
%  file is read from spec/modules/<id>/.
%
%  RULES
%    * A macro that spec/notation.md catalogues is a verbatim copy of it
%      (part A). Changing one is a notation change: bump spec/notation.md.
%    * Part B holds the module-local macros. A macro defined by two modules
%      with the same expansion appears once, under the first module (in the
%      order below) that defines it. Two macros with the same name and
%      DIFFERENT expansions are both kept under a module prefix, and the
%      claims write the prefixed names:
%        \Chain  -> \redChain (reductions, \mathfrak{C}: the reduction chain)
%                   \pcChain  (protocol-compiler, \mathfrak{R}: the two-segment
%                             chain; \Chsrc and \Chrec are built on it)
%        \Nmsm   -> \redNmsm  (reductions, \mathcal{N}: a POVM measurement
%                             channel)
%                   \pcNmsm   (\mathcal{N}_{\mathsf A}: Alice's measurement
%                             channel, written by
%                             transmitter/thm:certificate-sufficiency,
%                             transmitter/rem:xreg and
%                             reductions/lem:source-transfer)
%    * Deliberately omitted: a \provenance{tag}{details} macro. Provenance
%      is a header field of every claim file (spec/claim-format.md section
%      3); a claim body that calls such a macro fails to compile, which is
%      the intended signal.
%
%  REQUIRED PACKAGES, loaded by the assembler's preamble before this file:
%    amsmath, amssymb, amsfonts, mathtools, dsfont, braket, physics, xcolor
%    (braket before physics, so that physics wins for \ket, \bra, \tr, ...).
%
%  This file contains only definitions: no \usepackage, no document body.
% =============================================================================


% =============================================================================
%  PART A -- notation macros, the ones spec/notation.md catalogues, each a
%  verbatim copy of its catalogue row; grouped under numbered headings.
% =============================================================================

% -----------------------------------------------------------------------------
% 1. State spaces and operator sets (a "system label" -> spaces of operators)
% -----------------------------------------------------------------------------
\newcommand{\Hs}[1]{\mathcal{H}(#1)}        % Hilbert space of a system label
\newcommand{\Lin}[1]{\mathcal{L}(#1)}       % bounded linear operators on H(#1)
\newcommand{\Linn}[2]{\mathcal{L}(#1,#2)}   % linear maps H(#1) -> H(#2)
\newcommand{\Pos}[1]{\mathrm{Pos}(#1)}      % positive semidefinite operators on H(#1)
\newcommand{\Density}{\mathrm{D}}           % density operators: \Density(A) = D(A)
\newcommand{\Seq}[1]{S_{=}(#1)}             % normalized states (trace = 1)
\newcommand{\Sleq}[1]{S_{\leq}(#1)}         % subnormalized states (0 <= trace <= 1)
\newcommand{\id}{\mathds{1}}                % identity operator
\newcommand{\idon}[1]{\mathds{1}_{#1}}      % identity on a named system, e.g. \idon{C}

% DEPRECATED: applies \mathcal to both arguments itself, so it takes bare
% letters, not system labels. Prefer \Linn; do not use in new contracts.
\providecommand{\Lop}[2]{\mathcal{L}(\mathcal{#1}, \mathcal{#2})}


% -----------------------------------------------------------------------------
% 2. Composite systems
% -----------------------------------------------------------------------------
\newcommand{\dsum}{\oplus}                  % direct sum of system labels
% Tensor product: use \otimes directly.
% n-fold direct sum of an operator space: \Lin{A}^{\oplus n}


% -----------------------------------------------------------------------------
% 3. Events, conditioning and classical registers (for STATES)
%    Convention:  \wedge  -> subnormalized "and-event" component, e.g. \rho_{A\wedge x}
%                 |        -> normalized conditional state,        e.g. \rho_{A|x}
% (Write these directly in subscripts; no extra macros needed.)
% -----------------------------------------------------------------------------


% -----------------------------------------------------------------------------
% 4. Photon-number subspaces / direct-sum blocks
%    Convention:  operators carry the restriction "on top", in brackets;
%                 states carry it as a \wedge subscript (see Sec. 3).
% -----------------------------------------------------------------------------
\newcommand{\Pnum}[1]{P_{#1}}               % photon-number projector (system from context)
\newcommand{\Pnumon}[2]{P_{#1}^{#2}}        % with explicit system superscript, e.g. P_n^C
\newcommand{\blk}[2]{#1^{(#2)}}             % operator block, e.g. \blk{M}{\leqN} = M^{(<=N)}
\newcommand{\leqN}{\leq N}                  % use as \blk{M}{\leqN}
\newcommand{\gtN}{> N}                      % use as \blk{M}{\gtN}
\newcommand{\eqn}{= n}                      % use as \blk{M}{\eqn}


% -----------------------------------------------------------------------------
% 5. System labels used in the protocol
% -----------------------------------------------------------------------------
\newcommand{\Chan}{C}                       % quantum-channel (Eve) input to Bob
\newcommand{\Prb}{P}                        % probe input
\newcommand{\Flag}{{\mathrm{Flg}}}          % flag system (EXISTS ONLY AFTER squashing)
\newcommand{\Loss}{{\mathrm{Lss}}}          % loss system
\newcommand{\Vc}{{\mathrm{Vcm}}}            % vacuum system
\newcommand{\Nout}{N_{\mathrm{out}}}
\newcommand{\Nin}{N_{\mathrm{in}}}


% -----------------------------------------------------------------------------
% 6. Alphabets
% -----------------------------------------------------------------------------
\newcommand{\Yset}{\mathcal{Y}}             % Bob's (coarse-grained) detection alphabet
\newcommand{\Yraw}{\mathcal{Y}_{\mathrm{raw}}}% raw click patterns {0,1}^4
\newcommand{\Yscset}{\mathcal{Y}_{\singlec}}% single-click outcomes
\newcommand{\Ymcset}{\mathcal{Y}_{\mc}}     % multi-click outcomes
\newcommand{\Tset}{\mathcal{T}_{\mathrm{set}}}% private setting set (the probe settings in the optical instance)
\newcommand{\Xset}{\mathcal{X}}             % Alice's alphabet
\newcommand{\Dset}{\{H,V,+,-\}}             % polarization / detector labels
\newcommand{\pass}{\top}                    % pass sifting


% -----------------------------------------------------------------------------
% 7. Mode / outcome abbreviations
% -----------------------------------------------------------------------------
\newcommand{\on}{{\mathrm{on}}}
\newcommand{\off}{{\mathrm{off}}}
\newcommand{\dc}{{\mathrm{dc}}}             % dark count
\newcommand{\sq}{\mathrm{sq}}               % squashed
\newcommand{\mc}{\mathrm{mc}}               % multi-click
\newcommand{\nc}{\mathrm{nc}}               % no-click
\newcommand{\singlec}{\mathrm{sc}}          % single-click
\newcommand{\vac}{\mathrm{vac}}             % vacuum
\newcommand{\kout}{{k_\mathrm{out}}}
\newcommand{\kin}{{k_\mathrm{in}}}


% -----------------------------------------------------------------------------
% 8. Maps, channels and miscellaneous operators
% -----------------------------------------------------------------------------
\newcommand{\Prob}{\mathrm{Prob}}           % probability (use \Prob[\cdot] or \Prob(\cdot))
\newcommand{\lmin}{\lambda_{\min}}          % minimum eigenvalue
\newcommand{\Wflag}{W_{\Flag}}              % flag-space weight
\newcommand{\transpose}{\mathsf{T}}
\newcommand{\innerp}[2]{\left\langle #1 \vert #2 \right\rangle}
\renewcommand{\vec}[1]{\boldsymbol{#1}}     % NB: overrides the LaTeX arrow accent


% -----------------------------------------------------------------------------
% 9. Key-rate numerics
% -----------------------------------------------------------------------------
\newcommand{\cG}{\mathcal{G}}               % protocol map
\newcommand{\cZ}{\mathcal{Z}}               % pinching map
\newcommand{\Amap}{\mathcal{A}_{\mathrm{ann}}} % announcement + measurement map
\newcommand{\Amapnb}{\mathcal{A}_{\mathrm{ann}|\mathrm{nb}}}
\newcommand{\simpl}{}                       % empty suffix for simplified block maps
\newcommand{\acSet}{\mathcal{S}_{\mathrm{acc}}} % acceptance set
\newcommand{\testE}{\mathcal{E}_{\mathrm{test}}} % test-statistics map
\newcommand{\ecLeak}{\mathrm{leak}^{\mathrm{EC}}_{\mathrm{obs}}}
\newcommand{\pgen}{p_{\mathrm{gen}}}
\newcommand{\askeyr}{R_{\infty}}             % asymptotic key rate
\newcommand{\ppass}{p_{\mathrm{pass}}}       % sifting pass probability


% -----------------------------------------------------------------------------
% 10. Finite-size analysis / Renyi entropies
% -----------------------------------------------------------------------------

% --- Round types: test vs. key generation (register T, alphabet \Fset) --------
\newcommand{\tst}{\mathrm{test}}            % test round
\newcommand{\gen}{\mathrm{gen}}             % key-generation round
\newcommand{\Fset}{\mathcal{T}}             % round-type alphabet {test, gen} (not \Tset)
\newcommand{\Chat}{\widehat{C}}             % acceptance-test register
\newcommand{\Chatset}{\widehat{\mathcal{C}}}% its alphabet

% --- Basis / key extraction functions ----------------------------------------
\newcommand{\bA}{b_{A}}                     % Alice's basis-extraction function
\newcommand{\bB}{b_{B}}                     % Bob's basis-extraction function
\newcommand{\kB}{k_{B}}                     % Bob's key-extraction function

% --- Maps --------------------------------------------------------------------
\newcommand{\Mprot}{\mathcal{M}}            % single-round protocol map
\newcommand{\Phimap}{\boldsymbol{\Phi}}     % test-statistics map
\newcommand{\Pdeph}{\mathcal{P}}            % dephasing (pinching) map on \Blreg

% --- Isometries --------------------------------------------------------------
\newcommand{\Vmsm}{V_{\mathrm{msm}}}        % measurement isometry
\newcommand{\Vann}{V_{\mathrm{ann}}}        % announcement isometry
\newcommand{\VI}{V_{I}}                     % Eve's copy of the announcements
\newcommand{\VChat}{V_{\Chat}}              % acceptance-data isometry
\newcommand{\VR}{V_{R}}                     % sifting + key map -> raw key R
\newcommand{\VZ}{V_{Z}}                     % raw key -> key register Z
\newcommand{\Wprot}{W}                      % full protocol isometry

% --- Renyi entropies ---------------------------------------------------------
% SECURITY-CRITICAL: see spec/notation.md §4. Never change an arrow to
% make an inequality work out; never convert silently between the sandwiched
% (tilde) and Petz (overline) families.
\newcommand{\Hsandup}{\widetilde{H}^{\uparrow}_{\alpha}}   % sandwiched, up
\newcommand{\Hsanddown}{\widetilde{H}^{\downarrow}_{\alpha}}% sandwiched, down
\newcommand{\Hpetzup}{\overline{H}^{\uparrow}_{\alpha}}    % Petz, up
\newcommand{\Hpetzdown}{\overline{H}^{\downarrow}_{\alpha}}% Petz, down
\newcommand{\hup}{h^{\uparrow}_{\alpha}}                   % single-round rate
\newcommand{\hqkd}{h^{\uparrow,\mathrm{QKD}}_{\alpha}}     % single-round QKD rate
\newcommand{\hlow}{\underline{h}_{\alpha}}                 % lower bound on \hqkd
% --- spec/notation.md 3.5 and 4.5 -------------------------------------------
% \hlb  is the first-order linearization lower bound on \hlow: the headline
%       key rate (the `first-order` tier of foundations/def:tiers). NOT \hlow
%       itself.
% \tnum is the instance's one declared numerical tolerance. Its
%       near-collision with the declared marginal \tau_A is disclosed in
%       spec/notation.md 3.5; bare \tau is reserved and denotes neither.
\newcommand{\hlb}{\underline{h}^{\mathrm{FO}}_{\alpha}}    % first-order lower bound on \hlow
\newcommand{\tnum}{\tau_{\mathrm{num}}}                    % declared numerical tolerance


% =============================================================================
%  PART B -- module-local macros, grouped by module. Deprecated aliases,
%  defined so that a text written with them compiles, used in no new contract: \Lop,
%  \testE, \Pdeph, \pass, \ecLeak (spec/notation.md section 3.9).
% =============================================================================

% --- meat-core -----------------------------------------------------------------
\newcommand{\SecSet}{\mathbb{S}}            % sector index set of Q
\newcommand{\Qsys}{Q}                       % declared sector-carrying measured system
\newcommand{\Vset}{\mathcal{V}}             % declared joint outcome alphabet of the Q-POVM
\newcommand{\CAset}{\mathcal{C}_{A}}        % Alice announcement alphabet
\newcommand{\CBset}{\mathcal{C}_{B}}        % Bob announcement alphabet
\newcommand{\Party}{\mathfrak{P}}           % declared party partition
\newcommand{\Oset}[1]{\mathcal{O}_{#1}}     % readout alphabet of party #1
\newcommand{\rd}[1]{\pi_{#1}}               % readout map of party #1
\newcommand{\Sadm}{\mathcal{S}_{A}}         % admissible Alice subspace
\newcommand{\nomarg}{\mathsf{none}}         % "no trusted marginal declared"
\newcommand{\Zset}{\mathcal{Z}_{\perp}}     % raw-key alphabet incl. \perp
\newcommand{\Zval}{\mathcal{Z}}             % key alphabet without \perp
\newcommand{\Etil}{\widetilde{E}}           % single-round purifying register
\newcommand{\Ehat}{\widehat{E}}             % Eve's n-round register
\newcommand{\OmAT}{\Omega_{\mathrm{AT}}}    % acceptance event
\newcommand{\epa}{\varepsilon_{\mathrm{PA}}}
\newcommand{\eev}{\varepsilon_{\mathrm{EV}}}
\newcommand{\lEC}{\lambda_{\mathrm{EC}}}
\newcommand{\Decl}{\mathfrak{D}}            % protocol declaration
\newcommand{\CList}{\mathfrak{L}}           % declared constraint list
\newcommand{\PinS}{\mathcal{P}_{\SecSet}}   % sector pinching on Q
\newcommand{\Feas}{\mathcal{W}}             % feasible set of block states
\newcommand{\Chkey}{\Chatset_{\mathrm{key}}}% key-live acceptance values
\newcommand{\Kann}{\mathcal{K}}             % announcement-class alphabet
\newcommand{\hold}{\mathsf{h}}              % declared key holder (A or B)
\newcommand{\sfD}{D}                        % classical KL divergence (base 2)
\newcommand{\Dtwo}{D}                       % quantum relative entropy (base 2)
\newcommand{\Dsand}{\widetilde{D}_{\alpha}} % sandwiched Renyi divergence
\newcommand{\idmap}{\mathrm{id}}            % identity channel
\newcommand{\Aev}{\mathcal{A}}              % Eve's per-round attack channel
\newcommand{\AttSet}{\mathfrak{A}}          % licensed attack-channel set
\newcommand{\Chacc}{\Chatset_{\mathrm{acc}}}% acceptance support of \acSet
\newcommand{\epsstar}{\varepsilon^{\star}}      % distance to an exact POVM
\newcommand{\dfH}{\Delta_{\mathrm{herm}}}       % Hermiticity defect
\newcommand{\dfP}{\Delta_{\mathrm{pos}}}        % negative part
\newcommand{\dfC}{\Delta_{\mathrm{cpl}}}        % completeness defect
\newcommand{\lmoq}{g}                       % acceptance-side gradient block
\newcommand{\lmoJ}{G}                       % feasible-side gradient block
\newcommand{\bacc}{\beta_{\mathrm{acc}}}    % certified acceptance-side minimum
\newcommand{\Itset}{\mathfrak{I}}           % finite index set of iterates (C22)

% --- reductions ----------------------------------------------------------------
\newcommand{\Bsys}{B}                        % physical (pre-squashing) system
\newcommand{\Tcl}[1]{\mathcal{T}(#1)}        % trace-class operators on H(#1)
\newcommand{\Ins}{\mathcal{I}}               % generalized-instrument branch map
\newcommand{\redNmsm}{\mathcal{N}}              % POVM measurement channel
\newcommand{\Eprot}{\mathcal{E}}             % n-round protocol channel
\newcommand{\Grnd}{\mathsf{G}}               % one round map of a protocol
\newcommand{\Kfin}{\mathsf{K}}               % end-of-protocol processing
\newcommand{\Crit}{\mathsf{sec}}             % value of a security criterion
\newcommand{\Pflag}{\Pi_{\Flag}}             % flag-sector projector on Q
\newcommand{\Pleq}{\Pi_{\leqN}}              % preserved-block projector
\newcommand{\Pgt}{\Pi_{\gtN}}                % above-cutoff projector
\newcommand{\Sph}{S^{\mathrm{ph}}}           % family of physical inputs
\newcommand{\Srch}{S^{\mathrm{rch}}}         % reachable squashed states
\newcommand{\Ssq}{S^{\sq}}                   % squashed admissible set
\newcommand{\Ffss}{F^{\mathrm{FSS}}}         % flag-state-squasher target
\newcommand{\Fwp}{F^{\mathrm{WP}}}           % multi-bound WPFSS target
\newcommand{\Mwp}{M^{\mathrm{WP}}}           % renormalized above-cutoff POVM
\newcommand{\Mdec}{{M'}^{\mathrm{ph}}}       % declared physical family on B
\newcommand{\Ppvm}{P^{\Flag}}                % flag PVM
\newcommand{\Lmp}{{\Lambda_{\mathrm{mp}}}}   % measure-and-prepare channel
\newcommand{\Lwp}{{\Lambda^{\mathrm{WP}}}}   % multi-bound WPFSS channel
\newcommand{\Lfss}{{\Lambda^{\mathrm{FSS}}}} % flag-state-squasher channel
\newcommand{\lmax}{\lambda_{\max}}           % maximum eigenvalue
\newcommand{\Lnoise}{\widetilde{\Lambda}}    % flag-noise channel
\newcommand{\cg}{\varpi}                     % declared coarse-graining map
\newcommand{\MCK}{\mathsf{MC}}               % multiclick raw-pattern set
\newcommand{\redChain}{\mathfrak{C}}            % reduction chain
\newcommand{\Lnk}{\ell}                      % a single typed link
\newcommand{\dsq}{\preceq}                      % the squashing relation
\newcommand{\opn}[1]{\left\lVert #1\right\rVert}           % operator (spectral) norm
\newcommand{\mxn}[1]{\left\lVert #1\right\rVert_{\max}}    % entrywise max modulus
\newcommand{\Fld}{\mathbb{F}}                   % exact ordered substrate field
\newcommand{\Fldi}{\mathbb{F}_{i}}              % its Gaussian extension F + iF
\newcommand{\lat}[1]{\big\lfloor #1\big\rfloor_{\kappa}} % round DOWN onto the lattice

% --- protocol-compiler ---------------------------------------------------------
\newcommand{\Spec}{\mathfrak{S}}             % protocol specification (compiler input)
\newcommand{\pcChain}{\mathfrak{R}}            % reduction chain (two typed segments)
\newcommand{\Chsrc}{\pcChain_{\mathrm{src}}}   % its Alice-side source-map segment
\newcommand{\Chrec}{\pcChain_{\mathrm{rec}}}   % its receiver-side squashing segment
\newcommand{\rvirt}{\rho^{\mathrm{v}}}       % virtual (PRE-source-map) signal state
\newcommand{\Fph}{F^{\mathrm{ph}}}           % physical (pre-reduction) POVM element
\newcommand{\Vph}{\mathcal{V}^{\mathrm{ph}}} % physical outcome alphabet
\newcommand{\Tpub}{\mathcal{T}_{\mathrm{pub}}}% public-choice alphabet
\newcommand{\Uset}{\mathcal{U}}              % residual outcome alphabet
\newcommand{\idxm}{\mathsf{idx}}             % label-to-index map
\newcommand{\Asig}{A'}                       % transmitted (signal) system -- the REAL one
\newcommand{\Avirt}{A''}                     % virtual transmitted system (source descriptor of Xi_1)
\newcommand{\Abar}{\bar{A}}                  % Alice's measured register
\newcommand{\Ahat}{\hat{A}}                  % Alice's shield register
\newcommand{\pcNmsm}{\mathcal{N}_{\mathsf A}}  % Alice's measurement channel (see the RULES block)
\newcommand{\Phip}{\Phi^{+}}                 % maximally entangled vector label
\newcommand{\Smap}{\mathfrak{s}}             % the emitted sector map (fraktur: a compile-time object)

% --- receiver ------------------------------------------------------------------
\newcommand{\Modes}{\mathbb{M}}              % declared ordered mode set
\newcommand{\Inp}{\mathbb{I}}                % declared input modes
\newcommand{\InpC}{\mathbb{I}_{\Chan}}       % channel input modes
\newcommand{\InpS}{\mathbb{I}_{\Prb}}        % probe (source) input modes
\newcommand{\Env}{\mathbb{E}}                % dilation (environment) modes
\newcommand{\Port}[1]{\mathbb{P}_{#1}}       % mode set of detector port #1
\newcommand{\Emd}[1]{E_{#1}}                 % coordinate projector on modes #1
\newcommand{\Gmat}{G}                        % instrument (transfer) matrix
\newcommand{\Fock}{\mathfrak{F}}             % Fock / second-quantization functor
\newcommand{\Symn}[1]{\mathrm{Sym}^{#1}}     % symmetric tensor power
\newcommand{\normord}[1]{{:}\mspace{2mu}#1\mspace{2mu}{:}}   % normal ordering
\newcommand{\rk}{\operatorname{rank}}        % rank (physics.sty owns \rank)
\newcommand{\SCK}{\mathsf{SC}}               % single-click pattern set
\newcommand{\Recv}{\mathfrak{R}}             % receiver instrument tuple
\newcommand{\Cat}{\mathfrak{C}}              % component catalogue
\newcommand{\Grph}{\mathfrak{E}}             % declared element graph
\newcommand{\defic}{\mathbb{L}}              % loss deficiency operator
\newcommand{\Pvac}[1]{\Pi_{\vac}(#1)}        % joint no-click element on ports #1
\newcommand{\pmc}{p_{\mc}}                   % uniform high-photon mc weight
\newcommand{\Bnd}{\beta}                     % the scalar bound function beta(n)
\newcommand{\Rot}{R}                         % planar rotation matrix
\newcommand{\Wplate}{W}                      % retarder (wave-plate) matrix
\newcommand{\Bsp}{B}                         % beamsplitter core matrix

% --- statistics ----------------------------------------------------------------
\newcommand{\honest}{p^{\mathrm{hon}}}          % honest per-round distribution
\newcommand{\ecomp}{\varepsilon_{\mathrm{comp}}}% completeness budget
\newcommand{\simplex}[1]{\Delta(#1)}            % probability simplex
\newcommand{\Chatcon}{\Chatset_{\mathrm{con}}}  % constrained outcomes
\newcommand{\Ireg}[3]{\mathcal{I}_{#1}(#2,#3)}  % regularized incomplete beta
\newcommand{\Ncount}[1]{\mathsf{N}_{#1}}        % per-outcome count
\newcommand{\qlo}{\mathbf{q}^{-}}               % box lower bounds
\newcommand{\qhi}{\mathbf{q}^{+}}               % box upper bounds

\makeatletter
\newcommand*{\ClaimHeadFont}{\bfseries}   % switched by prominence
\newcommand*{\ClaimNoteFont}{\mdseries}
\newcommand*{\ClaimShortId}{}
\newcommand*{\SetShortId}[1]{\renewcommand*{\ClaimShortId}{#1}}
\newcommand*{\RestatedName}{}
% Upright bodies throughout: italic is reserved for non-normative text.
% \@endparpenalty\@M keeps the meta line on the statement's page.
\newtheoremstyle{claim}{\topsep}{\topsep}{\normalfont\@endparpenalty\@M}{}%
  {\ClaimHeadFont}{.}{.5em}{\thmname{#1}\thmnumber{ #2}\thmnote{ {\ClaimNoteFont(#3)}}}
\newtheoremstyle{claimass}{\topsep}{\topsep}{\normalfont\@endparpenalty\@M}{}%
  {\ClaimHeadFont}{.}{.5em}{\thmname{#1}\thmnumber{ #2}%
  \ifx\ClaimShortId\@empty\else\ [\textsf{\ClaimShortId}]\fi%
  \thmnote{ {\ClaimNoteFont(#3)}}}
\newtheoremstyle{claimremark}{.5\topsep}{.5\topsep}{\small\@endparpenalty\@M}{}%
  {\itshape}{.}{.5em}{\thmname{#1}\thmnumber{ #2}\thmnote{ {\upshape(#3)}}}
\newtheoremstyle{claimreq}{\topsep}{\topsep}{\normalfont\@endparpenalty\@M}{}%
  {\color{reqAccent}\scshape}{.}{.5em}%
  {\thmname{#1}\thmnumber{ #2}\thmnote{ {\upshape(#3)}}}
\makeatother
\theoremstyle{claim}
\newtheorem{theorem}{Theorem}[section]
\newtheorem{lemma}[theorem]{Lemma}
\newtheorem{proposition}[theorem]{Proposition}
\newtheorem{corollary}[theorem]{Corollary}
\newtheorem{fact}[theorem]{Fact}
\newtheorem{definition}[theorem]{Definition}
\newtheorem*{restatedclaim}{\RestatedName}
\theoremstyle{claimass}
\newtheorem{assumption}[theorem]{Assumption}
\newtheorem*{restatedassumption}{\RestatedName}
\theoremstyle{claimremark}
\newtheorem{remark}[theorem]{Remark}
\newtheorem{example}[theorem]{Example}
\theoremstyle{claimreq}
\newtheorem{requirement}[theorem]{Requirement}
\newtheorem*{restatedrequirement}{\RestatedName}
%% --- 8. Prominence: main results a margin rule and a bold heading ----------
\newtcolorbox{mainclaim}{enhanced jigsaw, breakable, blanker,
  borderline west={1.5pt}{-8pt}{mainRule}, left=0pt, right=0pt, top=0pt, bottom=0pt,
  boxsep=0pt, before skip=\topsep, after skip=\topsep, parbox=false,
  before upper={\renewcommand*{\ClaimNoteFont}{\bfseries}}}
\newenvironment{technicalclaim}{\renewcommand*{\ClaimHeadFont}{\mdseries\color{techGrey}}%
  \renewcommand*{\ClaimNoteFont}{}}{}
%% --- 9. Non-normative text, the meta line, navigation ------------------------
\makeatletter
% Tie what follows to this paragraph: no page break, no list penalty before
% the next theorem head (the \@nobreak route a section heading uses).
\newcommand{\keepwithnext}{\par\nobreak\vspace{3pt}\@nobreaktrue
  \everypar{\if@nobreak\@nobreakfalse\clubpenalty\@M\else\clubpenalty\@clubpenalty\everypar{}\fi}}
\newcommand{\claimsummary}[2]{% #1 claim id (named destination), #2 summary
  \par\addvspace{\bigskipamount}\Needspace{6\baselineskip}%
  {\small\itshape\color{nnGrey}\noindent\hypertarget{claim-#1}{}#2\par}\keepwithnext}
\newcommand{\claimanchor}[1]{% a claim without a summary: the destination alone
  \par\addvspace{\bigskipamount}\Needspace{6\baselineskip}%
  \phantomsection\hypertarget{claim-#1}{}\keepwithnext}
\newcommand{\claimleadin}[1]{{\color{nnGrey}\noindent#1\par}\keepwithnext}
\newcommand{\claimmeta}[1]{\par\nobreak\vspace{2pt}%
  {\footnotesize\color{metaGrey}\raggedright\interlinepenalty\@M\noindent#1\par}%
  \addvspace{\medskipamount}}
% After an assumption: the sentence generated from its verification field; after
% any claim with one: its plain-text reading (user_statement). Both
% non-normative, both grey.
\newcommand{\claimverification}[1]{\par\nobreak\vspace{2pt}%
  {\footnotesize\color{metaGrey}\noindent#1\par}\nobreak}
\newcommand{\claimreading}[1]{\par\nobreak\vspace{2pt}%
  {\footnotesize\color{nnGrey}\noindent\emph{Reading:}\enspace#1\par}\nobreak}
\makeatother
\newcommand{\claimbookmark}[4]{% #1 id, #2 kind word, #3 label, #4 plain title
  \bookmark[dest=claim-#1,level=2]{#2 \getrefnumber{#3} #4}}
\newcommand{\metasep}{\unskip\enspace\textperiodcentered\enspace\ignorespaces}
\newcommand{\metalabel}[1]{\emph{#1}\enspace}
\newcommand{\sectionintro}[1]{{\color{nnGrey}\noindent#1\par}\smallskip}
\newcommand{\sectionimports}[1]{{\footnotesize\color{metaGrey}\raggedright\noindent\metalabel{Imports}#1\par}\smallskip}
\newenvironment{overview}{\par\color{nnGrey}\itshape\noindent\ignorespaces}{\par}
% References: the whole "Def. 2.1" is the link; ids are the fallback text of
% every link into another PDF (a relative URI to a named destination).
\newcommand{\mref}[2]{\hyperref[#2]{#1~\ref*{#2}}}
\newcommand{\aref}[2]{\hyperref[#2]{\textsf{#1}}}
\DeclareUrlCommand{\claimid}{\urlstyle{sf}\def\UrlBreaks{\do\/\do\:\do\-\do\.\do\_}%
  \def\UrlBigBreaks{}\def\UrlNoBreaks{}}
\newcommand{\extref}[2]{\hyperref{#1}{}{claim-#2}{\claimid{#2}}}
\newcommand{\extaref}[3]{\hyperref{#1}{}{claim-#2}{\textsf{#3}}}
\newcommand{\coderef}[2]{\href{#1}{\claimid{#2}}}
% Claim-body commands (spec/claim-format.md section 4): \xref{<full id>} prints
% the kind, title and id of a claim of another contract, linked to its stable
% destination in that contract's PDF (one \DeclareXref per cited claim below);
% \term[gloss]{phrase} marks a defined phrase for the index (the gloss feeds
% the Terms table); \tref[phrase]{text} is a use of a defined term. The
% assembler rewrites both in the body into \termdef (the definition's named
% destination term-<slug>) and \termref / \exttermref (a link to it, in this
% or another contract's PDF); the fallbacks below serve text it copies
% elsewhere, such as the digest. \iface is an interface key.
\makeatletter
\newcommand{\DeclareXref}[4]{% id, kind word, title, contract PDF
  \@namedef{xref@#1}{\hyperref{#4}{}{claim-#1}{#2 ``#3'' ({\small\claimid{#1}})}}}
\newcommand{\xref}[1]{\@ifundefined{xref@#1}%
  {\textbf{[unresolved \nolinkurl{#1}]}\@latex@warning{Unresolved xref #1}}%
  {\@nameuse{xref@#1}}}
\makeatother
\DeclareUrlCommand{\iface}{\urlstyle{tt}}
\newcommand{\term}[2][]{\emph{#2}\index{#2}}
\newcommand{\tref}[2][]{#2}
\newcommand{\termdef}[2]{\hypertarget{#1}{}\emph{#2}\index{#2}}
\newcommand{\termref}[2]{\hyperlink{#1}{#2}}
\newcommand{\exttermref}[3]{\hyperref{#1}{}{#2}{#3}}
\hypersetup{pdftitle={Acceptance statistics (statistics) -- statements digest 0.10.0}}

% The digest: the contract's labels under the prefix C-, its numbers checked.
\externaldocument[C-]{contract}
\theoremstyle{claim}\newtheorem*{digestclaim}{\RestatedName}
\theoremstyle{claimass}\newtheorem*{digestassumption}{\RestatedName}
\theoremstyle{claimreq}\newtheorem*{digestrequirement}{\RestatedName}
\makeatletter
% #1 label, #2 the number the assembler computed: a local label carrying it, so
% that a reference between two digest items stays in the digest; the build
% stops when the contract's own number differs.
\newcommand{\digestlabel}[2]{\phantomsection
  \edef\digest@theirs{\getrefnumber{C-#1}}\def\digest@ours{#2}%
  \ifx\digest@theirs\digest@ours\else
    \PackageError{digest}{#1 is #2 here but \digest@theirs\space in the contract}{}\fi
  \def\@currentlabel{#2}\label{#1}}
\makeatother
% A display's local label (amsmath records the \tag, i.e. the contract's number).
\newcommand{\digesteqlabel}[1]{\label{#1}}

% Claims of other contracts cited here (\xref, [[...]]).
\DeclareXref{foundations/def:exact-island}{Definition}{The exact island}{../foundations/contract.pdf}
\DeclareXref{foundations/req:consumer-decides-exactly}{Requirement}{Exact decision by every consumer}{../foundations/contract.pdf}
\DeclareXref{foundations/req:rounding-direction}{Requirement}{Rounding direction of emitted bounds}{../foundations/contract.pdf}
\DeclareXref{meat-core/def:admissible-acceptance-set}{Definition}{Admissible acceptance set}{../meat-core/contract.pdf}
\DeclareXref{meat-core/def:feasible}{Definition}{Feasible set and computable objective}{../meat-core/contract.pdf}
\DeclareXref{meat-core/def:postprocessing-model}{Definition}{Fixed-length post-processing model}{../meat-core/contract.pdf}
\DeclareXref{meat-core/frank-wolfe/def:first-order-setting}{Definition}{Ambient program, pairings, acceptance support and linear minimisation}{../meat-core/contract.pdf}
\DeclareXref{meat-core/frank-wolfe/rem:acc-side}{Remark}{Exact minimisation on the acceptance side}{../meat-core/contract.pdf}
\DeclareXref{meat-core/prop:minimal-divergence}{Proposition}{The minimal divergence as the value with every pair dropped}{../meat-core/contract.pdf}
\DeclareXref{meat-core/rem:hlow-range}{Remark}{Reading an infinite value}{../meat-core/contract.pdf}
\DeclareXref{meat-core/req:infinite-value}{Requirement}{Propagation of an infinite value}{../meat-core/contract.pdf}
\DeclareXref{meat-core/req:unreachable-acceptance-set}{Requirement}{The check of the acceptance set against the feasible statistics}{../meat-core/contract.pdf}
\DeclareXref{protocol-compiler/ass:coin}{Assumption}{Public choices as private coins of declared bias}{../protocol-compiler/contract.pdf}
\DeclareXref{protocol-compiler/def:declaration}{Definition}{Protocol declaration}{../protocol-compiler/contract.pdf}
\DeclareXref{protocol-compiler/def:feasible}{Definition}{Feasible set}{../protocol-compiler/contract.pdf}
\DeclareXref{protocol-compiler/def:spec}{Definition}{Protocol specification}{../protocol-compiler/contract.pdf}
\DeclareXref{reductions/ass:tsm}{Assumption}{Output-state test statistics}{../reductions/contract.pdf}

\begin{document}
\begin{center}
{\LARGE\bfseries Acceptance statistics (statistics)}\\[4pt]
{\large Statements digest, version 0.10.0}
\end{center}
\noindent{\small Every definition, fact, assumption, result and requirement of the \href{contract.pdf}{contract}, with the contract's numbers; remarks, examples, technical steps and all proofs are left out. Each heading opens its statement in the contract.}

\part*{Mathematics}\pdfbookmark[0]{Mathematics}{digest-part-1}

\section*{1\quad Frequencies and acceptance geometries}\pdfbookmark[1]{1 Frequencies and acceptance geometries}{digest-group-statistics/def:freq}

% claim statistics/def:freq
\digestlabel{def:freq}{1.1}\pdfbookmark[2]{Definition 1.1 Frequency vector and counts}{digest-statistics/def:freq}
\renewcommand*{\RestatedName}{\hyperref{contract.pdf}{}{claim-statistics/def:freq}{Definition~1.1}}%
\begin{digestclaim}[Frequency vector and counts]
For an outcome string $\hat c_1^n=(\hat c_1,\dots,\hat c_n)\in\Chatset^{\,n}$
and $\hat c\in\Chatset$ put
\begin{equation*}\tag{\getrefnumber{C-eq:freq}}\digesteqlabel{eq:freq}
  \Ncount{\hat c}(\hat c_1^n)\coloneqq
  \big|\{\,j\in\{1,\dots,n\}:\hat c_j=\hat c\,\}\big|,
  \qquad
  \operatorname{freq}_{\hat c_1^n}(\hat c)\coloneqq
  \tfrac{1}{n}\,\Ncount{\hat c}(\hat c_1^n):
\end{equation*}
the \emph{count} of the outcome $\hat c$ and the
\emph{frequency vector} of the string. Then
$\operatorname{freq}_{\hat c_1^n}\in\simplex{\Chatset}$ for every string, and
its coordinates lie in $\tfrac1n\{0,1,\dots,n\}$.
\end{digestclaim}

% claim statistics/lem:freq-in-simplex
\digestlabel{lem:freq-in-simplex}{1.2}\pdfbookmark[2]{Lemma 1.2 Frequency vectors in the probability simplex}{digest-statistics/lem:freq-in-simplex}
\renewcommand*{\RestatedName}{\hyperref{contract.pdf}{}{claim-statistics/lem:freq-in-simplex}{Lemma~1.2}}%
\begin{digestclaim}[Frequency vectors in the probability simplex]
Let $n\ge1$, and write $\simplex{\Chatset}$ for the set of
$\mathbf{q}\in\mathbb{R}^{\Chatset}$ with non-negative coordinates summing
to one.
\begin{enumerate}[label=(\roman*),leftmargin=2.2em]
\item For every outcome string $\hat c_1^n\in\Chatset^{\,n}$, every count
$\Ncount{\hat c}(\hat c_1^n)$ of Definition~\ref{def:freq} is a non-negative
integer, at most $n$ of them are non-zero, and
$\sum_{\hat c\in\Chatset}\Ncount{\hat c}(\hat c_1^n)=n$. Hence
$\operatorname{freq}_{\hat c_1^n}\in\simplex{\Chatset}$.
\item For every set $S\subseteq\mathbb{R}^{\Chatset}$ and every random string
$\Chat_1^n$ with values in $\Chatset^{\,n}$, the sets
$\{\operatorname{freq}_{\Chat_1^n}\in S\cap\simplex{\Chatset}\}$ and
$\{\operatorname{freq}_{\Chat_1^n}\in S\}$ of sample points are equal. In
particular, whenever either is an event,
\begin{equation*}\tag{\getrefnumber{C-eq:simplex-free}}\digesteqlabel{eq:simplex-free}
  \Pr\big[\operatorname{freq}_{\Chat_1^n}\in S\cap\simplex{\Chatset}\big]
  =\Pr\big[\operatorname{freq}_{\Chat_1^n}\in S\big].
\end{equation*}
\end{enumerate}
\end{digestclaim}

% claim statistics/def:matched-alphabet
\digestlabel{def:matched-alphabet}{1.3}\pdfbookmark[2]{Definition 1.3 Alphabet and round count matched to the consuming declaration}{digest-statistics/def:matched-alphabet}
\renewcommand*{\RestatedName}{\hyperref{contract.pdf}{}{claim-statistics/def:matched-alphabet}{Definition~1.3}}%
\begin{digestclaim}[Alphabet and round count matched to the consuming declaration]
An acceptance alphabet $\Chatset$ and a round count $n$ are a
\emph{matched alphabet and
round count} when $\Chatset$ is the finite, non-empty acceptance alphabet of
the same declaration $\Decl$ (\xref{protocol-compiler/def:declaration}) the
consumer solves, the codomain of $\phi$, with a fixed enumeration, and
$n$ is the same total round count.
\end{digestclaim}

% claim statistics/def:taxonomy
\digestlabel{def:taxonomy}{1.4}\pdfbookmark[2]{Definition 1.4 Admissible geometries and the canonical box form}{digest-statistics/def:taxonomy}
\renewcommand*{\RestatedName}{\hyperref{contract.pdf}{}{claim-statistics/def:taxonomy}{Definition~1.4}}%
\begin{digestclaim}[Admissible geometries and the canonical box form]
An \emph{admissible geometry} is a pair consisting of a \emph{tag}
and its \emph{data}, drawn from exactly the following two cases, each with
its data, the set $\Gamma$ the data denote and its
\emph{validation predicate}. There is no third tag.
\begin{description}[leftmargin=1.9em]
\item[$\mathrm{singleton}$.] Data: a vector
$\mathbf{q}^{\ast}\in\mathbb{Q}^{\Chatset}$. Denotation
$\Gamma\coloneqq\{\mathbf{q}^{\ast}\}$. \emph{Validation predicate:}
$q^{\ast}(\hat c)\ge0$ for every $\hat c$ and $\sum_{\hat c}q^{\ast}(\hat
c)=1$.
\item[$\mathrm{polytope}$.] Data: an integer row count $m_{\mathrm{row}}\ge0$,
a matrix $B\in\mathbb{Q}^{\,m_{\mathrm{row}}\times\Chatset}$ and a vector
$\mathbf{b}\in\mathbb{Q}^{\,m_{\mathrm{row}}}$. Denotation
\begin{equation*}\tag{\getrefnumber{C-eq:polytope}}\digesteqlabel{eq:polytope}
  \Gamma\coloneqq\{\mathbf{q}\in\simplex{\Chatset}:\ B\mathbf{q}\le\mathbf{b}\},
\end{equation*}
the inequality componentwise. \emph{Validation predicate:} feasibility of the
linear system $\{\mathbf{q}\ge0,\ \sum_{\hat c}q(\hat c)=1,\
B\mathbf{q}\le\mathbf{b}\}$ over the rationals.
\end{description}
A $\mathrm{polytope}$ geometry is in \emph{canonical box form} when it carries, in addition, a
pair $\qlo,\qhi\in\mathbb{Q}^{\Chatset}$ with $0\le\qlo\le\qhi$
componentwise, and its $(B,\mathbf{b})$ are exactly
\begin{equation*}\tag{\getrefnumber{C-eq:boxrows}}\digesteqlabel{eq:boxrows}
  m_{\mathrm{row}}=2m,\qquad
  \begin{aligned}
    &\text{row }i:\ \ \delta_{\hat c^{(i)}}^{\transpose},
      &&\text{right-hand side } q^{+}(\hat c^{(i)}),\\
    &\text{row }m+i:\ \ -\delta_{\hat c^{(i)}}^{\transpose},
      &&\text{right-hand side } -q^{-}(\hat c^{(i)}),
  \end{aligned}
  \qquad i=1,\dots,m,
\end{equation*}
where $m\coloneqq|\Chatset|$, $\hat c^{(1)},\dots,\hat c^{(m)}$ is the
declaration's enumeration of $\Chatset$ and $\delta_{\hat c}$ is the standard
basis vector of $\mathbb{R}^{\Chatset}$ at $\hat c$. Row indices are integers,
never alphabet elements; equivalently
\begin{equation*}\tag{\getrefnumber{C-eq:boxset}}\digesteqlabel{eq:boxset}
  \Gamma=\{\mathbf{q}\in\simplex{\Chatset}:\
  q^{-}(\hat c)\le q(\hat c)\le q^{+}(\hat c)\ \ \forall\hat c\in\Chatset\}.
\end{equation*}
\end{digestclaim}

% claim statistics/thm:geometry-validation
\digestlabel{thm:geometry-validation}{1.5}\pdfbookmark[2]{Theorem 1.5 Validity of the admissible geometries}{digest-statistics/thm:geometry-validation}
\begin{mainclaim}
\renewcommand*{\RestatedName}{\hyperref{contract.pdf}{}{claim-statistics/thm:geometry-validation}{Theorem~1.5}}%
\begin{digestclaim}[Validity of the admissible geometries]
Let $\Chatset$ and $n$ be a matched alphabet and round count
(Definition~\ref{def:matched-alphabet}), and let $(\text{tag},\text{data})$
be an admissible geometry with denotation $\Gamma$.
\begin{enumerate}[label=(\roman*),leftmargin=2.2em]
\item If the validation predicate holds, $\Gamma$ is a non-empty, closed,
convex and compact subset of $\simplex{\Chatset}$. In particular
\xref{meat-core/def:admissible-acceptance-set}, that $\acSet$ is a closed convex subset of the
probability simplex on $\Chatset$, holds by construction for every admissible
geometry whose validation predicate holds.
\item For the tag $\mathrm{singleton}$, the validation predicate is necessary
and sufficient for $\Gamma$ to be a non-empty subset of $\simplex{\Chatset}$.
\item For the tag $\mathrm{polytope}$ in canonical box form with
$0\le\qlo\le\qhi$, the validation predicate, the condition
$\Gamma\neq\varnothing$ and the closed-form condition
\begin{equation*}\tag{\getrefnumber{C-eq:boxfeas}}\digesteqlabel{eq:boxfeas}
  \sum_{\hat c\in\Chatset}q^{-}(\hat c)\ \le\ 1\ \le\
  \sum_{\hat c\in\Chatset}q^{+}(\hat c)
\end{equation*}
are equivalent.
\end{enumerate}
\end{digestclaim}
\end{mainclaim}

\section*{2\quad Binomial tails}\pdfbookmark[1]{2 Binomial tails}{digest-group-statistics/thm:binomial-tails}

% claim statistics/thm:binomial-tails
\digestlabel{thm:binomial-tails}{2.1}\pdfbookmark[2]{Theorem 2.1 Binomial tails in closed form}{digest-statistics/thm:binomial-tails}
\renewcommand*{\RestatedName}{\hyperref{contract.pdf}{}{claim-statistics/thm:binomial-tails}{Theorem~2.1}}%
\begin{digestclaim}[Binomial tails in closed form]
Let $n\ge1$, $p\in[0,1]$, and let $\Ncount{}$ be a binomial random variable
with parameters $(n,p)$, i.e.\
$\Pr[\Ncount{}=i]=\binom{n}{i}p^{i}(1-p)^{n-i}$ for $0\le i\le n$. Write
\begin{equation*}\tag{\getrefnumber{C-eq:ireg}}\digesteqlabel{eq:ireg}
  \Ireg{x}{a}{b}\coloneqq
  \frac{\int_0^{x}t^{a-1}(1-t)^{b-1}\,\mathrm{d}t}
       {\int_0^{1}t^{a-1}(1-t)^{b-1}\,\mathrm{d}t},
  \qquad x\in[0,1],\ a,b\ge1 .
\end{equation*}
Then
\begin{align*}
  \Pr[\Ncount{}\ge k]&=\frac{n!}{(k-1)!\,(n-k)!}
    \int_0^{p}t^{k-1}(1-t)^{n-k}\,\mathrm{d}t
    =\Ireg{p}{k}{n-k+1}, &&1\le k\le n,
    \tag{\getrefnumber{C-eq:uppertail}}\digesteqlabel{eq:uppertail}\\
  \Pr[\Ncount{}\le k]&=1-\Pr[\Ncount{}\ge k+1]
    =\Ireg{1-p}{n-k}{k+1}, &&0\le k\le n-1,
    \tag{\getrefnumber{C-eq:lowertail}}\digesteqlabel{eq:lowertail}
\end{align*}
together with the boundary values $\Pr[\Ncount{}\ge k]=1$ for $k\le0$,
$\Pr[\Ncount{}\ge k]=0$ for $k>n$, $\Pr[\Ncount{}\le k]=0$ for $k<0$ and
$\Pr[\Ncount{}\le k]=1$ for $k\ge n$. Moreover $p\mapsto\Pr[\Ncount{}\ge k]$
is non-decreasing on $[0,1]$ for every $k\ge1$.
\end{digestclaim}

% claim statistics/cor:tail-forms
\digestlabel{cor:tail-forms}{2.2}\pdfbookmark[2]{Corollary 2.2 The two strict binomial tails}{digest-statistics/cor:tail-forms}
\renewcommand*{\RestatedName}{\hyperref{contract.pdf}{}{claim-statistics/cor:tail-forms}{Corollary~2.2}}%
\begin{digestclaim}[The two strict binomial tails]
Let $n\ge1$, $p\in[0,1]$ and let $\Ncount{}$ be binomial with parameters
$(n,p)$, as in Theorem~\ref{thm:binomial-tails}. Then, for $0\le k\le n$,
\begin{equation*}\tag{\getrefnumber{C-eq:tails-used}}\digesteqlabel{eq:tails-used}
\begin{aligned}
  \Pr[\Ncount{}<k]&=\Pr[\Ncount{}\le k-1]
  =\begin{cases}0,&k=0,\\[2pt] \Ireg{1-p}{n-k+1}{k},&1\le k\le n,\end{cases}\\[4pt]
  \Pr[\Ncount{}>k]&=\Pr[\Ncount{}\ge k+1]
  =\begin{cases}\Ireg{p}{k+1}{n-k},&0\le k\le n-1,\\[2pt] 0,&k=n.\end{cases}
\end{aligned}
\end{equation*}
\end{digestclaim}

% claim statistics/lem:exact-strict-tails
\digestlabel{lem:exact-strict-tails}{2.3}\pdfbookmark[2]{Lemma 2.3 Exact strict tails at a rational success probability}{digest-statistics/lem:exact-strict-tails}
\renewcommand*{\RestatedName}{\hyperref{contract.pdf}{}{claim-statistics/lem:exact-strict-tails}{Lemma~2.3}}%
\begin{digestclaim}[Exact strict tails at a rational success probability]
Let $n\ge1$, let $p\in[0,1]\cap\mathbb{Q}$ and let $\Ncount{}$ be binomial
with parameters $(n,p)$, as in Theorem~\ref{thm:binomial-tails}. For
$0\le k\le n$,
\begin{equation*}\tag{\getrefnumber{C-eq:tails-sums}}\digesteqlabel{eq:tails-sums}
  \Pr[\Ncount{}<k]=\sum_{i=0}^{k-1}\binom{n}{i}p^{i}(1-p)^{n-i},
  \qquad
  \Pr[\Ncount{}>k]=\sum_{i=k+1}^{n}\binom{n}{i}p^{i}(1-p)^{n-i},
\end{equation*}
an empty sum being $0$, and:
\begin{enumerate}[label=(\roman*),leftmargin=2.2em]
\item both strict tails are rational numbers, obtained from $n$, $k$ and $p$
by finitely many additions and multiplications of rationals, and they are
equal to the closed forms \eqref{eq:tails-used} of
Corollary~\ref{cor:tail-forms};
\item for every rational $\varepsilon$, whether
$\Pr[\Ncount{}<k]\le\varepsilon$ holds, and whether
$\Pr[\Ncount{}>k]\le\varepsilon$ holds, is decided exactly by evaluating the
rational number \eqref{eq:tails-sums} and comparing two rationals, which is a
comparison of two integers.
\end{enumerate}
\end{digestclaim}

\section*{3\quad The exact box oracle}\pdfbookmark[1]{3 The exact box oracle}{digest-group-statistics/thm:box-oracle}

% claim statistics/thm:box-oracle
\digestlabel{thm:box-oracle}{3.1}\pdfbookmark[2]{Theorem 3.1 Acceptance support and the exact box oracle}{digest-statistics/thm:box-oracle}
\begin{mainclaim}
\renewcommand*{\RestatedName}{\hyperref{contract.pdf}{}{claim-statistics/thm:box-oracle}{Theorem~3.1}}%
\begin{digestclaim}[Acceptance support and the exact box oracle]
Let $\Chatset$ and $n$ be a matched alphabet and round count
(Definition~\ref{def:matched-alphabet}), let $\acSet$ be a non-empty
canonical box, \eqref{eq:boxset}, with
$0\le\qlo\le\qhi$, and let $g\in\mathbb{Q}^{\Chatset}$.
\begin{enumerate}[label=(\roman*),leftmargin=2.2em]
\item \emph{(acceptance support)} The set
$\Chatset^{+}\coloneqq\{\hat c\in\Chatset: q^{+}(\hat c)>0\}$ contains
$\operatorname{supp}\mathbf{q}$ for every $\mathbf{q}\in\acSet$; hence it
contains the acceptance support
$\Chacc=\bigcup_{\mathbf{q}\in\acSet}\operatorname{supp}\mathbf{q}$ of
\xref{meat-core/frank-wolfe/def:first-order-setting}, which it need not equal: the box
constraints together with the simplex constraint can force $q(\hat c)=0$ at an
outcome with $q^{+}(\hat c)>0$. For a $\mathrm{singleton}$ geometry, which
declares no box, the union of the supports is
$\operatorname{supp}\mathbf{q}^{\ast}$, the acceptance support itself.
\item \emph{(exact minimum in closed form)} Enumerate
$\Chatset=\{c_1,\dots,c_m\}$ so that $g(c_1)\le\dots\le g(c_m)$. Put
$r\coloneqq1-\sum_{\hat c}q^{-}(\hat c)$ and
$D_i\coloneqq q^{+}(c_i)-q^{-}(c_i)\ge0$, let $j^{\ast}$ be the least index
with $\sum_{i\le j^{\ast}}D_i\ge r$, and define
\begin{equation*}\tag{\getrefnumber{C-eq:greedy}}\digesteqlabel{eq:greedy}
  q^{\mathrm{opt}}(c_i)\coloneqq
  \begin{cases}
    q^{+}(c_i), & i<j^{\ast},\\[2pt]
    q^{-}(c_i)+\big(r-\sum_{i'<j^{\ast}}D_{i'}\big), & i=j^{\ast},\\[2pt]
    q^{-}(c_i), & i>j^{\ast}.
  \end{cases}
\end{equation*}
Then $j^{\ast}$ exists, $\mathbf{q}^{\mathrm{opt}}\in\acSet$, and
\begin{equation*}\tag{\getrefnumber{C-eq:exactmin}}\digesteqlabel{eq:exactmin}
  \langle g,\mathbf{q}^{\mathrm{opt}}\rangle
  \;=\;\min_{\mathbf{q}\in\acSet}\ \langle g,\mathbf{q}\rangle ,
\end{equation*}
the minimum being attained. All quantities are exact rationals when
$g,\qlo,\qhi$ are.
\item \emph{(dual witness)} With $B,\mathbf{b}$ as in
\eqref{eq:boxrows}, $\lambda\coloneqq g(c_{j^{\ast}})$, $v\coloneqq-\lambda$ and
$\mathbf{u}\coloneqq(\mathbf{u}^{+},\mathbf{u}^{-})$ given by
$u^{+}(\hat c)\coloneqq\max\{\lambda-g(\hat c),0\}$ and
$u^{-}(\hat c)\coloneqq\max\{g(\hat c)-\lambda,0\}$, one has
$\mathbf{u}\ge0$,
\begin{equation*}\tag{\getrefnumber{C-eq:dualfeas}}\digesteqlabel{eq:dualfeas}
  g(\hat c)+(B^{\transpose}\mathbf{u})(\hat c)+v=0
  \qquad\text{for every }\hat c\in\Chatset,
\end{equation*}
in particular the left-hand side is $\ge0$, and
\begin{equation*}\tag{\getrefnumber{C-eq:dualvalue}}\digesteqlabel{eq:dualvalue}
  -\langle\mathbf{b},\mathbf{u}\rangle-v
  \;=\;\langle g,\mathbf{q}^{\mathrm{opt}}\rangle
  \;=\;\min_{\mathbf{q}\in\acSet}\langle g,\mathbf{q}\rangle .
\end{equation*}
The lower bound $-\langle\mathbf{b},\mathbf{u}\rangle-v$ on the minimum that
this witness gives is therefore \emph{tight} for the canonical box, and every
downward-rounded evaluation of $\langle g,\mathbf{q}^{\mathrm{opt}}\rangle$ is
a lower bound on $\min_{\mathbf{q}\in\acSet}\langle g,\mathbf{q}\rangle$.
\end{enumerate}
\end{digestclaim}
\end{mainclaim}

\section*{4\quad The honest envelope}\pdfbookmark[1]{4 The honest envelope}{digest-group-statistics/def:envelope}

% claim statistics/def:envelope
\digestlabel{def:envelope}{4.1}\pdfbookmark[2]{Definition 4.1 Honest envelope and its envelope set}{digest-statistics/def:envelope}
\renewcommand*{\RestatedName}{\hyperref{contract.pdf}{}{claim-statistics/def:envelope}{Definition~4.1}}%
\begin{digestclaim}[Honest envelope and its envelope set]
An \emph{honest envelope} on $\Chatset$ is a pair
${\honest}^{-},{\honest}^{+}\in\mathbb{Q}^{\Chatset}$ of exact rational
vectors with
\begin{equation*}\tag{\getrefnumber{C-eq:envelope}}\digesteqlabel{eq:envelope}
  0\ \le\ {\honest}^{-}(\hat c)\ \le\ {\honest}^{+}(\hat c)\ \le\ 1
  \quad\forall\hat c\in\Chatset,
  \qquad
  \sum_{\hat c\in\Chatset}{\honest}^{-}(\hat c)\ \le\ 1\ \le\
  \sum_{\hat c\in\Chatset}{\honest}^{+}(\hat c).
\end{equation*}
Its \emph{envelope set} is
\begin{equation*}\tag{\getrefnumber{C-eq:envelope-set}}\digesteqlabel{eq:envelope-set}
  \mathsf{H}\ \coloneqq\ \big\{\mathbf{p}\in\simplex{\Chatset}:\
  {\honest}^{-}(\hat c)\le p(\hat c)\le{\honest}^{+}(\hat c)\ \
  \forall\hat c\in\Chatset\big\},
\end{equation*}
the set of honest laws the envelope admits. An honest envelope with
${\honest}^{-}={\honest}^{+}$ is a \emph{degenerate envelope}; then
\eqref{eq:envelope} forces the common value to be a law
$\honest\in\simplex{\Chatset}\cap\mathbb{Q}^{\Chatset}$ and
$\mathsf{H}=\{\honest\}$. A single honest law $\honest$ (the
\emph{point form}) is read as the degenerate envelope
$(\honest,\honest)$. The pair $({\honest}^{-},{\honest}^{+})$ is a
canonical box form in the sense of Definition~\ref{def:taxonomy}
(rational, $0\le{\honest}^{-}\le{\honest}^{+}$), $\mathsf{H}$ is its
denotation \eqref{eq:boxset}, and the second condition of
\eqref{eq:envelope} is its closed-form feasibility condition
\eqref{eq:boxfeas}.
\end{digestclaim}

% claim statistics/lem:hon-degenerate
\digestlabel{lem:hon-degenerate}{4.2}\pdfbookmark[2]{Lemma 4.2 The envelope assumption at a degenerate envelope}{digest-statistics/lem:hon-degenerate}
\renewcommand*{\RestatedName}{\hyperref{contract.pdf}{}{claim-statistics/lem:hon-degenerate}{Lemma~4.2}}%
\begin{digestclaim}[The envelope assumption at a degenerate envelope]
Let $\honest\in\simplex{\Chatset}\cap\mathbb{Q}^{\Chatset}$ and let
$\Chat_1,\dots,\Chat_n$ be $\Chatset$-valued random variables on one
probability space. The following are equivalent:
\begin{enumerate}[label=(\alph*),leftmargin=2.2em]
\item the statement of Assumption~\ref{ass:hon-env} holds for the degenerate
envelope $(\honest,\honest)$ of Definition~\ref{def:envelope}: there is a
filtration $\mathcal{F}_0\subseteq\dots\subseteq\mathcal{F}_n$ with $\Chat_j$
$\mathcal{F}_j$-measurable and
$\Pr[\Chat_j=\hat c\,|\,\mathcal{F}_{j-1}]=\honest(\hat c)$ almost surely,
for every $j$ and $\hat c$;
\item the statement of Assumption~\ref{ass:hon} holds: $\Chat_1,\dots,\Chat_n$
are independent and identically distributed with law $\honest$.
\end{enumerate}
\end{digestclaim}

% claim statistics/def:conservative-evaluation
\digestlabel{def:conservative-evaluation}{4.3}\pdfbookmark[2]{Definition 4.3 Exact data and conservative evaluation}{digest-statistics/def:conservative-evaluation}
\renewcommand*{\RestatedName}{\hyperref{contract.pdf}{}{claim-statistics/def:conservative-evaluation}{Definition~4.3}}%
\begin{digestclaim}[Exact data and conservative evaluation]
An evaluation of the acceptance-set construction is a
\emph{conservative evaluation} when the honest envelope's vectors
${\honest}^{-},{\honest}^{+}$ (Definition~\ref{def:envelope}; for a single
declared law, $\honest$ itself), $\ecomp$ and the allocation are exact
rationals, and every tail comparison an endpoint rule makes
(Definition~\ref{def:endpoint-rule}; for the exact rule, those of
Definition~\ref{def:interval}) is decided either exactly or against a bound
that rounds the tail \emph{upward}.
\end{digestclaim}

\section*{5\quad Assumptions}\pdfbookmark[1]{5 Assumptions}{digest-group-statistics/ass:hon}

% claim statistics/ass:hon
\digestlabel{ass:hon}{5.1}\pdfbookmark[2]{Assumption 5.1 Honest i.i.d. rounds}{digest-statistics/ass:hon}
\SetShortId{A-HON}
\renewcommand*{\RestatedName}{\hyperref{contract.pdf}{}{claim-statistics/ass:hon}{Assumption~5.1}}%
\begin{digestassumption}[Honest i.i.d. rounds]
(i) On an honest run the acceptance statistics $\Chat_1,\dots,\Chat_n$ are
independent and identically distributed with law $\honest$ on $\Chatset$.
\end{digestassumption}

% claim statistics/ass:hon-env
\digestlabel{ass:hon-env}{5.2}\pdfbookmark[2]{Assumption 5.2 Honest rounds inside a declared envelope}{digest-statistics/ass:hon-env}
\SetShortId{A-HON-ENV}
\renewcommand*{\RestatedName}{\hyperref{contract.pdf}{}{claim-statistics/ass:hon-env}{Assumption~5.2}}%
\begin{digestassumption}[Honest rounds inside a declared envelope]
Let $({\honest}^{-},{\honest}^{+})$ be the declared honest envelope
(Definition~\ref{def:envelope}). On an honest run there is a filtration
$\mathcal{F}_0\subseteq\mathcal{F}_1\subseteq\dots\subseteq\mathcal{F}_n$ on
the underlying probability space such that
(i) $\Chat_j$ is $\mathcal{F}_j$-measurable for every $j\in\{1,\dots,n\}$;
(ii) for every such $j$ and every $\hat c\in\Chatset$, almost surely,
\begin{equation*}\tag{\getrefnumber{C-eq:hon-env}}\digesteqlabel{eq:hon-env}
  {\honest}^{-}(\hat c)\ \le\
  \Pr\big[\Chat_j=\hat c\ \big|\ \mathcal{F}_{j-1}\big]\ \le\
  {\honest}^{+}(\hat c).
\end{equation*}
\end{digestassumption}

% claim statistics/ass:alphabet
\digestlabel{ass:alphabet}{5.3}\pdfbookmark[2]{Assumption 5.3 An acceptance value recorded in every round}{digest-statistics/ass:alphabet}
\SetShortId{A-ALPHABET}
\renewcommand*{\RestatedName}{\hyperref{contract.pdf}{}{claim-statistics/ass:alphabet}{Assumption~5.3}}%
\begin{digestassumption}[An acceptance value recorded in every round]
(i) In every one of the $n$ rounds of the run, an acceptance value in
$\Chatset$ is recorded.
\end{digestassumption}

% claim statistics/ass:reach
\digestlabel{ass:reach}{5.4}\pdfbookmark[2]{Assumption 5.4 An honest law produced by a feasible state}{digest-statistics/ass:reach}
\SetShortId{A-REACH}
\renewcommand*{\RestatedName}{\hyperref{contract.pdf}{}{claim-statistics/ass:reach}{Assumption~5.4}}%
\begin{digestassumption}[An honest law produced by a feasible state]
(i) There exists $\omega^{\mathrm{hon}}\in\Feas$, the feasible set of the
consuming declaration (\xref{protocol-compiler/def:feasible}), with
$\Phimap[\omega^{\mathrm{hon}}]\in\mathsf{H}$, the envelope set of the
declared honest envelope (Definition~\ref{def:envelope}); for a single
declared law $\honest$ (the degenerate envelope, $\mathsf{H}=\{\honest\}$)
this reads $\Phimap[\omega^{\mathrm{hon}}]=\honest$.
\end{digestassumption}

\part*{Engine requirements}\pdfbookmark[0]{Engine requirements}{digest-part-2}

\section*{6\quad Requirements: the emitted geometry}\pdfbookmark[1]{6 Requirements: the emitted geometry}{digest-group-statistics/req:geometry-tag}

% claim statistics/req:geometry-tag
\digestlabel{req:geometry-tag}{6.1}\pdfbookmark[2]{Requirement 6.1 Acceptance sets under exactly the two geometry tags}{digest-statistics/req:geometry-tag}
\renewcommand*{\RestatedName}{\hyperref{contract.pdf}{}{claim-statistics/req:geometry-tag}{Requirement~6.1}}%
\begin{digestrequirement}[Acceptance sets under exactly the two geometry tags]
The engine emits every acceptance set under one of the two tags of
Definition~\ref{def:taxonomy} (\iface{geometry_tag}) together with the data of
that tag, and a consumer presented with an acceptance set under any
other tag refuses it rather than guessing.
\end{digestrequirement}

% claim statistics/req:box-bounds-key
\digestlabel{req:box-bounds-key}{6.2}\pdfbookmark[2]{Requirement 6.2 Recorded bounds of a canonical box}{digest-statistics/req:box-bounds-key}
\renewcommand*{\RestatedName}{\hyperref{contract.pdf}{}{claim-statistics/req:box-bounds-key}{Requirement~6.2}}%
\begin{digestrequirement}[Recorded bounds of a canonical box]
The engine emits a geometry in canonical box form with its pair
$(\qlo,\qhi)$ recorded under \iface{box_bounds}, beside the rows
\eqref{eq:boxrows} that the pair determines.
\end{digestrequirement}

% claim statistics/req:acceptance-support
\digestlabel{req:acceptance-support}{6.3}\pdfbookmark[2]{Requirement 6.3 The emitted acceptance support}{digest-statistics/req:acceptance-support}
\renewcommand*{\RestatedName}{\hyperref{contract.pdf}{}{claim-statistics/req:acceptance-support}{Requirement~6.3}}%
\begin{digestrequirement}[The emitted acceptance support]
The engine emits under \iface{acceptance_support} the set
$\Chatset^{+}=\{\hat c\in\Chatset:q^{+}(\hat c)>0\}$ of
Theorem~\ref{thm:box-oracle}\,(i) for a geometry in canonical box form, and
the set $\operatorname{supp}\mathbf{q}^{\ast}$ for a $\mathrm{singleton}$
geometry. For a box the emitted set is a superset of the acceptance support
$\Chacc$ of \xref{meat-core/frank-wolfe/def:first-order-setting} and is never to be
substituted for it; for a singleton the emitted set and $\Chacc$ coincide.
\end{digestrequirement}

% claim statistics/req:matched-alphabet
\digestlabel{req:matched-alphabet}{6.4}\pdfbookmark[2]{Requirement 6.4 Alphabet and round count taken from the consuming declaration}{digest-statistics/req:matched-alphabet}
\renewcommand*{\RestatedName}{\hyperref{contract.pdf}{}{claim-statistics/req:matched-alphabet}{Requirement~6.4}}%
\begin{digestrequirement}[Alphabet and round count taken from the consuming declaration]
The engine docks an acceptance set into the declaration a consumer
solves only on a matched alphabet and round count
(Definition~\ref{def:matched-alphabet}).
\begin{enumerate}[label=(\roman*),leftmargin=2.2em]
\item A set it builds, it builds on the acceptance alphabet, in its declared
  enumeration, of the declaration the consumer solves and on the round count
  of the one run parameter the consumer reads, and before the set is docked
  it compares the set's alphabet and enumeration with the declaration's. A
  disagreement is a hard error of the consumer's input validation: it is a
  refusal code of neither module, nothing is repaired, and no result is
  produced.
\item A set that the declaration itself carries, with no alphabet of its
  own, it reads positionally: its $k$-th coordinate is the $k$-th acceptance
  value in the declaration's declared enumeration, and its round count is
  the one run parameter the consumer reads, as in (i). A set whose number of coordinates differs
  from the size of the acceptance alphabet is refused as ill formed, and no
  result is produced.
\end{enumerate}
\end{digestrequirement}

\part*{Binomial box construction}\pdfbookmark[0]{Binomial box construction}{digest-part-3}

\section*{7\quad Admissible inputs, endpoint rules and the guard}\pdfbookmark[1]{7 Admissible inputs, endpoint rules and the guard}{digest-group-statistics/binomial-box/def:endpoint-rule}

% claim statistics/binomial-box/def:endpoint-rule
\digestlabel{def:endpoint-rule}{7.1}\pdfbookmark[2]{Definition 7.1 Admissible inputs, endpoint rules and guarded endpoints}{digest-statistics/binomial-box/def:endpoint-rule}
\renewcommand*{\RestatedName}{\hyperref{contract.pdf}{}{claim-statistics/binomial-box/def:endpoint-rule}{Definition~7.1}}%
\begin{digestclaim}[Admissible inputs, endpoint rules and guarded endpoints]
Fix $n\ge1$ and let $\Ncount{}$ be binomial with parameters $(n,p)$, where
$p\in[0,1]$ is rational. An \emph{admissible endpoint rule} $R$ assigns to every
$(n,p,\varepsilon^{-},\varepsilon^{+})$ with rational
$\varepsilon^{-},\varepsilon^{+}\in[0,1)$ two integers
$R^{-}=R^{-}(n,p,\varepsilon^{-})$ and $R^{+}=R^{+}(n,p,\varepsilon^{+})$
such that
\begin{enumerate}[label=(\alph*),leftmargin=2.2em]
\item \emph{(admissibility)} $R^{-},R^{+}\in\{0,\dots,n\}$ and
\begin{equation*}\tag{\getrefnumber{C-eq:rule-admissible}}\digesteqlabel{eq:rule-admissible}
  \Pr\big[\Ncount{}<R^{-}\big]\le\varepsilon^{-},
  \qquad
  \Pr\big[\Ncount{}>R^{+}\big]\le\varepsilon^{+};
\end{equation*}
\item \emph{(zero budget)} if $0<p<1$, then $R^{-}=0$ at $\varepsilon^{-}=0$
and $R^{+}=n$ at $\varepsilon^{+}=0$;
\item \emph{(degenerate mass)} $R^{-}=R^{+}=0$ at $p=0$ and $R^{-}=R^{+}=n$
at $p=1$, for every $\varepsilon^{\pm}$.
\end{enumerate}
The rule is \emph{monotone in the allocation} when, the other arguments fixed,
$\varepsilon^{-}\mapsto R^{-}$ is non-decreasing and
$\varepsilon^{+}\mapsto R^{+}$ is non-increasing.

An \emph{admissible input} consists of
\begin{enumerate}[label=(\roman*),leftmargin=2.2em]
\item a finite, non-empty acceptance alphabet $\Chatset$ with a fixed
enumeration;
\item a round count, an integer $n\ge1$;
\item an honest input on $\Chatset$: exactly one of a law
$\honest\in\simplex{\Chatset}$, the point form, and an honest envelope
$({\honest}^{-},{\honest}^{+})$ (Definition~\ref{def:envelope}); a law in the
point form is read as the degenerate envelope $(\honest,\honest)$;
\item a completeness budget $\ecomp\in[0,1)$;
\item a set $\Chatcon\subseteq\Chatset$ of constrained outcomes;
\item an allocation
$\varepsilon^{-},\varepsilon^{+}:\Chatcon\to\mathbb{Q}_{\ge0}$ with
$\sum_{\hat c\in\Chatcon}(\varepsilon^{-}(\hat c)+\varepsilon^{+}(\hat
c))\le\ecomp$;
\item exactness: every coordinate of the honest input, $\ecomp$ and every
value of the allocation is an exact rational;
\item a choice of rule, the exact binomial quantile rule of
Definition~\ref{def:interval} or the Bernstein rule of
Proposition~\ref{prop:rule-bernstein}, and a choice of one of two
evaluations, by the closed forms or by the enclosures of
Proposition~\ref{prop:bracketed-endpoints}, which on the exact rule decides
its tail comparisons: by the enclosures, it gives the rule
$R_{\mathrm{encl}}$ of that proposition in place of the exact rule;
\item consistency with the consumer: $\Chatset$, its enumeration and
$n$ are those of the declaration the consumer solves, a
matched alphabet and round count (Definition~\ref{def:matched-alphabet}).
\end{enumerate}
The clauses fix the form of the data --- an enumerated set $\Chatset$, an
integer $n$, one honest input that is a vector (the point form) or a pair of
vectors (the envelope form), a number $\ecomp$, a set $\Chatcon$ with a pair
of vectors indexed by it, and a named rule and evaluation --- and conditions
on data of that form.

Fix an admissible input, with its honest envelope
$({\honest}^{-},{\honest}^{+})$. For $\hat c\in\Chatcon$ write
$p^{-}_{\hat c}\coloneqq{\honest}^{-}(\hat c)$ and
$p^{+}_{\hat c}\coloneqq{\honest}^{+}(\hat c)$, let $\Ncount{\hat c}^{-}$ and
$\Ncount{\hat c}^{+}$ be binomial with parameters $(n,p^{-}_{\hat c})$ and
$(n,p^{+}_{\hat c})$ respectively, and abbreviate
$R^{-}\coloneqq R^{-}(n,p^{-}_{\hat c},\varepsilon^{-}(\hat c))$ and
$R^{+}\coloneqq R^{+}(n,p^{+}_{\hat c},\varepsilon^{+}(\hat c))$: the lower
endpoint is the rule's at the envelope's lower value, the upper endpoint the
rule's at its upper value. The \emph{guarded endpoints} of $R$ at $\hat c$, after the
\emph{honest-containment guard}, are
\begin{equation*}\tag{\getrefnumber{C-eq:guard}}\digesteqlabel{eq:guard}
  k^{-}(\hat c)\coloneqq\min\big\{R^{-},\ \lfloor n\,p^{-}_{\hat c}\rfloor\big\},
  \qquad
  k^{+}(\hat c)\coloneqq\max\big\{R^{+},\ \lceil n\,p^{+}_{\hat c}\rceil\big\}.
\end{equation*}
With the degenerate envelope, $p^{-}_{\hat c}=p^{+}_{\hat c}=\honest(\hat c)
\eqqcolon p_{\hat c}$, $\Ncount{\hat c}^{-}$ and $\Ncount{\hat c}^{+}$ have
the one law of $\Ncount{\hat c}$, binomial $(n,p_{\hat c})$, and
\eqref{eq:guard} reads $k^{-}(\hat c)=\min\{R^{-},\lfloor np_{\hat
c}\rfloor\}$, $k^{+}(\hat c)=\max\{R^{+},\lceil np_{\hat c}\rceil\}$ with both
rule endpoints at $p_{\hat c}$.

The \emph{side-wise reading} is the following. By
\eqref{eq:guard}, $k^{-}(\hat c)$ depends on the envelope only through
$p^{-}_{\hat c}$ and $k^{+}(\hat c)$ only through $p^{+}_{\hat c}$. A claim of
this submodule that fixes one value $p_{\hat c}$ and one count
$\Ncount{\hat c}$ binomial $(n,p_{\hat c})$ and speaks of the guarded
endpoints --- Definition~\ref{def:interval} and
Propositions~\ref{prop:rule-exact}, \ref{prop:rule-bernstein} and
\ref{prop:bracketed-endpoints} do --- states each assertion about the lower
endpoint (one naming $k^{-}$, $R^{-}$, $\varepsilon^{-}$ and the lower tail)
with $p_{\hat c}\coloneqq p^{-}_{\hat c}$ and
$\Ncount{\hat c}\coloneqq\Ncount{\hat c}^{-}$, and each assertion about the
upper endpoint with $p_{\hat c}\coloneqq p^{+}_{\hat c}$ and
$\Ncount{\hat c}\coloneqq\Ncount{\hat c}^{+}$; with the degenerate envelope
the two readings coincide and are the claim as written.
\end{digestclaim}

% claim statistics/binomial-box/lem:guard
\digestlabel{lem:guard}{7.2}\pdfbookmark[2]{Lemma 7.2 The honest-containment guard for every rule}{digest-statistics/binomial-box/lem:guard}
\renewcommand*{\RestatedName}{\hyperref{contract.pdf}{}{claim-statistics/binomial-box/lem:guard}{Lemma~7.2}}%
\begin{digestclaim}[The honest-containment guard for every rule]
Let $R$ be an admissible endpoint rule (Definition~\ref{def:endpoint-rule})
and $\hat c\in\Chatcon$; write $p^{\mp}\coloneqq p^{\mp}_{\hat c}$ for the
envelope's values at $\hat c$,
$R^{-}\coloneqq R^{-}(n,p^{-},\varepsilon^{-}(\hat c))$,
$R^{+}\coloneqq R^{+}(n,p^{+},\varepsilon^{+}(\hat c))$, and let
$k^{\mp}(\hat c)$ be the guarded endpoints \eqref{eq:guard}. Then:
\begin{enumerate}[label=(\roman*),leftmargin=2.2em]
\item \emph{(ordering)}
\begin{equation*}\tag{\getrefnumber{C-eq:ordering}}\digesteqlabel{eq:ordering}
  0\ \le\ k^{-}(\hat c)\ \le\ \lfloor n p^{-}\rfloor
  \ \le\ \lceil n p^{+}\rceil\ \le\ k^{+}(\hat c)\ \le\ n ;
\end{equation*}
\item \emph{(the budgets are met)} the maps $k\mapsto\Pr[\Ncount{\hat c}^{-}<k]$
and $k\mapsto\Pr[\Ncount{\hat c}^{+}>k]$ are non-decreasing and non-increasing
respectively on the integers; consequently
\begin{equation*}\tag{\getrefnumber{C-eq:budget-met}}\digesteqlabel{eq:budget-met}
  \Pr\big[\Ncount{\hat c}^{-}<k^{-}(\hat c)\big]\le\varepsilon^{-}(\hat c),
  \qquad
  \Pr\big[\Ncount{\hat c}^{+}>k^{+}(\hat c)\big]\le\varepsilon^{+}(\hat c),
\end{equation*}
and more generally \emph{every} pair $(\tilde k^{-},\tilde k^{+})$ with
$0\le\tilde k^{-}\le R^{-}$ and $R^{+}\le\tilde k^{+}\le n$ satisfies
\eqref{eq:budget-met} in place of $k^{\mp}(\hat c)$: enlarging the rule's
interval is always admissible;
\item \emph{(exact cost of the guard)} relative to the rule's own interval
$[R^{-},R^{+}]$, the guard \eqref{eq:guard} enlarges the interval by exactly
\begin{equation*}\tag{\getrefnumber{C-eq:guard-cost}}\digesteqlabel{eq:guard-cost}
  R^{-}-k^{-}(\hat c)
    =\max\big\{0,\ R^{-}-\lfloor np^{-}\rfloor\big\},
  \qquad
  k^{+}(\hat c)-R^{+}
    =\max\big\{0,\ \lceil np^{+}\rceil-R^{+}\big\}
\end{equation*}
counts on the lower and the upper side respectively;
\item \emph{(distance from the honest value)} measured from the envelope's
value on each side, the guard costs at most $1/n$ \emph{unconditionally}:
\begin{equation*}\tag{\getrefnumber{C-eq:guard-from-p}}\digesteqlabel{eq:guard-from-p}
  0\le p^{-}-\frac{k^{-}(\hat c)}{n}
    \le\max\Big\{p^{-}-\frac{R^{-}}{n},\ \frac1n\Big\},
  \qquad
  0\le \frac{k^{+}(\hat c)}{n}-p^{+}
    \le\max\Big\{\frac{R^{+}}{n}-p^{+},\ \frac1n\Big\},
\end{equation*}
so on whichever side the guard is active, the emitted endpoint lies within
$1/n$ of that side's envelope value.
\end{enumerate}
With the degenerate envelope, $p^{-}=p^{+}=p_{\hat c}$, both
$\Ncount{\hat c}^{\mp}$ are binomial $(n,p_{\hat c})$, and (i)--(iv) are
statements about that single law.
\end{digestclaim}

\section*{8\quad The exact rule}\pdfbookmark[1]{8 The exact rule}{digest-group-statistics/binomial-box/def:interval}

% claim statistics/binomial-box/def:interval
\digestlabel{def:interval}{8.1}\pdfbookmark[2]{Definition 8.1 The exact binomial quantile rule}{digest-statistics/binomial-box/def:interval}
\renewcommand*{\RestatedName}{\hyperref{contract.pdf}{}{claim-statistics/binomial-box/def:interval}{Definition~8.1}}%
\begin{digestclaim}[The exact binomial quantile rule]
Fix an admissible input (Definition~\ref{def:endpoint-rule}), with its
round count $n\ge1$, its honest envelope $({\honest}^{-},{\honest}^{+})$
(Definition~\ref{def:envelope}; a single law $\honest$ is the
degenerate envelope), its constrained outcomes
$\Chatcon\subseteq\Chatset$, its completeness budget $\ecomp\in[0,1)$ and its
budget allocation $\varepsilon^{-},\varepsilon^{+}:\Chatcon\to\mathbb{Q}_{\ge0}$
with $\sum_{\hat c\in\Chatcon}(\varepsilon^{-}(\hat c)+\varepsilon^{+}(\hat
c))\le\ecomp$; in particular $0\le\varepsilon^{\pm}(\hat c)\le\ecomp<1$ for
every $\hat c\in\Chatcon$. For $\hat c\in\Chatcon$ let $p_{\hat c}$ be the
rule's argument --- by the side-wise reading of
Definition~\ref{def:endpoint-rule}, $p_{\hat c}=p^{-}_{\hat c}$ in
\eqref{eq:kappaminus} and $p_{\hat c}=p^{+}_{\hat c}$ in \eqref{eq:kappaplus},
both equal to $\honest(\hat c)$ for the degenerate envelope --- and let
$\Ncount{\hat c}$ be binomial with parameters $(n,p_{\hat c})$. Define the
\emph{extremal endpoints}
\begin{align*}
  \kappa^{-}(\hat c)&\coloneqq\max\big\{k\in\{0,\dots,n\}:\
    \Pr[\Ncount{\hat c}<k]\le\varepsilon^{-}(\hat c)\big\},
    \tag{\getrefnumber{C-eq:kappaminus}}\digesteqlabel{eq:kappaminus}\\
  \kappa^{+}(\hat c)&\coloneqq\min\big\{k\in\{0,\dots,n\}:\
    \Pr[\Ncount{\hat c}>k]\le\varepsilon^{+}(\hat c)\big\}.
    \tag{\getrefnumber{C-eq:kappaplus}}\digesteqlabel{eq:kappaplus}
\end{align*}
The \emph{exact binomial quantile
rule} is the endpoint rule
$R^{\mp}_{\mathrm{ex}}(n,p_{\hat c},\varepsilon^{\mp}(\hat c))
\coloneqq\kappa^{\mp}(\hat c)$. Its guarded endpoints are those of
\eqref{eq:guard} with $R^{\mp}=\kappa^{\mp}(\hat c)$, $\kappa^{-}(\hat c)$
taken at $p^{-}_{\hat c}$ and $\kappa^{+}(\hat c)$ at $p^{+}_{\hat c}$, i.e.\
\begin{equation*}
  k^{-}(\hat c)=\min\big\{\kappa^{-}(\hat c),\ \lfloor n\,p^{-}_{\hat c}\rfloor\big\},
  \qquad
  k^{+}(\hat c)=\max\big\{\kappa^{+}(\hat c),\ \lceil n\,p^{+}_{\hat c}\rceil\big\}.
\end{equation*}
\end{digestclaim}

% claim statistics/binomial-box/prop:rule-exact
\digestlabel{prop:rule-exact}{8.2}\pdfbookmark[2]{Proposition 8.2 Properties of the exact binomial quantile rule}{digest-statistics/binomial-box/prop:rule-exact}
\renewcommand*{\RestatedName}{\hyperref{contract.pdf}{}{claim-statistics/binomial-box/prop:rule-exact}{Proposition~8.2}}%
\begin{digestclaim}[Properties of the exact binomial quantile rule]
Assume a conservative evaluation
(Definition~\ref{def:conservative-evaluation}). With the notation of
Definition~\ref{def:interval}, for every $\hat c\in\Chatcon$:
\begin{enumerate}[label=(\roman*),leftmargin=2.2em]
\item the two sets in \eqref{eq:kappaminus}--\eqref{eq:kappaplus} are non-empty
subsets of the finite set $\{0,\dots,n\}$, so $\kappa^{\pm}(\hat c)$ exist and
are unique;
\item \emph{(the exact rule is an admissible endpoint rule)} the pair
$R^{\mp}_{\mathrm{ex}}=\kappa^{\mp}(\hat c)$ satisfies
Definition~\ref{def:endpoint-rule}\,(a)--(c): $\kappa^{\pm}(\hat c)\in
\{0,\dots,n\}$ and
\begin{equation*}\tag{\getrefnumber{C-eq:exact-admissible}}\digesteqlabel{eq:exact-admissible}
  \Pr\big[\Ncount{\hat c}<\kappa^{-}(\hat c)\big]\le\varepsilon^{-}(\hat c),
  \qquad
  \Pr\big[\Ncount{\hat c}>\kappa^{+}(\hat c)\big]\le\varepsilon^{+}(\hat c);
\end{equation*}
if $0<p_{\hat c}<1$ then $\kappa^{-}(\hat c)=0$ at $\varepsilon^{-}(\hat c)=0$
and $\kappa^{+}(\hat c)=n$ at $\varepsilon^{+}(\hat c)=0$; and
$\kappa^{-}(\hat c)=\kappa^{+}(\hat c)=0$ at $p_{\hat c}=0$,
$\kappa^{-}(\hat c)=\kappa^{+}(\hat c)=n$ at $p_{\hat c}=1$. Moreover the
rule is monotone in the allocation: $\kappa^{-}(\hat c)$ is non-decreasing in
$\varepsilon^{-}(\hat c)$ and $\kappa^{+}(\hat c)$ is non-increasing in
$\varepsilon^{+}(\hat c)$. Consequently Lemma~\ref{lem:guard} applies to the
exact rule, and its guarded endpoints $k^{\mp}(\hat c)$ satisfy
\eqref{eq:ordering} and \eqref{eq:budget-met};
\item \emph{(extremality and admissible relaxation)} $\kappa^{-}(\hat c)$ is
the largest and $\kappa^{+}(\hat c)$ the smallest integer in $\{0,\dots,n\}$
satisfying the respective inequality of \eqref{eq:exact-admissible}, and
\emph{every} pair $(\tilde k^{-},\tilde k^{+})$ with
$0\le\tilde k^{-}\le\kappa^{-}(\hat c)$ and
$\kappa^{+}(\hat c)\le\tilde k^{+}\le n$ satisfies \eqref{eq:budget-met} in
place of $k^{\mp}(\hat c)$. Enlarging the interval is therefore always
admissible; shrinking it beyond $[\kappa^{-},\kappa^{+}]$ is not;
\item the defining inequalities are the closed forms
\eqref{eq:tails-used} of Corollary~\ref{cor:tail-forms}, so $\kappa^{\pm}(\hat c)$
are computable from $n$, $p_{\hat c}$ and the allocation alone
(Lemma~\ref{lem:exact-strict-tails});
\item \emph{(degenerate honest mass)} if $p_{\hat c}=0$ then
$k^{-}(\hat c)=k^{+}(\hat c)=0$, and if $p_{\hat c}=1$ then
$k^{-}(\hat c)=k^{+}(\hat c)=n$; in both cases $\varepsilon^{-}(\hat c)$ and
$\varepsilon^{+}(\hat c)$ are unspent, both tails in \eqref{eq:budget-met}
being exactly zero;
\item \emph{(exact cost of the guard for this rule)} write $p=p_{\hat c}$.
Relative to the raw quantile interval
$[\kappa^{-}(\hat c),\kappa^{+}(\hat c)]$, the guard \eqref{eq:guard}
enlarges the interval by exactly \eqref{eq:guard-cost} with
$R^{\mp}=\kappa^{\mp}(\hat c)$, and neither enlargement is bounded by one
count: for every $n\ge1$ and every $p\in(0,1)$, an allocation with
$\varepsilon^{-}(\hat c)\ge1-p^{\,n}$ --- admissible whenever $\ecomp$
permits, since $1-p^{\,n}<1$ --- gives $\kappa^{-}(\hat c)=n$ and a lower
enlargement of $n-\lfloor np\rfloor$ counts, that is, at least $1-p$ in
frequency \emph{uniformly in $n$}; symmetrically
$\varepsilon^{+}(\hat c)\ge1-(1-p)^{n}$ gives $\kappa^{+}(\hat c)=0$ and an
upper enlargement of $\lceil np\rceil$ counts. The exact conditions are
\begin{align*}
  \kappa^{-}(\hat c)-k^{-}(\hat c)=0
    &\iff \varepsilon^{-}(\hat c)
      <\Pr\big[\Ncount{\hat c}\le\lfloor np\rfloor\big],
    \tag{\getrefnumber{C-eq:guard-lo0}}\digesteqlabel{eq:guard-lo0}\\
  \kappa^{-}(\hat c)-k^{-}(\hat c)\le1
    &\iff \varepsilon^{-}(\hat c)
      <\Pr\big[\Ncount{\hat c}\le\lfloor np\rfloor+1\big],
    \tag{\getrefnumber{C-eq:guard-lo1}}\digesteqlabel{eq:guard-lo1}\\
  k^{+}(\hat c)-\kappa^{+}(\hat c)=0
    &\iff \varepsilon^{+}(\hat c)
      <\Pr\big[\Ncount{\hat c}\ge\lceil np\rceil\big],
    \tag{\getrefnumber{C-eq:guard-up0}}\digesteqlabel{eq:guard-up0}\\
  k^{+}(\hat c)-\kappa^{+}(\hat c)\le1
    &\iff \varepsilon^{+}(\hat c)
      <\Pr\big[\Ncount{\hat c}\ge\lceil np\rceil-1\big].
    \tag{\getrefnumber{C-eq:guard-up1}}\digesteqlabel{eq:guard-up1}
\end{align*}
Measured instead from the honest value, the guard costs at most $1/n$
unconditionally, \eqref{eq:guard-from-p} of Lemma~\ref{lem:guard}\,(iv);
\item \emph{(evaluation against upper bounds)} let $U^{-}_k$ and $U^{+}_k$,
for $k\in\{0,\dots,n\}$, be rationals determined by $n$, $p_{\hat c}$ and
$k$ alone, with
$U^{-}_k\ge\Pr[\Ncount{\hat c}<k]$, $U^{+}_k\ge\Pr[\Ncount{\hat c}>k]$,
$U^{-}_0=0$ and $U^{+}_n=0$, and with every $U^{\mp}_k$ equal to its tail
when $p_{\hat c}\in\{0,1\}$. Put
\begin{equation*}
  \tilde\kappa^{-}(\hat c)\coloneqq\max\big\{k\in\{0,\dots,n\}:\
    U^{-}_{j}\le\varepsilon^{-}(\hat c)\text{ for }1\le j\le k\big\},
  \qquad
  \tilde\kappa^{+}(\hat c)\coloneqq\min\big\{k\in\{0,\dots,n\}:\
    U^{+}_{k}\le\varepsilon^{+}(\hat c)\big\}.
\end{equation*}
Then $\tilde R^{\mp}=\tilde\kappa^{\mp}(\hat c)$ is an admissible endpoint
rule (Definition~\ref{def:endpoint-rule}\,(a)--(c)), monotone in the
allocation, with $\tilde\kappa^{-}(\hat c)\le\kappa^{-}(\hat c)$ and
$\tilde\kappa^{+}(\hat c)\ge\kappa^{+}(\hat c)$, and it is the exact rule
when every $U^{\mp}_k$ equals its tail. Consequently Lemma~\ref{lem:guard}
applies to it.
\end{enumerate}
\end{digestclaim}

\section*{9\quad The construction map}\pdfbookmark[1]{9 The construction map}{digest-group-statistics/binomial-box/def:construction}

% claim statistics/binomial-box/def:construction
\digestlabel{def:construction}{9.1}\pdfbookmark[2]{Definition 9.1 The construction map}{digest-statistics/binomial-box/def:construction}
\renewcommand*{\RestatedName}{\hyperref{contract.pdf}{}{claim-statistics/binomial-box/def:construction}{Definition~9.1}}%
\begin{digestclaim}[The construction map]
Let
$(\Chatset,{\honest}^{-},{\honest}^{+},n,\ecomp,\Chatcon,\varepsilon^{-},\varepsilon^{+})$
be the data of an admissible input (Definition~\ref{def:endpoint-rule}),
where $({\honest}^{-},{\honest}^{+})$ is its honest envelope
(Definition~\ref{def:envelope}) --- an input $\honest$ in the
point form being read as the degenerate envelope
$(\honest,\honest)$ --- and let $R$ be an admissible endpoint rule
(Definition~\ref{def:endpoint-rule}). With the guarded endpoints
$k^{\pm}$ of $R$, \eqref{eq:guard}, define
\begin{equation*}\tag{\getrefnumber{C-eq:boxbounds}}\digesteqlabel{eq:boxbounds}
  q^{-}(\hat c)\coloneqq
  \begin{cases}k^{-}(\hat c)/n,&\hat c\in\Chatcon,\\[2pt] 0,&\hat c\notin\Chatcon,\end{cases}
  \qquad
  q^{+}(\hat c)\coloneqq
  \begin{cases}k^{+}(\hat c)/n,&\hat c\in\Chatcon,\\[2pt] 1,&\hat c\notin\Chatcon.\end{cases}
\end{equation*}
The \emph{construction map} $\mathsf{Acc}_R$ --- written $\mathsf{Acc}$
where the rule is clear or is the default --- returns the
$\mathrm{polytope}$ geometry in canonical box form with these
$(\qlo,\qhi)$, i.e.\ the data \eqref{eq:boxrows}, and
\begin{equation*}\tag{\getrefnumber{C-eq:accset}}\digesteqlabel{eq:accset}
  \acSet\ \coloneqq\ \{\mathbf{q}\in\simplex{\Chatset}:\
  q^{-}(\hat c)\le q(\hat c)\le q^{+}(\hat c)\ \ \forall\hat c\in\Chatset\}.
\end{equation*}
\end{digestclaim}

% claim statistics/binomial-box/thm:box-construction
\digestlabel{thm:box-construction}{9.2}\pdfbookmark[2]{Theorem 9.2 Properties of the constructed box}{digest-statistics/binomial-box/thm:box-construction}
\begin{mainclaim}
\renewcommand*{\RestatedName}{\hyperref{contract.pdf}{}{claim-statistics/binomial-box/thm:box-construction}{Theorem~9.2}}%
\begin{digestclaim}[Properties of the constructed box]
On every admissible input (Definition~\ref{def:endpoint-rule}) and for every
admissible endpoint rule $R$, under a conservative evaluation
(Definition~\ref{def:conservative-evaluation}):
\begin{enumerate}[label=(\roman*),leftmargin=2.2em]
\item $\mathsf{Acc}_R$ is well defined, and all of its output data are exact
rationals: each $q^{\pm}(\hat c)$ is either an integer multiple of $1/n$ or one
of $0,1$, and the entries of $B$ lie in $\{0,\pm1\}$;
\item $0\le\qlo\le\qhi\le\mathbf{1}$ componentwise, so the pair
$(\qlo,\qhi)$ meets the conditions of the canonical box form of
Definition~\ref{def:taxonomy};
\item $\mathsf{H}\subseteq\acSet$, where $\mathsf{H}$ is the
envelope set \eqref{eq:envelope-set} of the input envelope: every
honest law the envelope admits is accepted, and for the degenerate envelope
this reads $\honest\in\acSet$. The envelope set is non-empty; consequently
$\acSet\neq\varnothing$ and, by Theorem~\ref{thm:geometry-validation}\,(iii),
$\sum_{\hat c}q^{-}(\hat c)\le1\le\sum_{\hat c}q^{+}(\hat c)$;
\item $\acSet$ is a non-empty, closed, convex, compact subset of
$\simplex{\Chatset}$; \xref{meat-core/def:admissible-acceptance-set} holds for the output of
$\mathsf{Acc}_R$ by construction, for every admissible input and every
admissible rule.
\end{enumerate}
\end{digestclaim}
\end{mainclaim}

\section*{10\quad Completeness}\pdfbookmark[1]{10 Completeness}{digest-group-statistics/binomial-box/lem:envelope-dominance}

% claim statistics/binomial-box/lem:envelope-dominance
\digestlabel{lem:envelope-dominance}{10.1}\pdfbookmark[2]{Lemma 10.1 Binomial dominance under bounded conditional success probabilities}{digest-statistics/binomial-box/lem:envelope-dominance}
\renewcommand*{\RestatedName}{\hyperref{contract.pdf}{}{claim-statistics/binomial-box/lem:envelope-dominance}{Lemma~10.1}}%
\begin{digestclaim}[Binomial dominance under bounded conditional success probabilities]
Let $\mathcal{F}_0\subseteq\mathcal{F}_1\subseteq\dots\subseteq\mathcal{F}_n$
be a filtration on a probability space, $n\ge1$, and let
$\chi_1,\dots,\chi_n$ be $\{0,1\}$-valued random variables with $\chi_j$
$\mathcal{F}_j$-measurable. Let $0\le a\le b\le1$ and suppose that, for every
$j\in\{1,\dots,n\}$, almost surely
\begin{equation*}\tag{\getrefnumber{C-eq:dominance-hyp}}\digesteqlabel{eq:dominance-hyp}
  a\ \le\ P_j\ \coloneqq\ \Pr\big[\chi_j=1\ \big|\ \mathcal{F}_{j-1}\big]\ \le\ b .
\end{equation*}
Put $S\coloneqq\sum_{j=1}^{n}\chi_j$, and for $m\in\{0,\dots,n\}$,
$r\in[0,1]$ and an integer $k$ write
$G^{r}_{m}(k)\coloneqq\Pr[\mathrm{B}^{r}_{m}\ge k]$, where
$\mathrm{B}^{r}_{m}$ is binomial with parameters $(m,r)$ ($\mathrm{B}^{r}_{0}=0$).
Then for every integer $k$
\begin{equation*}\tag{\getrefnumber{C-eq:dominance}}\digesteqlabel{eq:dominance}
  \Pr[S\ge k]\ \le\ \Pr\big[\mathrm{B}^{b}_{n}\ge k\big],
  \qquad
  \Pr[S\le k]\ \le\ \Pr\big[\mathrm{B}^{a}_{n}\le k\big],
\end{equation*}
and equivalently
\begin{equation*}\tag{\getrefnumber{C-eq:dominance-strict}}\digesteqlabel{eq:dominance-strict}
  \Pr[S>k]\ \le\ \Pr\big[\mathrm{B}^{b}_{n}>k\big],
  \qquad
  \Pr[S<k]\ \le\ \Pr\big[\mathrm{B}^{a}_{n}<k\big].
\end{equation*}
\end{digestclaim}

% claim statistics/binomial-box/thm:completeness-bound
\digestlabel{thm:completeness-bound}{10.2}\pdfbookmark[2]{Theorem 10.2 The completeness bound}{digest-statistics/binomial-box/thm:completeness-bound}
\begin{mainclaim}
\renewcommand*{\RestatedName}{\hyperref{contract.pdf}{}{claim-statistics/binomial-box/thm:completeness-bound}{Theorem~10.2}}%
\begin{digestclaim}[The completeness bound]
Assume a conservative evaluation
(Definition~\ref{def:conservative-evaluation}), and assume A-HON-ENV
(Assumption~\ref{ass:hon-env}) for the
honest envelope $({\honest}^{-},{\honest}^{+})$ of the input: on an
honest run there is a filtration
$\mathcal{F}_0\subseteq\dots\subseteq\mathcal{F}_n$ with $\Chat_j$
$\mathcal{F}_j$-measurable and
${\honest}^{-}(\hat c)\le\Pr[\Chat_j=\hat c\,|\,\mathcal{F}_{j-1}]\le
{\honest}^{+}(\hat c)$ almost surely, for every $j$ and $\hat c$. Let
$\acSet$ be the output of $\mathsf{Acc}_R$
(Definition~\ref{def:construction}) for an admissible endpoint rule
$R$ (Definition~\ref{def:endpoint-rule}). Then
\begin{equation*}\tag{\getrefnumber{C-eq:completeness}}\digesteqlabel{eq:completeness}
  \Pr\big[\operatorname{freq}_{\Chat_1^n}\notin\acSet\big]
  \ \le\ \varepsilon^{\star}
  \ \coloneqq\ \sum_{\hat c\in\Chatcon}
  \Big(\Pr\big[\Ncount{\hat c}^{-}<k^{-}(\hat c)\big]
      +\Pr\big[\Ncount{\hat c}^{+}>k^{+}(\hat c)\big]\Big)
  \ \le\ \ecomp ,
\end{equation*}
where $\Ncount{\hat c}^{-}$ and $\Ncount{\hat c}^{+}$ are binomial with
parameters $(n,{\honest}^{-}(\hat c))$ and $(n,{\honest}^{+}(\hat c))$
respectively (Definition~\ref{def:endpoint-rule}). Equivalently, the honest
run passes the acceptance test, the event $\OmAT$ of
\xref{meat-core/def:postprocessing-model}, with probability
$\Pr[\OmAT]=\Pr[\operatorname{freq}_{\Chat_1^n}\in\acSet]
\ge1-\varepsilon^{\star}\ge1-\ecomp$. For the degenerate envelope both
$\Ncount{\hat c}^{\mp}$ are binomial $(n,\honest(\hat c))$.
\end{digestclaim}
\end{mainclaim}

\section*{11\quad Honest containment and infeasibility}\pdfbookmark[1]{11 Honest containment and infeasibility}{digest-group-statistics/binomial-box/thm:honest-containment}

% claim statistics/binomial-box/thm:honest-containment
\digestlabel{thm:honest-containment}{11.1}\pdfbookmark[2]{Theorem 11.1 Honest containment and infeasibility}{digest-statistics/binomial-box/thm:honest-containment}
\renewcommand*{\RestatedName}{\hyperref{contract.pdf}{}{claim-statistics/binomial-box/thm:honest-containment}{Theorem~11.1}}%
\begin{digestclaim}[Honest containment and infeasibility]
Let $\acSet$ be the output of $\mathsf{Acc}_R$ on an admissible input
(Definition~\ref{def:endpoint-rule}), for any
admissible endpoint rule $R$, under a conservative evaluation
(Definition~\ref{def:conservative-evaluation}).
\begin{enumerate}[label=(\roman*),leftmargin=2.2em]
\item $\mathsf{H}\subseteq\acSet$, where $\mathsf{H}$ is the
envelope set \eqref{eq:envelope-set} of the input envelope, for every
admissible $\ecomp\in[0,1)$, every $\Chatcon$ and every admissible
allocation; for the degenerate envelope, $\honest\in\acSet$.
\item Assume A-REACH (Assumption~\ref{ass:reach}): there exists
$\omega^{\mathrm{hon}}\in\Feas$, the feasible set of the consuming
declaration $\Decl$ (\xref{protocol-compiler/def:feasible}), with
$\mathbf{p}^{\mathrm{hon}}\coloneqq\Phimap[\omega^{\mathrm{hon}}]\in\mathsf{H}$
(for the degenerate envelope, $\Phimap[\omega^{\mathrm{hon}}]=\honest$). Then
the pair $(\mathbf{p}^{\mathrm{hon}},\omega^{\mathrm{hon}})$ satisfies
$\mathbf{p}^{\mathrm{hon}}\in\acSet$, $\omega^{\mathrm{hon}}\in\Feas$ and
$\mathbf{p}^{\mathrm{hon}}\ll\Phimap[\omega^{\mathrm{hon}}]$. Consequently the
condition
\begin{equation*}\tag{\getrefnumber{C-eq:infeasible}}\digesteqlabel{eq:infeasible}
  \text{\emph{no} }\mathbf{q}\in\acSet\text{ is absolutely continuous with
  respect to }\Phimap[\omega]\text{ for any }\omega\in\Feas
\end{equation*}
is \emph{false} for this $\acSet$.
\item \emph{(refusals come from the input)} The validation predicate
of Theorem~\ref{thm:geometry-validation}\,(iii) is decidable in closed form, and on every
admissible input it is satisfied by the output of $\mathsf{Acc}_R$
unconditionally, for every admissible rule. Hence the construction never
fails on an admissible input, and a refusal can arise only from an input that
is not admissible. On data of the form of an admissible input
(Definition~\ref{def:endpoint-rule}), the violations of its conditions that
are decidable from this module's own inputs are exactly:
\begin{enumerate}[label=(\arabic*),leftmargin=2.2em]
\item $|\Chatset|=0$;
\item $|\Chatset|$ not finite;
\item $n<1$;
\item an honest input in the point form with
$\honest\notin\simplex{\Chatset}$, including a coordinate outside
$\Chatset$;
\item an honest input in the envelope form violating \eqref{eq:envelope}:
some entry of ${\honest}^{-}$ or ${\honest}^{+}$ outside $[0,1]$, some
${\honest}^{-}(\hat c)>{\honest}^{+}(\hat c)$,
$\sum_{\hat c}{\honest}^{-}(\hat c)>1$, $\sum_{\hat c}{\honest}^{+}(\hat
c)<1$, or a coordinate outside $\Chatset$;
\item $\ecomp\notin[0,1)$;
\item some $\varepsilon^{\pm}(\hat c)<0$;
\item $\sum_{\hat c\in\Chatcon}(\varepsilon^{-}(\hat c)+\varepsilon^{+}(\hat
c))>\ecomp$;
\item $\Chatcon\not\subseteq\Chatset$;
\item some coordinate of $\honest$ or of ${\honest}^{\mp}$, or $\ecomp$, or
some $\varepsilon^{\pm}(\hat c)$ is not an exact rational, violating
Definition~\ref{def:conservative-evaluation};
\item the requested endpoint rule, or the requested evaluation, is none of
those of clause (viii) of Definition~\ref{def:endpoint-rule}.
\end{enumerate}
Exactly one clause of an admissible input is \emph{not} decidable from this
module's inputs: clause (ix) of Definition~\ref{def:endpoint-rule}, the
condition of Definition~\ref{def:matched-alphabet} that $\Chatset$, its
enumeration and $n$ are those of the same declaration $\Decl$ the
consumer solves. That clause
is a relation between artefacts, not a property of an input.
\end{enumerate}
\end{digestclaim}

\section*{12\quad Monotonicity and limiting cases}\pdfbookmark[1]{12 Monotonicity and limiting cases}{digest-group-statistics/binomial-box/def:monotone-rule}

% claim statistics/binomial-box/def:monotone-rule
\digestlabel{def:monotone-rule}{12.1}\pdfbookmark[2]{Definition 12.1 Monotone allocation rule}{digest-statistics/binomial-box/def:monotone-rule}
\renewcommand*{\RestatedName}{\hyperref{contract.pdf}{}{claim-statistics/binomial-box/def:monotone-rule}{Definition~12.1}}%
\begin{digestclaim}[Monotone allocation rule]
A \emph{monotone allocation rule}
assigns to each budget $e\in[0,1)$ an allocation
$\varepsilon^{\pm}_{(e)}$ on $\Chatcon$ such that
$\sum_{\hat c}(\varepsilon^{-}_{(e)}(\hat c)+\varepsilon^{+}_{(e)}(\hat c))\le
e$ and each $e\mapsto\varepsilon^{\pm}_{(e)}(\hat c)$ is non-decreasing. The
\emph{uniform rule}
$\varepsilon^{\pm}_{(e)}(\hat c)=e/(2|\Chatcon|)$ is monotone.
\end{digestclaim}

% claim statistics/binomial-box/lem:widening
\digestlabel{lem:widening}{12.2}\pdfbookmark[2]{Lemma 12.2 Widening the box keeps its guarantees}{digest-statistics/binomial-box/lem:widening}
\renewcommand*{\RestatedName}{\hyperref{contract.pdf}{}{claim-statistics/binomial-box/lem:widening}{Lemma~12.2}}%
\begin{digestclaim}[Widening the box keeps its guarantees]
Let $\acSet$ be the output of $\mathsf{Acc}_R$
(Definition~\ref{def:construction}) on an admissible input
(Definition~\ref{def:endpoint-rule}) for an admissible endpoint rule
$R$, with box bounds $\qlo,\qhi$ \eqref{eq:boxbounds} and
guarded endpoints $k^{\mp}$ \eqref{eq:guard}. Let
$\tilde{\mathbf{q}}^{-},\tilde{\mathbf{q}}^{+}\in\mathbb{R}^{\Chatset}$
satisfy $\tilde q^{-}(\hat c)\le q^{-}(\hat c)$ and
$\tilde q^{+}(\hat c)\ge q^{+}(\hat c)$ for every $\hat c\in\Chatset$, and put
\begin{equation*}\tag{\getrefnumber{C-eq:widened-set}}\digesteqlabel{eq:widened-set}
  \acSet'\ \coloneqq\ \{\mathbf{q}\in\simplex{\Chatset}:\
  \tilde q^{-}(\hat c)\le q(\hat c)\le\tilde q^{+}(\hat c)\ \
  \forall\hat c\in\Chatset\}.
\end{equation*}
\begin{enumerate}[label=(\roman*),leftmargin=2.2em]
\item \emph{(containment)} $\acSet\subseteq\acSet'$.
\item \emph{(failure probability)} For all $\Chatset$-valued random
variables $\Chat_1,\dots,\Chat_n$ on one probability space,
\begin{equation*}
  \Pr\big[\operatorname{freq}_{\Chat_1^n}\notin\acSet'\big]
  \ \le\ \Pr\big[\operatorname{freq}_{\Chat_1^n}\notin\acSet\big].
\end{equation*}
\item \emph{(infima)} For every function
$\Psi:\simplex{\Chatset}\to[-\infty,+\infty]$,
\begin{equation*}\tag{\getrefnumber{C-eq:widened-inf}}\digesteqlabel{eq:widened-inf}
  \inf_{\mathbf{q}\in\acSet'}\Psi(\mathbf{q})
  \ \le\ \inf_{\mathbf{q}\in\acSet}\Psi(\mathbf{q}).
\end{equation*}
\item \emph{(integer endpoints)} Let
$\tilde k^{-},\tilde k^{+}:\Chatcon\to\mathbb{Z}$ satisfy
$0\le\tilde k^{-}(\hat c)\le k^{-}(\hat c)$ and
$k^{+}(\hat c)\le\tilde k^{+}(\hat c)\le n$ for every $\hat c\in\Chatcon$,
and let $\tilde{\mathbf{q}}^{\mp}$ be given by \eqref{eq:boxbounds} with
$\tilde k^{\mp}$ in place of $k^{\mp}$. Then $\tilde{\mathbf{q}}^{\mp}$
satisfy the hypothesis above, so (i)--(iii) hold; the pair
$(\tilde{\mathbf{q}}^{-},\tilde{\mathbf{q}}^{+})$ meets the conditions of the
canonical box form of Definition~\ref{def:taxonomy}, and its
denotation \eqref{eq:boxset} is $\acSet'$; and for every $\hat c\in\Chatcon$,
with $\Ncount{\hat c}^{\mp}$ as in Definition~\ref{def:endpoint-rule},
\begin{equation*}\tag{\getrefnumber{C-eq:widened-tails}}\digesteqlabel{eq:widened-tails}
  \Pr\big[\Ncount{\hat c}^{-}<\tilde k^{-}(\hat c)\big]
  \le\Pr\big[\Ncount{\hat c}^{-}<k^{-}(\hat c)\big]
  \le\varepsilon^{-}(\hat c),
  \qquad
  \Pr\big[\Ncount{\hat c}^{+}>\tilde k^{+}(\hat c)\big]
  \le\Pr\big[\Ncount{\hat c}^{+}>k^{+}(\hat c)\big]
  \le\varepsilon^{+}(\hat c).
\end{equation*}
Summed over $\Chatcon$, the tails at $\tilde k^{\mp}$ therefore total at most
the tails at $k^{\mp}$, and at most $\ecomp$.
\item \emph{(not the reverse)} Replacing $k^{-}(\hat c)$ by a larger integer
or $k^{+}(\hat c)$ by a smaller one keeps, in general, neither the budgets of
\eqref{eq:widened-tails} nor the tail sum at $k^{\mp}$ as a bound on the
probability that the frequency vector leaves the set. Take
$\Chatset=\{\hat c_0,\hat c_1\}$, the degenerate envelope at
$\honest=(\tfrac12,\tfrac12)$ (Definition~\ref{def:envelope}), $n=3$,
$\Chatcon=\{\hat c_0\}$,
$\varepsilon^{-}(\hat c_0)=\varepsilon^{+}(\hat c_0)=\tfrac25$,
$\ecomp=\tfrac45$ and the exact binomial quantile rule of
Definition~\ref{def:interval}. Then $k^{-}(\hat c_0)=1$, $k^{+}(\hat c_0)=2$,
and the tails at these endpoints sum to $\tfrac14$. Replacing
$k^{-}(\hat c_0)$ by $2$, or $k^{+}(\hat c_0)$ by $1$, gives a tail of
$\tfrac12>\tfrac25$ on the replaced side; and for $\Chat_1,\Chat_2,\Chat_3$
independent with law $\honest$, the frequency vector lies outside the set
\eqref{eq:accset} built from the replaced endpoints with probability
$\tfrac58>\tfrac14$, while it lies outside $\acSet$ with probability
$\tfrac14$.
\end{enumerate}
\end{digestclaim}

% claim statistics/binomial-box/thm:monotonicity-limits
\digestlabel{thm:monotonicity-limits}{12.4}\pdfbookmark[2]{Theorem 12.4 Monotonicity in the budget and limiting cases}{digest-statistics/binomial-box/thm:monotonicity-limits}
\renewcommand*{\RestatedName}{\hyperref{contract.pdf}{}{claim-statistics/binomial-box/thm:monotonicity-limits}{Theorem~12.4}}%
\begin{digestclaim}[Monotonicity in the budget and limiting cases]
Fix $\Chatset$, an honest envelope $({\honest}^{-},{\honest}^{+})$
with envelope set $\mathsf{H}$ (Definition~\ref{def:envelope}; a single
law $\honest$ is the degenerate envelope, $\mathsf{H}=\{\honest\}$),
$n$, $\Chatcon$ and an admissible endpoint rule $R$
(Definition~\ref{def:endpoint-rule}) under a conservative evaluation
(Definition~\ref{def:conservative-evaluation}), write $p^{\mp}_{\hat c}$ for the
envelope's values as there, and write $\acSet(e)$ for the output of
$\mathsf{Acc}_R$ at budget $e$ under a fixed monotone allocation rule
(Definition~\ref{def:monotone-rule}).
\begin{enumerate}[label=(\roman*),leftmargin=2.2em]
\item \emph{(budget monotonicity: the set shrinks as the budget grows)} If
$R$ is monotone in the allocation and $0\le e_1\le e_2<1$, then
$\acSet(e_2)\subseteq\acSet(e_1)$.
\item \emph{(consequence for any consumer infimum)} Under the hypothesis of
(i), for every function $\Psi:\simplex{\Chatset}\to[-\infty,+\infty]$ and
$e_1\le e_2$,
\begin{equation*}\tag{\getrefnumber{C-eq:inf-mono}}\digesteqlabel{eq:inf-mono}
  \inf_{\mathbf{q}\in\acSet(e_1)}\Psi(\mathbf{q})
  \ \le\ \inf_{\mathbf{q}\in\acSet(e_2)}\Psi(\mathbf{q}) .
\end{equation*}
In particular, taking
$\Psi(\mathbf{q})=\inf_{\omega\in\Feas}O_{\alpha}(\mathbf{q},\omega)$, the
quantity $\hlow$, which \xref{meat-core/def:feasible} defines as
$\inf_{\mathbf{q}\in\acSet}\inf_{\omega\in\Feas}O_{\alpha}$, is
\emph{non-decreasing} in $\ecomp$: a weaker completeness guarantee buys a
larger or equal entropy bound.
\item \emph{(limit $\ecomp=0$, every rule)} At $e=0$ every
$\varepsilon^{\pm}(\hat c)=0$ and
\begin{equation*}
\begin{aligned}
  \acSet(0)=\big\{\mathbf{q}\in\simplex{\Chatset}:\
  &q(\hat c)=0\ \text{for }\hat c\in\Chatcon\text{ with }p^{+}_{\hat c}=0,\\
  &q(\hat c)=1\ \text{for }\hat c\in\Chatcon\text{ with }p^{-}_{\hat
  c}=1\big\},
\end{aligned}
\end{equation*}
with $\varepsilon^{\star}=0$. If $p^{+}_{\hat c}>0$ and $p^{-}_{\hat c}<1$ for
every $\hat c\in\Chatcon$ (for the degenerate envelope: if
$0<\honest(\hat c)<1$ for every $\hat c\in\Chatcon$), then
$\acSet(0)=\simplex{\Chatset}$: a zero completeness budget buys the whole
simplex and no constraint at all.
\item \emph{(limit $n\to\infty$, the exact rule)} Let $R$ be the exact rule
of Definition~\ref{def:interval}, let $\hat c\in\Chatcon$ and write
$p^{\mp}=p^{\mp}_{\hat c}$. Then
\begin{equation*}\tag{\getrefnumber{C-eq:width}}\digesteqlabel{eq:width}
  0\ \le\ q^{+}(\hat c)-p^{+}\ \le\
  \sqrt{\frac{p^{+}(1-p^{+})}{n\,\varepsilon^{+}(\hat c)}}+\frac1n,
  \qquad
  0\ \le\ p^{-}-q^{-}(\hat c)\ \le\
  \sqrt{\frac{p^{-}(1-p^{-})}{n\,\varepsilon^{-}(\hat c)}}+\frac1n ,
\end{equation*}
the first whenever $0<p^{+}<1$ and $\varepsilon^{+}(\hat c)>0$, the second
whenever $0<p^{-}<1$ and $\varepsilon^{-}(\hat c)>0$.
Consequently, if $\Chatcon=\Chatset$, $0<p^{\mp}_{\hat c}<1$ and
$\varepsilon^{\pm}(\hat c)>0$ for all $\hat c$, and the allocation is held
fixed as $n$ grows (for instance the uniform rule at a fixed $e>0$), then
every $\mathbf{q}\in\acSet$ satisfies
$p^{-}_{\hat c}-O(n^{-1/2})\le q(\hat c)\le p^{+}_{\hat c}+O(n^{-1/2})$
uniformly in $\hat c$: $\acSet$ lies coordinatewise within $O(n^{-1/2})$ of
the box $\prod_{\hat c}[p^{-}_{\hat c},p^{+}_{\hat c}]$, and it contains
$\mathsf{H}$ by Theorem~\ref{thm:box-construction}\,(iii). For the degenerate envelope,
$\sup_{\mathbf{q}\in\acSet}\|\mathbf{q}-\honest\|_{\infty}=O(n^{-1/2})$ and
$\acSet$ converges to $\{\honest\}$ in Hausdorff distance.
\item \emph{(no monotonicity in $n$, with a counterexample; the exact rule)}
For the exact rule, $\acSet$ is \emph{not} monotone in $n$, already for a
degenerate envelope. Take $\Chatset=\{\hat c_0,\hat c_1\}$, the degenerate
envelope at $\honest=(\tfrac12,\tfrac12)$, $\Chatcon=\{\hat c_0\}$ and
$\varepsilon^{-}(\hat c_0)=\varepsilon^{+}(\hat c_0)=\tfrac25$ (so
$\ecomp=\tfrac45$). Then
\begin{equation*}
\begin{aligned}
  n&=2:\quad [k^{-},k^{+}]=[1,1],\quad
   \acSet=\{(\tfrac12,\tfrac12)\};\\
  n&=3:\quad [k^{-},k^{+}]=[1,2],\quad
   \acSet=\{\mathbf{q}\in\simplex{\Chatset}:\tfrac13\le q(\hat
   c_0)\le\tfrac23\},
\end{aligned}
\end{equation*}
so $\acSet(n{=}2)\subsetneq\acSet(n{=}3)$: the set is strictly larger at the
larger round count.
\item \emph{(degenerate alphabet $|\Chatset|=1$, every rule)} If
$\Chatset=\{\hat c_0\}$ then, for every $n\ge1$, every $\ecomp\in[0,1)$,
every $\Chatcon\subseteq\Chatset$ and every admissible allocation,
\begin{equation*}
  {\honest}^{+}(\hat c_0)=1,
  \qquad
  \mathsf{H}=\acSet=\simplex{\Chatset}=\{\delta_{\hat c_0}\},
\end{equation*}
and $\varepsilon^{\star}=0$ whenever ${\honest}^{-}(\hat c_0)=1$, in
particular for the degenerate envelope, which is then
$\honest=\delta_{\hat c_0}$. The emitted set is therefore constant in all of
the construction's parameters on a one-letter alphabet, being at once the
whole simplex and a single point; for the degenerate envelope so are its box
data, and completeness is perfect. For a non-degenerate envelope
(${\honest}^{-}(\hat c_0)<1$, which \eqref{eq:envelope} permits here) the
lower box bound $q^{-}(\hat c_0)$ and the reported $\varepsilon^{\star}$ may be
positive and may vary with the parameters; the reported figure is still a
valid bound (by Theorem~\ref{thm:completeness-bound}), merely a loose one, the true honest
failure probability being $0$.
\end{enumerate}
\end{digestclaim}

\section*{13\quad Elementary enclosures}\pdfbookmark[1]{13 Elementary enclosures}{digest-group-statistics/binomial-box/lem:log-enclosure}

% claim statistics/binomial-box/lem:log-enclosure
\digestlabel{lem:log-enclosure}{13.1}\pdfbookmark[2]{Lemma 13.1 Rational bounds on the logarithm from integer square roots}{digest-statistics/binomial-box/lem:log-enclosure}
\renewcommand*{\RestatedName}{\hyperref{contract.pdf}{}{claim-statistics/binomial-box/lem:log-enclosure}{Lemma~13.1}}%
\begin{digestclaim}[Rational bounds on the logarithm from integer square roots]
Let $x>0$ be real and $m\ge1$ an integer. Then
\begin{equation*}\tag{\getrefnumber{C-eq:log-root-bounds}}\digesteqlabel{eq:log-root-bounds}
  m\big(1-x^{-1/m}\big)\ \le\ \ln x\ \le\ m\big(x^{1/m}-1\big),
\end{equation*}
with equality on either side if and only if $x=1$, and both slacks are bounded
by
\begin{equation*}\tag{\getrefnumber{C-eq:log-root-slack}}\digesteqlabel{eq:log-root-slack}
  \max\Big\{m\big(x^{1/m}-1\big)-\ln x,\ \ \ln x-m\big(1-x^{-1/m}\big)\Big\}
  \ \le\ \frac{(\ln x)^{2}}{2m}\,e^{|\ln x|/m}.
\end{equation*}
Consequently, if $l,u$ are rationals with $0<l\le x^{1/m}\le u$, then
\begin{equation*}\tag{\getrefnumber{C-eq:log-rational-bounds}}\digesteqlabel{eq:log-rational-bounds}
  m\Big(1-\frac1l\Big)\ \le\ \ln x\ \le\ m\,(u-1),
\end{equation*}
and for rational $x=a/b$ ($a,b$ positive integers) and $m=2^{j}$, such
$l,u$ of any prescribed positive width are computed by $j$ integer square
roots: with $\operatorname{isqrt}(N)\coloneqq\lfloor\sqrt N\rfloor$
for an integer $N\ge0$, one has for every rational $y=c/d>0$
\begin{equation*}\tag{\getrefnumber{C-eq:isqrt-bounds}}\digesteqlabel{eq:isqrt-bounds}
  \frac{\operatorname{isqrt}(c\,d\,4^{s})}{d\,2^{s}}
  \ \le\ \sqrt{y}\ <\
  \frac{\operatorname{isqrt}(c\,d\,4^{s})+1}{d\,2^{s}},
\end{equation*}
the two bounds differing by $1/(d\,2^{s})\le2^{-s}$, and iterating
\eqref{eq:isqrt-bounds} $j$ times, the lower bound always taken from the
previous lower bound and the upper from the previous upper, gives
$l\le x^{1/2^{j}}\le u$ whose width $u-l$ is an explicit rational the
caller reads off and drives below any prescribed bound by the precisions
$s$ of the steps.
\end{digestclaim}

\section*{14\quad The Bernstein rule}\pdfbookmark[1]{14 The Bernstein rule}{digest-group-statistics/binomial-box/lem:bernstein-tail}

% claim statistics/binomial-box/prop:rule-bernstein
\digestlabel{prop:rule-bernstein}{14.2}\pdfbookmark[2]{Proposition 14.2 The Bernstein endpoint rule}{digest-statistics/binomial-box/prop:rule-bernstein}
\renewcommand*{\RestatedName}{\hyperref{contract.pdf}{}{claim-statistics/binomial-box/prop:rule-bernstein}{Proposition~14.2}}%
\begin{digestclaim}[The Bernstein endpoint rule]
Assume a conservative evaluation
(Definition~\ref{def:conservative-evaluation}). Fix an integer $j\ge1$ and a
precision $s\ge0$ (parameters of the evaluation). The \emph{Bernstein rule}
$R_{\mathrm{B}}$ assigns to
$(n,p,\varepsilon^{-},\varepsilon^{+})$, with $v\coloneqq np(1-p)$:
\begin{itemize}[leftmargin=1.9em]
\item if $p=0$: $R^{-}_{\mathrm{B}}=R^{+}_{\mathrm{B}}=0$; if $p=1$:
$R^{-}_{\mathrm{B}}=R^{+}_{\mathrm{B}}=n$;
\item otherwise, for each sign separately: if $\varepsilon^{-}=0$ then
$R^{-}_{\mathrm{B}}=0$, and if $\varepsilon^{+}=0$ then
$R^{+}_{\mathrm{B}}=n$; if $\varepsilon^{\mp}>0$, let $L^{\mp}$ be the
rational upper bound $m(u-1)$ of \eqref{eq:log-rational-bounds} on
$\ln(1/\varepsilon^{\mp})$ with $m=2^{j}$ and $u$ the upper root enclosure
of \eqref{eq:isqrt-bounds} at precision $s$, let
\begin{equation*}\tag{\getrefnumber{C-eq:bernstein-endpoint-t}}\digesteqlabel{eq:bernstein-endpoint-t}
  t^{\mp}\ \coloneqq\ \text{a rational }\ \ge\
  \frac{L^{\mp}}{3}+\sqrt{\Big(\frac{L^{\mp}}{3}\Big)^{2}+2vL^{\mp}}
\end{equation*}
obtained by rounding the square root upward with \eqref{eq:isqrt-bounds},
and put
\begin{equation*}\tag{\getrefnumber{C-eq:bernstein-endpoints}}\digesteqlabel{eq:bernstein-endpoints}
  R^{-}_{\mathrm{B}}\ \coloneqq\ \max\big\{0,\ \lfloor np-t^{-}\rfloor\big\},
  \qquad
  R^{+}_{\mathrm{B}}\ \coloneqq\ \min\big\{n,\ \lceil np+t^{+}\rceil\big\}.
\end{equation*}
\end{itemize}
Then:
\begin{enumerate}[label=(\roman*),leftmargin=2.2em]
\item $R_{\mathrm{B}}$ is an admissible endpoint rule in the sense of
Definition~\ref{def:endpoint-rule}\,(a)--(c), computed from $n$, $p$ and the
allocation by finitely many operations on integers and rationals (integer
square roots included) and no others; Lemma~\ref{lem:guard} therefore
applies to it;
\item $R_{\mathrm{B}}$ is monotone in the allocation, provided every
rounding in $L^{\mp}$ and $t^{\mp}$ is a non-decreasing function of its
real argument. The integer-square-root rounding of \eqref{eq:isqrt-bounds}
is \emph{not} such a function of a rational argument as written (its excess over the
root depends on the denominator), but it becomes one when the argument is
first rounded up to the fixed dyadic grid $2^{-s}\mathbb{Z}$, i.e.\ replaced
by $\lceil y\,2^{s}\rceil/2^{s}$ before the root is taken; the evaluation
takes $u$ from \eqref{eq:isqrt-bounds} applied to that grid point (with
$d=2^{s}$ fixed) at every step of Lemma~\ref{lem:log-enclosure}'s
iteration, rounds the square root in \eqref{eq:bernstein-endpoint-t} the
same way, and takes $L^{\mp}$ and $t^{\mp}$ themselves as the least
multiples of $2^{-s}$ at or above the values so obtained. Every such step
is a composition of non-decreasing maps, and every rounding is upward, so
(i) is unaffected;
\item \emph{(half-width)} for $0<p<1$ and $\varepsilon^{\mp}>0$, the
guarded endpoints of $R_{\mathrm{B}}$ satisfy
\begin{equation*}\tag{\getrefnumber{C-eq:bernstein-width}}\digesteqlabel{eq:bernstein-width}
  0\ \le\ q^{+}(\hat c)-p\ \le\ \frac{t^{+}}{n}+\frac1n,
  \qquad
  0\ \le\ p-q^{-}(\hat c)\ \le\ \frac{t^{-}}{n}+\frac1n,
\end{equation*}
where, for the unrounded value of \eqref{eq:bernstein-endpoint-t},
\begin{equation*}\tag{\getrefnumber{C-eq:bernstein-halfwidth}}\digesteqlabel{eq:bernstein-halfwidth}
  \frac{t^{\mp}}{n}
  \ =\ \frac{L^{\mp}}{3n}+\sqrt{\Big(\frac{L^{\mp}}{3n}\Big)^{2}
        +\frac{2p(1-p)L^{\mp}}{n}}
\end{equation*}
is strictly decreasing in $n$ and is $O(n^{-1/2})$ at fixed $p$ and
$\varepsilon^{\mp}$. The rule never emits an interval narrower than the raw
quantile interval of the exact rule of Definition~\ref{def:interval}
(Proposition~\ref{prop:rule-exact}\,(iii)).
\end{enumerate}
\end{digestclaim}

\section*{15\quad The exact rule's local evaluation}\pdfbookmark[1]{15 The exact rule's local evaluation}{digest-group-statistics/binomial-box/lem:pmf-enclosure}

% claim statistics/binomial-box/prop:bracketed-endpoints
\digestlabel{prop:bracketed-endpoints}{15.3}\pdfbookmark[2]{Proposition 15.3 The enclosure evaluation of the exact rule}{digest-statistics/binomial-box/prop:bracketed-endpoints}
\renewcommand*{\RestatedName}{\hyperref{contract.pdf}{}{claim-statistics/binomial-box/prop:bracketed-endpoints}{Proposition~15.3}}%
\begin{digestclaim}[The enclosure evaluation of the exact rule]
Assume a conservative evaluation
(Definition~\ref{def:conservative-evaluation}). Let $n\ge1$, let
$p\in(0,1)$ and $\varepsilon^{+}\in(0,1)$ be rational, let
$\Ncount{}$ be binomial with parameters $(n,p)$, and let $\kappa^{+}$ be the
exact upper quantile of Definition~\ref{def:interval} at
$(n,p,\varepsilon^{+})$. Let
\begin{equation*}\tag{\getrefnumber{C-eq:tail-enclosure}}\digesteqlabel{eq:tail-enclosure}
  \tau^{-}(k)\ \le\ \Pr[\Ncount{}>k]\ \le\ \tau^{+}(k),\qquad k=0,\dots,n,
\end{equation*}
be an enclosure with $\tau^{-}(n)=\tau^{+}(n)=0$, each bound either a
rational or given through a rational bound on its logarithm, such as the
one Lemma~\ref{C-lem:tail-majorant}\,(ii)--(iii) yields with the mass point
enclosed by Lemma~\ref{C-lem:pmf-enclosure} and the logarithms by
Lemma~\ref{lem:log-enclosure}. Call $k$ \emph{certified admissible},
$\mathsf{C}^{+}(k)$, when the available upper bound proves
$\Pr[\Ncount{}>k]\le\varepsilon^{+}$: $\tau^{+}(k)\le\varepsilon^{+}$, or in
log space $U(k)\le\lambda^{-}$ for rationals $U(k)\ge\ln\Pr[\Ncount{}>k]$
and $\lambda^{-}\le\ln\varepsilon^{+}$; and \emph{certified inadmissible},
$\mathsf{D}^{+}(k)$, when the lower bound proves
$\Pr[\Ncount{}>k]>\varepsilon^{+}$: $\tau^{-}(k)>\varepsilon^{+}$, or in log
space $\Lambda(k)>\lambda^{+}$ for rationals
$\Lambda(k)\le\ln\Pr[\Ncount{}>k]$ and $\lambda^{+}\ge\ln\varepsilon^{+}$.
An index that is neither is \emph{undecided}. The predicate
$\mathsf{C}^{+}(n)$ always holds.

A \emph{bracketed upper endpoint} is an integer $\tilde k^{+}\in\{0,\dots,n\}$
with $\mathsf{C}^{+}(\tilde k^{+})$ and, if $\tilde k^{+}\ge1$,
$\lnot\mathsf{C}^{+}(\tilde k^{+}-1)$. The \emph{enclosure evaluation of the
exact rule} is the endpoint rule
$R_{\mathrm{encl}}$ whose $R^{+}_{\mathrm{encl}}$ is the bracketed upper
endpoint returned by a fixed deterministic search procedure that returns such
an endpoint for every input, whose $R^{-}_{\mathrm{encl}}$
is obtained in the same way for the lower tail through the mirror
$\Ncount{}'=n-\Ncount{}$ of Lemma~\ref{C-lem:tail-majorant}\,(v), and which
returns the values of Definition~\ref{def:endpoint-rule}\,(b) and (c) at a
zero budget and at $p\in\{0,1\}$. Then:
\begin{enumerate}[label=(\roman*),leftmargin=2.2em]
\item \emph{(admissible relaxation)} Every bracketed upper endpoint satisfies
$\Pr[\Ncount{}>\tilde k^{+}]\le\varepsilon^{+}$ and
$\kappa^{+}\le\tilde k^{+}$; mirror-wise every bracketed lower endpoint
$\tilde k^{-}$ satisfies $\Pr[\Ncount{}<\tilde k^{-}]\le\varepsilon^{-}$ and
$\tilde k^{-}\le\kappa^{-}$. Hence $R_{\mathrm{encl}}$ is an
admissible endpoint rule, Lemma~\ref{lem:guard} applies to it, and the
box
$\mathsf{Acc}_{R_{\mathrm{encl}}}$ emits (Definition~\ref{def:construction})
contains the box the exact rule emits: every $q^{-}$ is at most and every
$q^{+}$ at least the exact rule's.
\item \emph{(the bracket decided)} If $\tilde k^{+}\ge1$ and
$\mathsf{D}^{+}(\tilde k^{+}-1)$ holds, then $\tilde k^{+}=\kappa^{+}$; if
$\tilde k^{+}=0$ then $\tilde k^{+}=\kappa^{+}$. The mirror statement holds
for the lower side.
\item \emph{(conservatism in counts)} Suppose $\tilde k^{+}>\kappa^{+}$,
$\kappa^{+}\le n-2$ and $r_{\kappa^{+}+1}<1$, with $r_j$ the ratios
\eqref{C-eq:pmf-ratio}. Let $w\ge0$
be the slack of the failed certificate at $\tilde k^{+}-1$:
$w\coloneqq\ln\tau^{+}(\tilde k^{+}-1)-\ln\Pr[\Ncount{}>\tilde k^{+}-1]$, or
in log space
$w\coloneqq U(\tilde k^{+}-1)-\ln\Pr[\Ncount{}>\tilde k^{+}-1]
+(\ln\varepsilon^{+}-\lambda^{-})$. Then
\begin{equation*}\tag{\getrefnumber{C-eq:bracket-conservatism}}\digesteqlabel{eq:bracket-conservatism}
  \tilde k^{+}-\kappa^{+}\ <\ 1+\frac{w}{-\ln r_{\kappa^{+}+1}} .
\end{equation*}
The emitted endpoint is therefore the exact quantile up to a number of
counts bounded by the enclosure's slack divided by the per-count decay of
the tail at the quantile. The mirror statement holds for the lower side,
through $\Ncount{}'$ (Lemma~\ref{C-lem:tail-majorant}\,(v)): with
$\kappa^{-}=n-\kappa'^{+}$ and $\tilde k^{-}=n-\tilde k'^{+}$, where the
primed quantities are those of $\Ncount{}'$ at level $\varepsilon^{-}$,
$\kappa^{-}-\tilde k^{-}<1+w'/(-\ln r'_{\kappa'^{+}+1})$ under the mirrored
hypotheses. Consequently, wherever the bracket is decided the emitted
endpoint is the exact rule's, and so are its guarded endpoint and its box
bound; wherever it is undecided the guarded endpoint exceeds the exact rule's
by at most $\tilde k^{+}-\kappa^{+}$ counts and the box bound $q^{+}$ the
exact rule's by at most $(\tilde k^{+}-\kappa^{+})/n$, which
\eqref{eq:bracket-conservatism} bounds under the hypotheses above;
mirror-wise for $k^{-}$ and $q^{-}$. Where the hypothesis $r_{\kappa^{+}+1}<1$ fails (the
exact quantile at or below the mean, reachable at a near-unit tail budget,
Proposition~\ref{prop:rule-exact}\,(vi)) the endpoint is still admissible by
(i) and only this bound is unavailable; at $\kappa^{+}=n-1$ the excess
$\tilde k^{+}-\kappa^{+}$ is at most one count.
\end{enumerate}
\end{digestclaim}

\section*{16\quad Requirements: the request, its refusal and the reports}\pdfbookmark[1]{16 Requirements: the request, its refusal and the reports}{digest-group-statistics/binomial-box/req:rule-in-force}

% claim statistics/binomial-box/req:rule-in-force
\digestlabel{req:rule-in-force}{16.1}\pdfbookmark[2]{Requirement 16.1 The rule the request selects}{digest-statistics/binomial-box/req:rule-in-force}
\renewcommand*{\RestatedName}{\hyperref{contract.pdf}{}{claim-statistics/binomial-box/req:rule-in-force}{Requirement~16.1}}%
\begin{digestrequirement}[The rule the request selects]
The engine builds the box of Definition~\ref{def:construction} with the rule
that the request's \iface{endpoint_rule} and \iface{tail_evaluation} select,
and with no other: \iface{endpoint_rule: exact} with
\iface{tail_evaluation: sweep}, the exact binomial quantile rule of
Definition~\ref{def:interval} with its tail comparisons decided on the
closed forms, exactly or against upper bounds on the tails as
Proposition~\ref{prop:rule-exact}\,(vii) admits, which is the default; \iface{endpoint_rule: exact} with
\iface{tail_evaluation: enclosure}, the enclosure evaluation
$R_{\mathrm{encl}}$ of the exact rule
(Proposition~\ref{prop:bracketed-endpoints}), an admissible relaxation of it;
and \iface{endpoint_rule: bernstein}, the Bernstein rule of
Proposition~\ref{prop:rule-bernstein}, on which \iface{tail_evaluation}
selects only how the completeness figure is computed
(Requirement~\ref{req:completeness-report}).
\end{digestrequirement}

% claim statistics/binomial-box/req:request-defaults
\digestlabel{req:request-defaults}{16.2}\pdfbookmark[2]{Requirement 16.2 Defaults for the constrained outcomes and the allocation}{digest-statistics/binomial-box/req:request-defaults}
\renewcommand*{\RestatedName}{\hyperref{contract.pdf}{}{claim-statistics/binomial-box/req:request-defaults}{Requirement~16.2}}%
\begin{digestrequirement}[Defaults for the constrained outcomes and the allocation]
Where a request omits the constrained outcomes, the engine takes
$\Chatcon=\Chatset$ (\iface{constrained_outcomes}); where it omits the
allocation, the engine takes the uniform rule of
Definition~\ref{def:monotone-rule} at the request's budget,
$\varepsilon^{\pm}(\hat c)=\ecomp/(2|\Chatcon|)$ for every
$\hat c\in\Chatcon$ (\iface{budget_allocation}).
\end{digestrequirement}

% claim statistics/binomial-box/req:input-refusal
\digestlabel{req:input-refusal}{16.3}\pdfbookmark[2]{Requirement 16.3 Refusal of an input that violates a decidable clause}{digest-statistics/binomial-box/req:input-refusal}
\renewcommand*{\RestatedName}{\hyperref{contract.pdf}{}{claim-statistics/binomial-box/req:input-refusal}{Requirement~16.3}}%
\begin{digestrequirement}[Refusal of an input that violates a decidable clause]
Before it constructs anything, the engine decides exactly every condition of
an admissible input (Definition~\ref{def:endpoint-rule}) that its own inputs
determine on data of the form that definition fixes. On a violation it emits
\iface{infeasibility_report} naming the violated condition by its code, the
first violation in the order the interface fixes, and it emits no
\iface{acceptance_set}; it repairs no input and puts no default in place of a
violating value. The codes of the violations (1)--(11) of
Theorem~\ref{thm:honest-containment}\,(iii) are, in that order,
\iface{alphabet_empty}, \iface{alphabet_not_finite},
\iface{rounds_out_of_range}, \iface{honest_not_a_distribution},
\iface{honest_envelope_not_admissible}, \iface{budget_out_of_range},
\iface{allocation_negative}, \iface{allocation_exceeds_budget},
\iface{selector_not_a_subset}, \iface{inputs_not_rational} and
\iface{endpoint_rule_unknown}. Data that, with the defaults of
Requirements~\ref{req:request-defaults} and~\ref{req:rule-in-force} in place,
do not have that form are refused as well: the engine emits no
\iface{acceptance_set} for them and substitutes nothing for them.
\end{digestrequirement}

% claim statistics/binomial-box/req:completeness-report
\digestlabel{req:completeness-report}{16.4}\pdfbookmark[2]{Requirement 16.4 The reported completeness figure}{digest-statistics/binomial-box/req:completeness-report}
\renewcommand*{\RestatedName}{\hyperref{contract.pdf}{}{claim-statistics/binomial-box/req:completeness-report}{Requirement~16.4}}%
\begin{digestrequirement}[The reported completeness figure]
The engine reports under \iface{completeness_certified} the following figure.
On \iface{endpoint_rule: exact} with \iface{tail_evaluation: sweep} it
reports $\varepsilon^{\star}$ of Theorem~\ref{thm:completeness-bound} itself, and where it
cannot evaluate a tail exactly it rounds the tail upward and reports
$\min\{\ecomp,\tilde\varepsilon\}$ for the resulting
$\tilde\varepsilon\ge\varepsilon^{\star}$ (Remark~\ref{C-rem:epsstar}). On
every other route it reports the admissible figure of
Remark~\ref{C-rem:completeness-on-rules} that the route's row of
\iface{completeness_report} names: the allocation sum, or
$\min\{\ecomp,\tilde\varepsilon\}$ for an upward-rounded evaluation
$\tilde\varepsilon\ge\varepsilon^{\star}$ of the tails at the emitted
endpoints. It never reports a downward-rounded figure.
\end{digestrequirement}

% claim statistics/binomial-box/req:honest-model-record
\digestlabel{req:honest-model-record}{16.5}\pdfbookmark[2]{Requirement 16.5 The honest-model record beside every completeness figure}{digest-statistics/binomial-box/req:honest-model-record}
\renewcommand*{\RestatedName}{\hyperref{contract.pdf}{}{claim-statistics/binomial-box/req:honest-model-record}{Requirement~16.5}}%
\begin{digestrequirement}[The honest-model record beside every completeness figure]
With every completeness figure it reports, the engine records, in the record
\iface{honest_model}, the input and the evaluation under which the figure was
computed: the honest model (the law in the point form, or the
honest envelope), the round count $n$, the constrained outcomes
$\Chatcon$ and the allocation; the endpoint rule and the tail evaluation; and
the parameters in force, for the Bernstein rule its $j$ and $s$
(Proposition~\ref{prop:rule-bernstein}) and, wherever the tails are
evaluated by enclosures, the parameters of those enclosures
(Proposition~\ref{prop:bracketed-endpoints}). A completeness figure is
meaningful only relative to that record.
\end{digestrequirement}

\section*{17\quad Requirements: the enclosure evaluation}\pdfbookmark[1]{17 Requirements: the enclosure evaluation}{digest-group-statistics/binomial-box/req:search-procedure}

% claim statistics/binomial-box/req:search-procedure
\digestlabel{req:search-procedure}{17.1}\pdfbookmark[2]{Requirement 17.1 One fixed search for the enclosure evaluation}{digest-statistics/binomial-box/req:search-procedure}
\renewcommand*{\RestatedName}{\hyperref{contract.pdf}{}{claim-statistics/binomial-box/req:search-procedure}{Requirement~17.1}}%
\begin{digestrequirement}[One fixed search for the enclosure evaluation]
On the enclosure evaluation of the exact rule, the engine finds every
bracketed endpoint of Proposition~\ref{prop:bracketed-endpoints}, upper and
lower, by one fixed search procedure, total and deterministic, with fixed
parameters: the procedure \iface{search_procedure} records.
\end{digestrequirement}

% claim statistics/binomial-box/req:enclosure-substrate
\digestlabel{req:enclosure-substrate}{17.2}\pdfbookmark[2]{Requirement 17.2 Log-space arithmetic with directed rounding}{digest-statistics/binomial-box/req:enclosure-substrate}
\renewcommand*{\RestatedName}{\hyperref{contract.pdf}{}{claim-statistics/binomial-box/req:enclosure-substrate}{Requirement~17.2}}%
\begin{digestrequirement}[Log-space arithmetic with directed rounding]
On the enclosure evaluation the engine computes everything in log space with
directed rounding: the mass point through \eqref{C-eq:pmf-enclosure}, the window
through the rationals $A_T$ and $B_T$ of \eqref{C-eq:window-rationals}, whose
logarithms Lemma~\ref{lem:log-enclosure} encloses, and the budget through the
same lemma. It holds the products $\prod r_l$ inside $A_T$ and $B_T$ as
dyadic numbers with an explicit binary exponent, rounded in the direction
each bound requires (down for $A_T$ in the lower bound, up for $A_T+B_T$ in
the upper), not as exact rationals, and the Stirling remainders and the root
bounds as rationals of the size of a logarithm; the only numbers it ever
forms have the magnitude of logarithms.
\end{digestrequirement}

% claim statistics/binomial-box/req:endpoint-exact-record
\digestlabel{req:endpoint-exact-record}{17.3}\pdfbookmark[2]{Requirement 17.3 The bracket verdict per outcome and side}{digest-statistics/binomial-box/req:endpoint-exact-record}
\renewcommand*{\RestatedName}{\hyperref{contract.pdf}{}{claim-statistics/binomial-box/req:endpoint-exact-record}{Requirement~17.3}}%
\begin{digestrequirement}[The bracket verdict per outcome and side]
On the enclosure evaluation of the exact rule the engine records in the
diagnostic \iface{endpoint_exact}, per constrained outcome and per side,
whether the bracket was decided, that is whether the emitted endpoint is
proved equal to the exact quantile by
Proposition~\ref{prop:bracketed-endpoints}\,(ii). When it holds a
certified-inadmissible index $k_{\mathsf{D}}$ below the bracket with
$0<r_{k_{\mathsf{D}}+2}<1$, it may also report the bound
\eqref{eq:bracket-conservatism} in counts with $-\ln r_{k_{\mathsf{D}}+2}$ in
place of $-\ln r_{\kappa^{+}+1}$ and the computable upper bound
$U(\tilde k^{+}-1)-\Lambda(\tilde k^{+}-1)+(\lambda^{+}-\lambda^{-})$ on $w$
in place of $w$ (Remark~\ref{C-rem:local-evaluation}).
\end{digestrequirement}

\part*{Singleton acceptance set}\pdfbookmark[0]{Singleton acceptance set}{digest-part-4}

\section*{18\quad Completeness of a singleton acceptance set}\pdfbookmark[1]{18 Completeness of a singleton acceptance set}{digest-group-statistics/singleton/thm:singleton-completeness}

% claim statistics/singleton/thm:singleton-completeness
\digestlabel{thm:singleton-completeness}{18.1}\pdfbookmark[2]{Theorem 18.1 Completeness of a singleton acceptance set}{digest-statistics/singleton/thm:singleton-completeness}
\begin{mainclaim}
\renewcommand*{\RestatedName}{\hyperref{contract.pdf}{}{claim-statistics/singleton/thm:singleton-completeness}{Theorem~18.1}}%
\begin{digestclaim}[Completeness of a singleton acceptance set]
Assume A-HON, let $\Chatset$ and $n$ be a
matched alphabet and round count
(Definition~\ref{def:matched-alphabet}), and let
$\acSet=\{\mathbf{q}^{\ast}\}$ be a singleton geometry with
$\mathbf{q}^{\ast}\in\simplex{\Chatset}$. Call a completeness
budget $\ecomp$ \emph{admissible for this geometry} when $\ecomp\in[0,1)$ and
$\Pr[\operatorname{freq}_{\Chat_1^n}\notin\acSet]\le\ecomp$.
\begin{enumerate}[label=(\roman*),leftmargin=2.2em]
\item \emph{(exact accept probability)} With the convention $0^{0}=1$,
\begin{equation*}\tag{\getrefnumber{C-eq:singleton-acc}}\digesteqlabel{eq:singleton-acc}
  \Pr\big[\operatorname{freq}_{\Chat_1^n}\in\acSet\big]
  =\begin{dcases}
     \frac{n!}{\prod_{\hat c\in\Chatset}\big(n\,q^{\ast}(\hat c)\big)!}
     \prod_{\hat c\in\Chatset}\honest(\hat c)^{\,n\,q^{\ast}(\hat c)},
     &\text{if }n\,q^{\ast}(\hat c)\in\mathbb{Z}_{\ge0}\ \forall\hat c,\\
     0,&\text{otherwise;}
   \end{dcases}
\end{equation*}
\item \emph{(the least admissible budget, and when there is none)} the accept
probability \eqref{eq:singleton-acc} is $0$ if and only if either
$n\,q^{\ast}(\hat c)\notin\mathbb{Z}_{\ge0}$ for some $\hat c$, or
$\mathbf{q}^{\ast}\not\ll\honest$, that is, some $\hat c$ has
$q^{\ast}(\hat c)>0$ and $\honest(\hat c)=0$. In that case \emph{no}
$\ecomp\in[0,1)$ is admissible for this geometry, and no completeness
statement whatsoever can be made about it; a budget of $1$ lies outside
$[0,1)$. Otherwise the accept probability is strictly positive and the least
admissible completeness budget is
\begin{equation*}\tag{\getrefnumber{C-eq:singleton-budget}}\digesteqlabel{eq:singleton-budget}
  \ecomp=1-\Pr\big[\operatorname{freq}_{\Chat_1^n}\in\acSet\big]\ \in[0,1),
\end{equation*}
which does lie in $[0,1)$;
\item \emph{(when the budget is zero)} the budget \eqref{eq:singleton-budget}
equals $0$ if and only if $\honest$ is a point mass $\delta_{\hat c_0}$ for
some $\hat c_0\in\Chatset$ and $\mathbf{q}^{\ast}=\delta_{\hat c_0}$; in
particular it equals $0$ whenever $|\Chatset|=1$, and it is strictly positive
whenever $\honest$ gives positive mass to two distinct outcomes.
\end{enumerate}
\end{digestclaim}
\end{mainclaim}

\section*{19\quad Requirements}\pdfbookmark[1]{19 Requirements}{digest-group-statistics/singleton/req:singleton-budget-record}

% claim statistics/singleton/req:singleton-budget-record
\digestlabel{req:singleton-budget-record}{19.1}\pdfbookmark[2]{Requirement 19.1 The completeness budget of a declared singleton}{digest-statistics/singleton/req:singleton-budget-record}
\renewcommand*{\RestatedName}{\hyperref{contract.pdf}{}{claim-statistics/singleton/req:singleton-budget-record}{Requirement~19.1}}%
\begin{digestrequirement}[The completeness budget of a declared singleton]
An instance that declares a singleton acceptance set $\{\mathbf{q}^{\ast}\}$
records the completeness budget $\ecomp$ of \eqref{eq:singleton-budget}, or,
in the degenerate case of Theorem~\ref{thm:singleton-completeness}\,(ii) where the accept
probability is exactly $0$, records that \emph{no} admissible completeness
budget exists rather than a budget of $1$, which is the honest reading of a
geometry the honest run can never satisfy. In either case it claims no
completeness the set does not have.
\end{digestrequirement}

% claim statistics/singleton/req:singleton-point-form
\digestlabel{req:singleton-point-form}{19.2}\pdfbookmark[2]{Requirement 19.2 Point-form honest input for the singleton}{digest-statistics/singleton/req:singleton-point-form}
\renewcommand*{\RestatedName}{\hyperref{contract.pdf}{}{claim-statistics/singleton/req:singleton-point-form}{Requirement~19.2}}%
\begin{digestrequirement}[Point-form honest input for the singleton]
The singleton takes its honest input in the point form only: the
engine reads a degenerate envelope as its law, and refuses a
non-degenerate envelope with \iface{honest_envelope_not_degenerate} rather
than reducing it to one of its laws.
\end{digestrequirement}

\bibliographystyle{plain}
\bibliography{../../../refs/references}

\end{document}
